\documentclass[11pt,reqno]{amsart}

\usepackage[T1]{fontenc}
\usepackage{lmodern}
\usepackage{microtype}
\usepackage{amsmath,amssymb,mathtools,mathrsfs}
\usepackage{booktabs,longtable,array}
\usepackage{enumitem}
\usepackage{hyperref}
\usepackage{aliascnt}
\usepackage[nameinlink,noabbrev]{cleveref}
\usepackage{xcolor}

\hypersetup{
  colorlinks=true,
  linkcolor=blue!55!black,
  citecolor=blue!55!black,
  urlcolor=blue!55!black,
  pdftitle={The circular law for independent row slices down to logarithmic degree},
  pdfauthor={Griffin P. Hurt}
}

\setlist[enumerate]{leftmargin=2.2em,itemsep=2pt,topsep=4pt}
\setlist[itemize]{leftmargin=1.8em,itemsep=2pt,topsep=4pt}
\setcounter{tocdepth}{2}
\allowdisplaybreaks
\raggedbottom

\newtheorem{theorem}{Theorem}[section]
\newaliascnt{proposition}{theorem}
\newtheorem{proposition}[proposition]{Proposition}
\aliascntresetthe{proposition}
\newaliascnt{lemma}{theorem}
\newtheorem{lemma}[lemma]{Lemma}
\aliascntresetthe{lemma}
\newaliascnt{corollary}{theorem}
\newtheorem{corollary}[corollary]{Corollary}
\aliascntresetthe{corollary}
\newaliascnt{fact}{theorem}
\newtheorem{fact}[fact]{Fact}
\aliascntresetthe{fact}
\theoremstyle{definition}
\newaliascnt{definition}{theorem}
\newtheorem{definition}[definition]{Definition}
\aliascntresetthe{definition}
\newaliascnt{example}{theorem}
\newtheorem{example}[example]{Example}
\aliascntresetthe{example}
\theoremstyle{remark}
\newaliascnt{remark}{theorem}
\newtheorem{remark}[remark]{Remark}
\aliascntresetthe{remark}
\newtheorem*{conjecture}{Conjecture}
\newtheorem*{lemmafiveonea}{Lemma 5.1A}

\numberwithin{equation}{section}

\newcommand{\C}{\mathbb C}
\newcommand{\R}{\mathbb R}
\newcommand{\N}{\mathbb N}
\newcommand{\Z}{\mathbb Z}
\newcommand{\Pp}{\mathbb P}
\newcommand{\E}{\mathbb E}
\newcommand{\1}{\mathbf 1}
\newcommand{\Om}{\Omega}
\newcommand{\eps}{\varepsilon}
\newcommand{\wt}{\widetilde}
\newcommand{\wh}{\widehat}
\newcommand{\ol}{\overline}
\newcommand{\cK}{c_K}
\newcommand{\cP}{c_P}
\newcommand{\BL}{\mathrm{BL}}
\newcommand{\HS}{\mathrm{HS}}
\newcommand{\isv}{\mathrm{isv}}
\newcommand{\src}{\mathrm{src}}
\newcommand{\vol}{\operatorname{Vol}}
\newcommand{\dist}{\operatorname{dist}}
\newcommand{\diam}{\operatorname{diam}}
\newcommand{\supp}{\operatorname{supp}}
\newcommand{\spanv}{\operatorname{span}}
\newcommand{\rank}{\operatorname{rank}}
\newcommand{\tr}{\operatorname{tr}}
\newcommand{\Var}{\operatorname{Var}}
\newcommand{\Lip}{\operatorname{Lip}}
\newcommand{\Ext}{\operatorname{Ext}}
\newcommand{\Ber}{\operatorname{Bernoulli}}
\newcommand{\Mult}{\operatorname{Multinomial}}
\newcommand{\diag}{\operatorname{diag}}
\newcommand{\cone}[1]{\C^{\times}#1}
\newcommand{\logp}{\log^{+}}
\newcommand{\logm}{\log^{-}}
\newcommand{\ceil}[1]{\left\lceil #1\right\rceil}
\newcommand{\floor}[1]{\left\lfloor #1\right\rfloor}
\newcommand{\abs}[1]{\left|#1\right|}
\newcommand{\norm}[1]{\left\|#1\right\|}
\newcommand{\inner}[2]{\left\langle #1,#2\right\rangle}
\newcommand{\set}[1]{\left\{#1\right\}}
\newcommand{\given}{\,\middle|\,}
\newcommand{\dd}{\,d}
\newcommand{\eqdist}{\stackrel{d}{=}}
\newcommand{\littleoP}{o_{\Pp}(1)}
\newcommand{\bigOP}{O_{\Pp}(1)}

\DeclareMathOperator{\med}{med}
\DeclareMathOperator{\ESD}{ESD}

\crefname{theorem}{theorem}{theorems}
\crefname{proposition}{proposition}{propositions}
\crefname{lemma}{lemma}{lemmas}
\crefname{corollary}{corollary}{corollaries}
\crefname{fact}{fact}{facts}
\crefname{example}{example}{examples}
\crefname{definition}{definition}{definitions}
\crefname{remark}{remark}{remarks}
\crefname{section}{Section}{Sections}

\title[The circular law down to logarithmic degree]{The circular law for independent row slices down to logarithmic degree}
\author{Griffin P. Hurt}
\thanks{Copyright 2026 Griffin P. Hurt.
Licensed under \href{https://creativecommons.org/licenses/by/4.0/}{CC BY 4.0}.}
\email{gh7775@rit.edu}
\date{Revised October 10, 2026}
\subjclass[2020]{60B20, 15B52, 05C80, 60F05}
\keywords{circular law, sparse random matrices, row-regular matrices, intermediate singular values, random normals, small-ball probability}

\begin{document}

\begin{abstract}
Let $M_n$ be an $n\times n$ random $0$--$1$ matrix with independent rows, each uniform over vectors having exactly $d$ ones. We prove that, for an absolute constant $C_*>0$, the empirical spectral distribution of
\[
  \frac{M_n}{\sqrt{d(1-d/n)}}
\]
converges in probability to the circular law whenever $C_*\log n\le d$ and $d=o(n)$. This lowers the known $\log^{2+\eps}n$ threshold to sufficiently large logarithmic degree throughout the sparse range.

The main obstacle is the lack of column regularity, which removes the aggregate-profile identities available for random regular digraphs. In the lower polylogarithmic range, we replace those identities by an averaged multinomial row-profile estimate and combine it with a rectangular right-kernel analysis to control all but $O(n/d^2)$ of the smallest singular values. A regularized row-replacement argument then compares logarithmic determinants with those of an i.i.d. Bernoulli matrix. Above every fixed superlogarithmic cutoff, the same comparison is combined instead with a uniform polynomial smallest-singular-value bound. Girko hermitization and the sparse i.i.d. circular law complete the proof.
\end{abstract}

\maketitle
\begingroup
\emergencystretch=2em
\noindent\textbf{Revision note.}
This revision corrects the tensorization threshold in \cref{lem:tensorization},
makes the sufficiently large dimension condition explicit in
\cref{cor:uniform-polynomial}, and replaces the rounded-prefix comparison
in the proof of \cref{prop:Bernoulli-hermitization} by direct convergence
of the actual minor prefix. Accompanying Lean 4 proofs check all 29 named
results, including \cref{thm:main}, using the standard Lean/mathlib
foundations. The statement mappings, proof sources, and revision details
are available in the \href{https://github.com/gxh2932/circular-law-independent-row-slices}{repository}.
\par
\endgroup

\tableofcontents

\section{Introduction}

\subsection{The model and the main result}

For integers $1\le d\le n$, let
\[
  \Omega_{n,d}:=\set{x\in\set{0,1}^n:\ \sum_{j=1}^n x_j=d}.
\]
We consider an $n\times n$ random matrix $M_n$ whose rows are independent and each is uniformly distributed on $\Omega_{n,d}$. Thus every row has exactly $d$ nonzero entries, while the column sums fluctuate. We call this the \emph{independent row-slice model}.

For a square matrix $A$, write
\[
  \mu_A:=\frac1n\sum_{j=1}^n\delta_{\lambda_j(A)}
\]
for its empirical spectral distribution, with algebraic multiplicity. The circular law $\mu_{\mathrm{circ}}$ is the uniform probability measure on the unit disk in $\C$.

Put $p=d/n$. Each entry of $M_n$ has variance $p(1-p)$, so the natural bulk scale is
\[
  \sqrt{np(1-p)}=\sqrt{d(1-d/n)}.
\]
The exact row sums also give
\[
  M_n\1=d\1.
\]
Thus $d$ is a deterministic eigenvalue, which becomes a deterministic outlier at
\[
  \sqrt{\frac{d}{1-d/n}}
\]
under the bulk normalization. This single forced eigenvalue contributes only $1/n$ to the empirical spectral distribution, so it does not affect the limiting law.

The circular law for the row-slice model was established by Li, Litvak, and Yu~\cite{LLY26a} under the condition $d\ge \log^{2+\eps}n$, $d=o(n)$. Their smallest-singular-value theorem already reaches the logarithmic invertibility threshold $d\ge C\log n$. The gap between these two ranges is caused by the logarithmic contribution of intermediate singular values. The purpose of this paper is to close that gap throughout the sparse range above a sufficiently large logarithmic constant.

\begin{theorem}[Main theorem]\label{thm:main}
There exists an absolute constant $C_*>0$ with the following property. Let $d=d(n)$ be any integer sequence satisfying
\[
  C_*\log n\le d,
  \qquad
  d=o(n).
\]
If $M_n$ is the independent row-slice matrix described above, then
\[
  \mu_{M_n/\sqrt{d(1-d/n)}}\xrightarrow{\Pp}\mu_{\mathrm{circ}}.
\]
\end{theorem}

The constant $C_*$ is chosen large enough for the support-diameter estimate in the max-weight lemma, the imported smallest-singular-value and restricted-norm estimates, and the rectangular witness and net arguments in \cref{sec:kernels}; see \cref{rem:large-log-constant,rem:final-Cstar}. The proof does not establish the result for every positive constant multiplying $\log n$.

The proof uses two complementary regimes. For every fixed $\sigma\in(0,1)$, the new structural and intermediate-singular-value argument is uniform for
\begin{equation}\label{eq:lower-band-intro}
  C_*\log n\le d\le (\log n)^{1+\sigma}.
\end{equation}
For every fixed $\eps>0$, a separate superlogarithmic bridge treats
\begin{equation}\label{eq:superlog-intro}
  (\log n)^{1+\eps}\le d,
  \qquad d=o(n).
\end{equation}
The two regimes use the same regularized row-replacement comparison but different controls on the bottom singular values:
\begin{center}
\begin{tabular}{p{0.40\textwidth}p{0.23\textwidth}p{0.28\textwidth}}
\toprule
Degree regime & Prefix codimension & Bottom-tail input\\
\midrule
$C_*\log n\le d\le(\log n)^{1+\sigma}$ & $r=\ceil{n/\sqrt d}$ & new intermediate-singular-value bound\\
$(\log n)^{1+\eps}\le d$, $d=o(n)$ & $r=\ceil{n/d^{a_\eps}}$ & uniform polynomial smallest-singular-value bound\\
\bottomrule
\end{tabular}
\end{center}
Here
\[
  a_\eps=\frac{2+\eps}{2(1+\eps)}.
\]
In both rows, the comparison transfers the relevant prefix log volume to an i.i.d. Bernoulli matrix; interlacing and hermitization then recover the square-matrix spectral law.

For example, take $\sigma=\eps=1/2$ and split an arbitrary sequence at $(\log n)^{3/2}$. The lower-band theorem applies below the cutoff and the superlogarithmic bridge applies above it. This two-regime argument yields \cref{thm:main}. More generally, any common fixed exponent may be used. For degrees above $\log^{2+\eta}n$, with any fixed $\eta\in(0,1)$, the circular-law theorem of Li--Litvak--Yu~\cite[Theorem~1.3]{LLY26a} also applies directly, providing an independent overlap with the bridge rather than a third logically necessary piece of the proof.

\subsection{Context}

Girko introduced the hermitization and logarithmic-potential approach to the circular law in~\cite{Gir85}; modern accounts include~\cite{BC12,TV10}. For matrices with independent entries, the classical sequence of results includes Bai's theorem~\cite{Bai97} and the optimal finite-variance circular law of Tao and Vu~\cite{TV10}. In the sparse i.i.d. Bernoulli model, the optimal divergence condition $np\to\infty$ is known from~\cite{RT19,SSS25}.

Circular-law questions with fixed linear constraints have a separate history. Nguyen and Vu proved the circular law for matrices whose rows are independent random $\{-1,1\}$ vectors with a prescribed row sum~\cite{NV13}. For directed matrices with both row and column sums fixed, the sparse circular law was established in~\cite{Coo19,LLTTY21}. Related invertibility and singularity results for combinatorial and sparse matrices include~\cite{Tra20,FKS22,LLY26a}. We abbreviate Litvak--Lytova--Tikhomirov--Tomczak-Jaegermann--Youssef as LLTTY.

The present independent row-slice model has exact row sums and independent rows, but no exact column regularity. Li, Litvak, and Yu proved its circular law when $d\ge\log^{2+\eps}n$, $d=o(n)$, and obtained smallest-singular-value estimates down to $d\ge C\log n$~\cite{LLY26a}. The remaining gap is therefore not an invertibility threshold; it is the logarithmic contribution of a block of intermediate singular values. In the doubly regular model, exact column regularity supplies deterministic identities for aggregate row profiles. Those identities disappear here because the column sums fluctuate, and replacing them is the main new task.

\subsection{Proof architecture and new ingredients}

The proof separates imported mechanisms from arguments specific to the independent row-slice model:
\begin{center}
\begin{tabular}{p{0.43\textwidth}p{0.49\textwidth}}
\toprule
Imported interfaces & New arguments proved here\\
\midrule
LLTTY discretization and gradual-vector decomposition & Averaging the one-row multinomial profile before tensorization\\
LLTTY fixed-subspace Gaussian-normal estimates & Direct rectangular exclusion of the shifted steep and almost-constant classes\\
Li--Litvak--Yu smallest-singular-value and restricted-norm estimates & Exact-slice double counting and the $O(n/d^2)$ intermediate-singular-value bound\\
Sparse i.i.d. Bernoulli circular law and hermitization interface & Arbitrary-codimension regularized row replacement and the two-regime synthesis\\
\bottomrule
\end{tabular}
\end{center}

Fix $\sigma\in(0,1)$ and work first in the lower band~\eqref{eq:lower-band-intro}. Set
\[
  m=\floor{n/d^3}-1,
  \qquad
  \gamma=\frac1{288},
  \qquad
  \rho=\exp[-C(n/m)^\gamma],
  \qquad
  \delta_{\mathrm{cl}}=(m/n)^{\gamma/2}.
\]
Thus $(n/m)^\gamma=d^{1/96+o(1)}$ and $\delta_{\mathrm{cl}}=d^{-1/192+o(1)}$. The lower-band proof has six steps.

\smallskip
\noindent\textbf{1. Averaged row profiles.}
For a canonical approximation $y$, let $B$ be the dyadic order of the truncated LLTTY weight attached to a uniformly chosen coordinate. A with-replacement row has a multinomial profile $Q$, and the LLTTY one-row small-ball estimate remains valid conditional on $Q$. Averaging before tensorization gives
\[
  \E_Q\operatorname{SB}(Q)
  \le Ce^{-cD(B)}\prod_q\wt w_q^{-p_q}+e^{-cd^{2/3}},
  \qquad
  D(B)=\E\abs{B-B'},
\]
where $B'$ is an independent copy. The key input is
\[
  \E 2^{-\max_{a\le d}B_a}
  \le 2^{-\E B}\exp\bigl(-c\E\abs{B-B'}\bigr)
\]
under the support-diameter condition proved in \cref{sec:profiles}. Conditioning the sampled coordinates to be distinct costs only $e^{O(d^2)}=e^{o(n)}$.

\smallskip
\noindent\textbf{2. Structural control of rectangular kernels.}
The averaged estimate is inserted into the LLTTY decomposition
\[
  S_L=\bigcup_{u=4}^{v}K_u(L)\cup P_v(L).
\]
In the $K_u$ cases, either the truncated weights already supply a factor $d^{-cn}$ or two constant-mass weight populations are separated by $cd^\kappa$. In the clustered $P_v$ case, the separation is comparable to $v$, where
\[
  v\asymp \frac{d^{1/96}}{\log d},
  \qquad
  \delta_{\mathrm{cl}}v\gg\log d.
\]
The resulting energy factor pays for the remaining profile entropy. A separate rectangular right-kernel argument excludes the codimension-dependent steep and almost-constant classes and identifies the remaining vectors with the exact normalized gradual class. This yields the row-slice structural theorem for codimensions $m$ and $m+1$.

\smallskip
\noindent\textbf{3. Random normals and intermediate singular values.}
The Gaussian-normal argument of~\cite[Section~4]{LLTTY21} is model-independent once the exposed row span has the structural property. Outside $O(n\delta_{\mathrm{cl}})$ exceptional coordinates, a Gaussian normal has no $\rho$-disk containing more than $m$ coordinates. Exact-slice double counting then gives one-row concentration
\[
  O(m/n)+e^{-cd}=O(d^{-3}).
\]
Leave-one-row-out distances and the negative second-moment identity give the following bound.

\begin{theorem}[Intermediate singular values]\label{thm:intro-intermediate}
Fix $\sigma\in(0,1)$ and assume~\eqref{eq:lower-band-intro}. For every fixed $z\in\C$, there are constants $C_{\isv}(z),C'_{\isv}(z)>0$ and integers $q_0=q_0(n,z)$ with
\[
  q_0\le C'_{\isv}(z)\frac{n}{d^2}
\]
such that
\[
  \Pp\set{
    s_{n-q_0}(M_n-z\sqrt d\,I_n)
    \ge e^{-C_{\isv}(z)d^{1/96}}
  }=1-o(1).
\]
\end{theorem}

The quantitative feature needed later is the deletion count $q_0=O(n/d^2)$, not merely $q_0=o(n)$. At the lower-band prefix codimension $r\asymp n/\sqrt d$, only the bottom $q_0$ singular values pay the full smallest-singular-value cost, while the remaining omitted singular values pay the intermediate scale. After division by $n$, the two negative-log contributions are
\[
  O\!\left(\frac{\log n\,\log\log n}{d^2}\right)=o(1),
  \qquad
  O\!\left(d^{-1/2}d^{1/96}\right)=O(d^{-47/96})=o(1).
\]
This is the point at which the $O(n/d^2)$ bound closes the circular-law argument at logarithmic degree.

\smallskip
\noindent\textbf{4. Regularized row replacement.}
For a prefix with $k=n-r$ rows, a Schur-complement hybrid reduces each row replacement to a regularized distance associated with a positive contraction of trace at least $r+1$. Bernoulli convex concentration followed by conditioning on the exact row sum yields
\begin{align*}
  \frac1n\E\abs{\Psi_\eta(H_{\mathrm{slice}})-\Psi_\eta(H_{\mathrm{iid}})}
  &\le C\sqrt{\frac{\log(ed)}{Q}}
  +C\sqrt d\log\frac U\eta\,e^{-cQ},\\
  Q&=p(1-p)(r+1).
\end{align*}
For the lower band, $r=\ceil{n/\sqrt d}$ gives $Q\asymp\sqrt d$.

\smallskip
\noindent\textbf{5. Lower-band hermitization.}
Only $q_0=O(n/d^2)$ singular values pay the full smallest-singular-value cost. The intermediate bound controls the rest of the omitted block. Weak singular-measure convergence, full logarithmic-potential convergence, and Hilbert--Schmidt control give the corresponding Bernoulli bottom-tail estimate. Interlacing then converts the prefix comparison into a square log-determinant comparison, and the replacement principle closes the lower band.

\smallskip
\noindent\textbf{6. The superlogarithmic bridge.}
For fixed $\eps>0$, the refined Li--Litvak--Yu smallest-singular-value threshold is uniformly polynomial when $d\ge(\log n)^{1+\eps}$. With
\[
  a_\eps=\frac{2+\eps}{2(1+\eps)},
  \qquad
  r=\ceil{n/d^{a_\eps}},
\]
one has
\[
  Q\asymp d^{1-a_\eps}\to\infty,
  \qquad
  d^{-a_\eps}\log n\le(\log n)^{-\eps/2}=o(1).
\]
Thus all $r$ omitted singular values may pay the polynomial lower bound, so the same replacement theorem covers the full superlogarithmic range. Its overlap with the lower-band theorem proves \cref{thm:main}.

The logical dependency of the lower-band argument may be summarized as
\[
\begin{gathered}
\text{row-local small ball}\Longrightarrow\text{averaged profiles},\\
\text{averaged profiles}\Longrightarrow\text{rectangular structure},\\
\text{rectangular structure}\Longrightarrow\text{Gaussian normals},\\
\text{Gaussian normals}\Longrightarrow\text{intermediate singular values},\\
\text{intermediate singular values}\Longrightarrow\text{regularized replacement},\\
\text{regularized replacement}\Longrightarrow\text{determinants}
\Longrightarrow\text{circular law}.
\end{gathered}
\]
Appendices~\ref{app:LLTTY-extraction} and~\ref{app:dependencies} give the source map and probability bookkeeping for each arrow.

\subsection{Organization}

\Cref{sec:preliminaries} fixes notation, the lower-band regime, and the imported deterministic and probabilistic ingredients. \Cref{sec:profiles} proves the averaged profile estimate. \Cref{sec:kernels} proves the rectangular right-kernel compatibility theorem for the exact LLTTY classes. \Cref{sec:structural} closes the $K_u$ and $P_v$ branches and derives the row-slice structural theorem. \Cref{sec:normals} treats Gaussian normals and exact-slice anti-concentration. \Cref{sec:intermediate} proves \cref{thm:intro-intermediate}. \Cref{sec:replacement} proves the generalized regularized row-replacement estimate. \Cref{sec:hermitization} proves the lower-band circular law, establishes the superlogarithmic bridge, and assembles \cref{thm:main}. \Cref{sec:further-directions} discusses the remaining sublogarithmic obstruction. The appendices give a lemma-by-lemma extraction of the LLTTY deterministic consequences and a dependency ledger for all imported results.

\section{Preliminaries and imported ingredients}\label{sec:preliminaries}

\subsection{Notation}

All logarithms are natural unless a base is displayed. Constants denoted by $c,C,C'$ are positive and may change from line to line. A constant carrying a subscript may depend on that subscript but not on $n$.

Fix a parameter $\sigma\in(0,1)$. Throughout the structural argument in \cref{sec:profiles,sec:kernels,sec:structural,sec:normals,sec:intermediate} and the lower-band hermitization argument, the integer degree satisfies
\begin{equation}\label{eq:lower-band}
  C_*\log n\le d\le(\log n)^{1+\sigma}.
\end{equation}
All constants in those arguments are uniform in $d$ on this interval; the threshold from which the estimates hold may depend on $\sigma$. The superlogarithmic bridge in \cref{subsec:superlog-bridge} is stated separately and does not assume~\eqref{eq:lower-band}. Rounding is suppressed only inside asymptotic relations such as $m=n/d^3+O(1)$.

\medskip
\noindent\textbf{Constant convention.}
Let $a_3^{\src},\cK,\cP,c_{4.10}>0$ denote fixed admissible constants from the cited LLTTY statements. Choose once and for all
\[
  0<a_3\le a_3^{\src}
\]
sufficiently small for the entropy and row-load inequalities in \cref{subsec:shifted-T1T2,subsec:almost-constant}. This is permitted because the source construction takes $a_3$ to be a sufficiently small absolute constant; throughout the manuscript the LLTTY classes and the scale $n_3=\floor{a_3n}$ are defined using this single final choice. Next choose the absolute normal parameter $a\le c_{4.10}$, and then enlarge the absolute constant $C_*$ so that the support-diameter estimate, the imported smallest-singular-value and restricted-norm estimates, and the rectangular witness and net estimates all hold. After these absolute choices, fix $\sigma\in(0,1)$ and a shift bound $Z_0<\infty$. Only the threshold
\[
  n\ge n_0(\sigma,Z_0)
\]
may depend on these two parameters. No later argument further decreases $a_3$ or enlarges a source constant.

For a vector $x\in\C^n$, let
\[
  x_1^*\ge x_2^*\ge\cdots\ge x_n^*
\]
be the nonincreasing rearrangement of $(|x_j|)_{j\le n}$. For a matrix $H$ with at least as many columns as rows, its singular values are denoted
\[
  s_1(H)\ge\cdots\ge s_{\min\{\#\text{rows},\#\text{columns}\}}(H)\ge0.
\]
For a $k\times n$ matrix $H$, define its row volume by
\[
  \vol(H):=\sqrt{\det(HH^*)}=\prod_{j=1}^k s_j(H).
\]
We use
\[
  \logp t:=\max\{\log t,0\},
  \qquad
  \logm t:=\max\{-\log t,0\}.
\]

For probability measures $\mu,\nu$ on $\C$, write
\[
  d_{\BL}(\mu,\nu)
  :=\sup\set{
    \abs{\int f\,d\mu-\int f\,d\nu}:
    \norm{f}_\infty\le1,
    \ \Lip(f)\le1
  }.
\]
For random probability measures, $\mu_n\Rightarrow\mu$ in probability means $d_{\BL}(\mu_n,\mu)\to0$ in probability. The same convention is used for differences of two random empirical spectral measures.

Our Hermitian inner product is
\[
  \inner{u}{v}=\sum_{j=1}^n u_j\overline{v_j},
\]
linear in the first variable. Matrix rows act bilinearly:
\[
  r\cdot x=\sum_{j=1}^n r_jx_j.
\]
Thus $r\cdot x=\inner{x}{\overline r}$. This convention is used throughout the right-kernel, Gaussian-normal, and leave-one-row-out sections.

A row slice is a random vector $X$ uniformly distributed on $\Omega_{n,d}$. We write $p=d/n$. A Bernoulli comparator row is denoted by
\[
  \boldsymbol\xi=(\xi_1,\ldots,\xi_n),
  \qquad
  \xi_j\stackrel{\mathrm{iid}}\sim\Ber(p).
\]
We also write
\[
  \C^\times:=\C\setminus\{0\},
  \qquad
  \C^\times\mathcal X:=\set{\lambda x:\lambda\in\C^\times,\ x\in\mathcal X}
\]
for the nonzero scalar cone generated by a set $\mathcal X\subset\C^n$.

\subsection{\texorpdfstring{Discretization and the LLTTY $\ell$-decomposition}{Discretization and the LLTTY ell-decomposition}}

We recall only the portions of the discretization from~\cite[Sections~5--7]{LLTTY19} that enter our proof.

\begin{definition}[$k$-vector]\label{def:k-vector}
For an integer $k\ge1$, a vector $y\in\C^n$ is a $k$-vector if $ky_j\in\Z+i\Z$ for every $j$. The canonical $k$-approximation of $x\in\C^n$ is
\[
  y_j=\frac1k\bigl(\floor{k\Re x_j}+i\floor{k\Im x_j}\bigr).
\]
\end{definition}

Given a $k$-vector $y$, its coordinates are partitioned into level sets and then into regular and spread $\ell$-parts by the source construction. At order $j\ge0$, let $I_j$ be the coordinates not assigned at an earlier order and let $\Lambda_j=\{y_i:i\in I_j\}$. For $\lambda\in\Lambda_j$, the level set $L(j,\lambda)$ is the entire remaining $\lambda$-level when that level has cardinality less than $2^{j+1}$, and otherwise is its leftmost subset of cardinality $2^j$. Thus every nonempty $L(j,\lambda)$ has cardinality between $2^{j-1}$ and $2^{j+1}$.

The spread values at order $j$ are chosen as follows. If
\[
  \diam(\Lambda_j)<d/k,
\]
then the spread family is empty and every level set of order $j$ is assigned to the regular part. Otherwise the canonical greedy procedure of~\cite[Section~5.2]{LLTTY19} selects a maximal pairwise $d/k$-separated family $\Lambda_j^S\subset\Lambda_j$, starting from a separated pair. In particular, $\Lambda_j^S$ then contains at least two values. The spread and regular parts of order $j$ are
\[
  \mathcal L_j^S=\bigcup_{\lambda\in\Lambda_j^S}L(j,\lambda),
  \qquad
  \mathcal L_j^R=\bigcup_{\lambda\in\Lambda_j\setminus\Lambda_j^S}L(j,\lambda).
\]
Consequently, a singleton family of represented values is always regular, every nonempty spread part has height at least two, and the distinct values represented in a spread part are pairwise $d/k$-separated. Every value represented in a regular part is either in a set of diameter less than $d/k$ (when $\mathcal L_j^S=\varnothing$) or lies at distance less than $d/k$ from the greedy spread family.

Empty parts are omitted and the nonempty spread and regular parts are re-enumerated as
\[
  \mathcal L^{(1)},\ldots,\mathcal L^{(r)},
  \qquad r\le3\log n.
\]
The height $h_q=h(\mathcal L^{(q)})$ is the number of level sets comprising $\mathcal L^{(q)}$. If $\mathcal L^{(q)}$ has order $j$, then
\begin{equation}\label{eq:part-size-height}
  2^{j-1}h_q\le s_q:=|\mathcal L^{(q)}|\le2^{j+1}h_q.
\end{equation}
We use the source's lexicographic tie-breaking and ordering of the nonempty parts, which make the decomposition canonical. After permuting coordinates, the types, orders, sizes, and heights of the parts agree up to relabeling. Only these invariants are used below.

For $p_q=s_q/n$ and a profile count $Q_q\in\{0,\ldots,d\}$, define the actual weights
\begin{equation}\label{eq:actual-weights}
  w_q(Q)=
  \begin{cases}
    h_qQ_q/d,&\mathcal L^{(q)}\text{ regular},\\
    h_q\sqrt{Q_q},&\mathcal L^{(q)}\text{ spread},
  \end{cases}
\end{equation}
and the truncated weights
\begin{equation}\label{eq:truncated-weights}
  \wt w_q=
  \begin{cases}
    h_qp_q,&s_q\ge nd^{-1/3},\ \mathcal L^{(q)}\text{ regular},\\
    h_q/d,&s_q<nd^{-1/3},\ \mathcal L^{(q)}\text{ regular},\\
    h_q\sqrt{dp_q},&s_q\ge nd^{-1/3},\ \mathcal L^{(q)}\text{ spread},\\
    h_q,&s_q<nd^{-1/3},\ \mathcal L^{(q)}\text{ spread}.
  \end{cases}
\end{equation}
In particular,
\begin{equation}\label{eq:weight-range}
  \frac{h_q}{d}\le\wt w_q\le n\sqrt d.
\end{equation}
The dyadic $w$-set of order $b$ is the union of those $\mathcal L^{(q)}$ for which
\[
  2^b\le\wt w_q<2^{b+1}.
\]
The conditional one-row estimator is
\begin{equation}\label{eq:SB-actual}
  \operatorname{SB}(Q):=\min\left(1,\min_{q:Q_q>0}w_q(Q)^{-1}\right),
\end{equation}
and its truncated analogue is
\begin{equation}\label{eq:SB-truncated}
  \operatorname{SB}(\wt Q):=\min\left(1,\min_{q:Q_q>0}\wt w_q^{-1}\right).
\end{equation}

\subsection{Author-derived deterministic consequences of the LLTTY decomposition}

The following result is extracted from~\cite[Lemmas~5.1, 5.4, 5.5, 7.2, 7.7--7.9 and Theorem~5.6]{LLTTY19}. Its $K_u$ and clustered-$P_v$ conclusions include intermediate estimates from the proofs of the cited source lemmas, not merely verbatim source statements. Appendix~\ref{app:LLTTY-extraction} gives a lemma-by-lemma derivation, with exact source-version, page, lemma, and equation locators. The source works in the degree range $d\le\exp(c\sqrt{\log n})$ for an absolute $c>0$. Under~\eqref{eq:lower-band}, every fixed polylogarithmic upper bound is eventually smaller than this quantity, so the global source restriction is automatic.

In the lower-band argument, set
\[
  m:=\floor{n/d^3}-1
\]
and call a codimension admissible if
\[
  L\in\{m,m+1\}
  =\set{\floor{n/d^3}-1,\floor{n/d^3}}.
\]
For all sufficiently large $n$, one has $m\ge1$, and every admissible codimension satisfies the source restriction
\[
  1\le L\le\frac n{d^3}.
\]

For each admissible codimension $L$, let $S_L$ denote the exact normalized gradual class of~\cite[Section~4]{LLTTY19}, formed with the codimension-dependent steep classes and the final constant $a_3$ fixed above. This is the source class used in \cref{sec:structural,sec:normals,sec:intermediate}; \cref{subsec:classes} reproduces its defining inequalities in full. Thus no statement below depends on a definition deferred to a later section.

For $x\in\C^n$, a scale $u>0$, and $\rho>0$, a set $J\subset[n]$ is a $\rho$-cluster of $x$ at scale $u$ if there is $\lambda\in\C$ such that
\[
  |x_j-\lambda|\le\rho u\qquad(j\in J).
\]
When $u=x^*_{n_3}$, we simply say that $x$ has a $\rho$-cluster of cardinality $|J|$.

For an integer $k\ge1$, let
\[
  \mathcal A_k(L):=\set{y:y\text{ is the canonical $k$-approximation of some }x\in S_L}.
\]
Fix the lexicographic order on $\C$, let $y^\sharp$ denote the coordinates of $y$ rearranged in nonincreasing lexicographic order, and, for $y,\wt y\in\mathcal A_k(L)$, set
\[
  y\sim_k\wt y\quad\Longleftrightarrow\quad y^\sharp=\wt y^\sharp.
\]
Thus an equivalence class consists precisely of normalized canonical approximants having the same coordinate multiset; this relation is not asserted on the unrestricted, infinite set of all $k$-vectors. For $\mathcal X\subset S_L$, write
\[
  \mathcal A_k(\mathcal X)
  :=\left\{y:\begin{array}{l}
    y\text{ is the canonical $k$-approximation}\\[-2pt]
    \text{of some }x\in\mathcal X
  \end{array}\right\}
  \subset\mathcal A_k(L).
\]
and use the restriction of $\sim_k$ to this set.

\begin{proposition}[Author-derived deterministic consequences of the LLTTY decomposition]\label{prop:LLTTY-package}
There are absolute constants $a_3,\cK,\cP,c_0,c_1>0$ such that, for every admissible codimension $L$ and the exact normalized gradual class $S_L$, there are families $K_u(L),P_u(L)\subset S_L$ with the following properties.
\begin{enumerate}[label=\textup{(\roman*)}]
\item For every integer $v\ge5$,
\[
  S_L=\bigcup_{u=4}^{v}K_u(L)\cup P_v(L).
\]

\item Suppose additionally that $d\le n^{1/3}$ and
\[
  1\le k\le\frac{\sqrt n}{d^{3/2}}.
\]
Then $\mathcal A_k(L)$ has at most $e^n$ equivalence classes under $\sim_k$. The same bound holds for the nonempty classes obtained by restricting $\sim_k$ to any $\mathcal A_k(\mathcal X)\subset\mathcal A_k(L)$. If $\mathcal C$ is such a class with $\ell$-part sizes $s_q$ and heights $h_q$, then
\begin{equation}\label{eq:class-cardinality}
  |\mathcal C|\le\frac{n!}{\prod_qs_q!}\prod_q h_q^{s_q}.
\end{equation}

\item If $I(b)=\{q:2^b\le\wt w_q<2^{b+1}\}$ and $t_b=\sum_{q\in I(b)}s_q$, then
\begin{equation}\label{eq:factorial-compression}
  \prod_{q\in I(b)}s_q!\ge t_b!e^{-Ct_b}.
\end{equation}
Moreover, for $J(b)\subset I(b)$ consisting of parts of size $<d^{-\kappa}n$, where
\[
  \kappa=\frac{\cK a_3}{16}<\frac13,
\]
one has
\begin{equation}\label{eq:small-bin-mass}
  \left|\bigcup_{q\in J(b)}\mathcal L^{(q)}\right|\le Cd^{-\kappa}n.
\end{equation}

\item If $y$ is a $d^u$-approximation of a vector in $K_u(L)$, then the total mass of spread $\ell$-parts is at least $\cK n_3$, where $n_3=\floor{a_3n}$. With
\[
  I=\{q:s_q\ge d^{-\kappa}n\},
\]
one of the two alternatives holds:
\begin{align}
  \sum_{q\in I}s_q&\ge n-\frac{\cK n_3}{4},\tag{K1}\label{eq:K1}\\
  \sum_{q\in I^c}s_q&\ge\frac{\cK n_3}{4}.\tag{K2}\label{eq:K2}
\end{align}
In alternative~\eqref{eq:K1},
\begin{equation}\label{eq:K1-product}
  \prod_qh_q^{s_q}\prod_q\wt w_q^{-s_q}\le C^n d^{-c_1n}.
\end{equation}
In alternative~\eqref{eq:K2}, there are integers $b_1\le b_2$ such that the coordinate distribution of the dyadic weight order has mass at least $c_0$ below $b_1$ and at least $c_0$ above $b_2$, and
\begin{equation}\label{eq:K2-separation}
  b_2-b_1\ge c_1d^\kappa.
\end{equation}

\item Fix $v\ge5$ and $\rho\le d^{-v}$. Assume that
\begin{equation}\label{eq:Pv-delta}
  0<\delta\le36\cP a_3
\end{equation}
and
\begin{equation}\label{eq:Pv-separation-hyp}
  \cP a_3(v-4)
  \ge2\log_2d+2\log_2\!\left(\frac{72\sqrt d}{\delta}\right)
     +2-\log_2(\cP a_3).
\end{equation}
If $x\in P_v(L)$ has a $\rho$-cluster of at least $\delta n$ coordinates and $y$ is its $d^v$-approximation, then the dyadic weight-order distribution has mass at least $\delta/36$ below some $b_1$, mass at least $c_0$ above some $b_2$, and
\begin{equation}\label{eq:Pv-separation}
  b_2-b_1\ge\frac{\cP a_3}{2}(v-4).
\end{equation}
\end{enumerate}
\end{proposition}

The source status of the five parts is summarized below. Page numbers refer to arXiv:1801.05575v4.
\begin{center}
\small
\begin{tabular}{p{0.07\textwidth}p{0.36\textwidth}p{0.47\textwidth}}
\toprule
Part & Status & Exact locator\\
\midrule
(i) & verbatim decomposition after notation change & p.~38, Theorem~5.6\\
(ii) & direct restatement and reindexing & pp.~34--35, Lemma~5.1 and (14)--(16); p.~37, Lemmas~5.4--5.5\\
(iii) & stated factorial bound plus proof-level cutoff consequence & p.~55, Lemma~7.2; p.~61, Case~2 in the proof of Lemma~7.7\\
(iv) & author-derived packaging of the two proof branches & pp.~60--61, proof of Lemma~7.7, equations (32)--(33), (38), Cases~1--2\\
(v) & author-derived combination of two named lemmas & p.~62, Lemma~7.8; pp.~62--63, Lemma~7.9 and its proof\\
\bottomrule
\end{tabular}
\end{center}

\begin{remark}\label{rem:source-map-package}
Appendix~\ref{app:LLTTY-extraction} shows in particular that \eqref{eq:K1-product} is the three-family calculation in the first case of~\cite[Lemma~7.7]{LLTTY19}, while the separated weight populations in~\eqref{eq:K2} and part~(v) are intermediate conclusions used in the proofs of~\cite[Lemmas~7.7 and 7.9]{LLTTY19}. The factorial estimate~\eqref{eq:factorial-compression} and the small-bin estimate~\eqref{eq:small-bin-mass} use exactly the source's truncated weights and $w$-sets; no change of binning is made.
\end{remark}

\subsection{Row-slice inputs from Li--Litvak--Yu}

We use the following consequences of~\cite{LLY26a}.

\begin{proposition}[Smallest singular values]\label{prop:LLY-smallest}
There are absolute constants $C,C_1,C_2>0$ such that, whenever $C\log n\le d\le n/2$ and $|\zeta|\le\sqrt d\log\log d$, both the independent row-slice matrix $M_n$ and the i.i.d. Bernoulli$(d/n)$ matrix $B_n$ satisfy
\[
  \Pp\set{s_n(M_n-\zeta I_n)\le n^{-\beta_n}}
  +\Pp\set{s_n(B_n-\zeta I_n)\le n^{-\wt\beta_n}}
  \le\frac{C_1}{\sqrt d},
\]
for exponents
\[
  2<\beta_n,\wt\beta_n\le C_2\log\log n.
\]
\end{proposition}

This combines~\cite[Theorem~1.7, Remark~1.9, Theorem~3.1 and Remark~3.2]{LLY26a}.

\begin{corollary}[Uniform polynomial scale above a superlogarithmic cutoff]\label{cor:uniform-polynomial}
Fix $\eps>0$. There exists a constant $B_\eps<\infty$ such that, whenever
\[
  (\log n)^{1+\eps}\le d\le n/2,
  \qquad
  |\zeta|\le\sqrt d\log\log d,
\]
one has
\begin{equation}\label{eq:uniform-polynomial}
  \Pp\set{s_n(M_n-\zeta I_n)\le n^{-B_\eps}}
  +\Pp\set{s_n(B_n-\zeta I_n)\le n^{-B_\eps}}
  \le\frac C{\sqrt d}.
\end{equation}
\end{corollary}

\begin{proof}
We make the extraction from the source explicit. Write
\[
  \Gamma(n,d):=\frac{2\log(6d)}{\log(d/(C_2\log n))}.
\]
In the sparse branch of~\cite[Remark~1.9]{LLY26a}, the admissible threshold, denoted there by $\alpha(n,d)$, satisfies
\[
  \alpha(n,d)\ge n^{-2\Gamma(n,d)}.
\]
If $d\ge(\log n)^{1+\eps}$, then for all sufficiently large $n$,
\[
  \log\frac d{C_2\log n}
  =\log d-\log\log n-\log C_2
  \ge\frac{\eps}{2(1+\eps)}\log d.
\]
Also $\log(6d)\le2\log d$ for large $d$. Consequently,
\[
  \Gamma(n,d)\le\frac{8(1+\eps)}{\eps}
\]
uniformly throughout the sparse branch. Hence the source threshold there is at least $n^{-B_\eps}$ for any fixed
\[
  B_\eps\ge\frac{16(1+\eps)}{\eps}.
\]
In the complementary dense branch, the same remark gives
\[
  \alpha(n,d)\ge\frac1{2nd^{3/2}}\ge n^{-3}
\]
for all sufficiently large $n$, since $d\le n/2$. Increasing $B_\eps$ if necessary therefore gives a single polynomial threshold throughout $d\le n/2$. Remark~3.2 of the same paper permits exactly the same threshold for the Bernoulli comparator. The probability estimate is the one in \cref{prop:LLY-smallest}.
\end{proof}

\begin{proposition}[Restricted operator norm in the range used below]\label{prop:restricted-norm}
Let $H$ be an $m\times n$ matrix with independent uniform $d$-subset rows. Assume
\[
  \frac n2\le m\le n,
  \qquad
  d\le\frac n2.
\]
For every fixed $A>0$ there is $C_A<\infty$ such that, whenever $d\ge C_A\log n$,
\[
  \norm{H-\E H}\le C_A\sqrt d
\]
with probability at least $1-n^{-A}$.
\end{proposition}

\begin{proof}
Apply~\cite[Theorem~4.1]{LLY26a} with its parameters $\gamma=A+4$ and $\beta=A+1$. The source failure probability is
\[
  nm^{-\gamma}+2n^{-\beta}.
\]
Since $m\ge n/2$, this is at most $n^{-A}$ for all sufficiently large $n$. Write
\[
  N=\max\{9d,2\gamma\log m\}.
\]
Because $m\le n$ and $d\ge C_A\log n$, enlarging $C_A$ if necessary gives $N=9d$. The source estimate
\[
  \norm{H-\E H}
  \le \frac{2d}{\sqrt N}+C_1\beta\sqrt N
\]
therefore yields $\norm{H-\E H}\le C_A\sqrt d$. This is the only range of the rectangular theorem used in the manuscript.
\end{proof}

The right-kernel exclusions needed later are not imported: their complete rectangular proof is given in \cref{sec:kernels}.

\subsection{Concentration and tensorization}

We use the L\'evy concentration function
\[
  \mathcal Q(\xi,t):=\sup_{\lambda\in\C}\Pp\set{|\xi-\lambda|\le t}.
\]
The following standard tensorization statement is the form used in~\cite[Lemma~2.7]{LLTTY19}.

\begin{lemma}[Tensorization]\label{lem:tensorization}
Let $\xi_1,\ldots,\xi_K$ be independent complex random variables. If
\[
  \mathcal Q(\xi_i,\tau t)\le A\tau^2
  \qquad\text{for all }\tau\ge1\text{ and all }i,
\]
then, for an absolute $C$ and every $\tau\ge1$,
\[
  \Pp\set{\sum_{i=1}^K|\xi_i|^2\le C\tau^2t^2K}
  \le(CA\tau^2)^K.
\]
The statement remains valid after deterministic shifts of the $\xi_i$.
\end{lemma}

We also use the negative second-moment identity. If $H$ is a full-row-rank $k\times n$ matrix with rows $R_i$, then
\begin{equation}\label{eq:negative-second-moment}
  \sum_{j=1}^k s_j(H)^{-2}
  =\sum_{i=1}^k\dist\bigl(R_i,\spanv(R_\ell:\ell\ne i)\bigr)^{-2}.
\end{equation}
The inequality needed later remains valid after zero singular values are removed, or after restricting to the surviving full-row-rank submatrix.

\section{An averaged row-profile small-ball estimate}\label{sec:profiles}

This section contains the principal model-specific input. The regular-digraph proof in~\cite{LLTTY19} conditions on a global profile matrix and then uses exact column regularity to cancel profile entropy. In the independent row-slice model, global profiles fluctuate. We instead average the one-row estimator over the multinomial profile before tensorizing over the independent rows.

\subsection{The with-replacement representation}

Let $y$ be a $k$-vector with nonempty $\ell$-parts $\mathcal L^{(1)},\ldots,\mathcal L^{(r)}$. Put
\[
  s_q=|\mathcal L^{(q)}|,
  \qquad
  p_q=s_q/n.
\]
Let $J_1,\ldots,J_d$ be i.i.d. uniform random variables on $[n]$ and define
\[
  \widehat X:=\sum_{a=1}^d e_{J_a}.
\]
Its profile relative to the $\ell$-decomposition is
\[
  Q_q:=\#\set{a\le d:J_a\in\mathcal L^{(q)}},
  \qquad
  Q\sim\Mult(d;p_1,\ldots,p_r).
\]
Conditional on $Q$, the samples belonging to distinct $\ell$-parts are independent, and within $\mathcal L^{(q)}$ they are i.i.d. uniform on that part. If $\psi_{q,a}$ denotes the corresponding level value of $y$, then
\begin{equation}\label{eq:conditional-block-sum}
  \widehat X\cdot y
  \eqdist
  \sum_{q=1}^r\sum_{a=1}^{Q_q}\psi_{q,a}
  \qquad\text{conditional on }Q.
\end{equation}

\begin{lemma}[Conditional row-local small ball]\label{lem:conditional-row-small-ball}
For every profile $Q$ and every $\tau\ge1$,
\[
  \mathcal Q\!\left(\widehat X\cdot y,\frac{\tau\sqrt d}{k}\,\middle|\,Q\right)
  \le C\tau^2\operatorname{SB}(Q).
\]
\end{lemma}

\begin{proof}
This is the row-local content of~\cite[Lemma~6.4]{LLTTY19}. We indicate why no global regularity enters. The blocks in~\eqref{eq:conditional-block-sum} are independent, so convolution reduces the estimate to any visited part $q$.

The level sets inside one $\ell$-part have cardinalities within a factor four, hence
\[
  \max_z\Pp\{\psi_{q,a}=z\}\le\frac4{h_q+3}.
\]
For a regular part, the distinct represented values form a $1/k$-separated subset of the complex plane. The discrete planar concentration estimate~\cite[Proposition~2.5]{LLTTY19} gives
\[
  \mathcal Q\!\left(\sum_{a=1}^{Q_q}\psi_{q,a},\frac{\tau\sqrt d}{k}\right)
  \le C\max\left\{\frac{d\tau^2}{h_qQ_q},\frac1{h_q}\right\}
  \le\frac{C\tau^2}{w_q(Q)}.
\]

For a spread part, the distinct represented values are $d/k$-separated and, by the source-compatible construction in \cref{sec:preliminaries}, one has $h_q\ge2$. In the Miroshnikov branch $h_q>32$, we may first assume
\[
  1\le\tau^2\le w_q(Q)=h_q\sqrt{Q_q},
\]
since otherwise $C\tau^2/w_q(Q)\ge1$, after taking $C\ge1$, and the desired probability bound is trivial. Under this reduction, $Q_q\le d$, the atom bound above, and the $d/k$-separation give, for all sufficiently large $d$,
\[
  \max_{a\le Q_q}\mathcal Q\!\left(\psi_{q,a},\frac{\tau\sqrt{2d}}{k}\right)
  \le\frac{8+64\tau^2/d}{h_q}\le\frac12.
\]
This verifies the atom hypothesis in the Miroshnikov estimate. Applying the two-dimensional Esseen estimate~\cite[Proposition~2.2]{LLTTY19} when $2\le h_q\le32$, and the Miroshnikov estimate~\cite[Proposition~2.3]{LLTTY19} when $h_q>32$, yields
\[
  \mathcal Q\!\left(\sum_{a=1}^{Q_q}\psi_{q,a},\frac{\tau\sqrt d}{k}\right)
  \le\frac{C\tau^2}{h_q\sqrt{Q_q}}
  =\frac{C\tau^2}{w_q(Q)}.
\]
The lower bound $h_q\ge2$ is essential here: a one-level block is deterministic conditional on $Q_q$ and therefore belongs to the regular part. For such a block $w_q(Q)=Q_q/d\le1$, so the preceding regular estimate and the truncation in~\eqref{eq:SB-actual} give the required trivial bound.

Taking the minimum over all visited parts, and truncating the bound by one, proves the lemma. The proof uses only the counts $Q_q$, the conditional independence in~\eqref{eq:conditional-block-sum}, and the deterministic geometry of the $\ell$-decomposition. The identities $\sum_iQ_{iq}=ds_q$ and exact column regularity do not occur.
\end{proof}

\subsection{Comparison with truncated weights}

Define the large-part set
\[
  \mathscr L:=\set{q:p_q\ge d^{-1/3}}.
\]
For $q\in\mathscr L$, the marginal $Q_q$ is binomial with mean $dp_q\ge d^{2/3}$. Let
\[
  \mathcal G:=\bigcap_{q\in\mathscr L}\set{Q_q\ge\tfrac12dp_q}.
\]
Since $|\mathscr L|\le d^{1/3}$, Chernoff's inequality gives
\begin{equation}\label{eq:profile-good-failure}
  \Pp(\mathcal G^c)
  \le d^{1/3}e^{-cd^{2/3}}
  \le e^{-c'd^{2/3}}.
\end{equation}
On $\mathcal G$, the definitions~\eqref{eq:actual-weights}--\eqref{eq:truncated-weights} imply
\[
  w_q(Q)\ge\frac12\wt w_q
\]
for every visited large regular part, the same inequality with $1/\sqrt2$ for every visited large spread part, and $w_q(Q)\ge\wt w_q$ for every visited small part. Hence
\begin{equation}\label{eq:actual-truncated-comparison}
  \E_Q\operatorname{SB}(Q)
  \le2\E_Q\operatorname{SB}(\wt Q)+e^{-cd^{2/3}}.
\end{equation}

\subsection{The max-weight inequality}

Let $B$ be a real random variable and let $B_1,\ldots,B_d$ be i.i.d. copies. The next lemma is the mechanism that converts occupancy of weight bins into both a geometric mean and an entropy-dispersion penalty.

\begin{lemma}[Max-weight lemma]\label{lem:max-weight}
There are absolute constants $\eps_0,c>0$ such that the following holds for every integer $d\ge2$. If
\[
  R:=\diam(\supp B)\le\eps_0d,
\]
then
\begin{equation}\label{eq:max-weight}
  \E2^{-\max_{a\le d}B_a}
  \le2^{-\E B}\exp\bigl(-c\E|B-B'|\bigr),
\end{equation}
where $B'$ is an independent copy of $B$.
\end{lemma}

\begin{proof}
Put $\mu=\E B$ and $D=\E|B-B'|$. Let $t=\floor{d/2}$ and pair the first $2t$ samples. Set
\[
  Z_j=\max(B_{2j-1},B_{2j}),
  \qquad1\le j\le t.
\]
Since $\max_{a\le d}B_a\ge t^{-1}\sum_{j=1}^tZ_j$, with $\theta=(\log2)/t$ we have
\begin{equation}\label{eq:max-weight-mgf}
  \E2^{-\max B_a}\le\bigl(\E e^{-\theta Z_1}\bigr)^t.
\end{equation}
The identity
\[
  \max(B,B')=\frac{B+B'+|B-B'|}{2}
\]
gives
\begin{equation}\label{eq:EZ1}
  \E Z_1=\mu+\frac D2.
\end{equation}
By Efron--Stein and the fact that the maximum is $1$-Lipschitz in each coordinate,
\[
  \Var(Z_1)\le\E(B-B')^2\le RD.
\]
The bounded-variable Bernstein moment estimate yields, whenever $\theta R\le1/2$,
\[
  \log\E e^{-\theta(Z_1-\E Z_1)}\le C\theta^2RD.
\]
Indeed, $|Z_1-\E Z_1|\le R$; expanding the exponential and summing the terms of order at least two gives
\[
  \E e^{-\theta(Z_1-\E Z_1)}\le1+C\theta^2\Var(Z_1),
\]
and $\Var(Z_1)\le RD$. Combining this with~\eqref{eq:max-weight-mgf} and~\eqref{eq:EZ1}, and using $t\ge d/3$, gives
\[
  \log\E2^{-\max B_a}
  \le-(\log2)\mu-\frac{\log2}{2}D+C\frac RdD.
\]
Choosing $\eps_0$ sufficiently small absorbs the last term and proves~\eqref{eq:max-weight}.
\end{proof}

\subsection{Averaging the profile}

Define $B$ by selecting an $\ell$-part according to coordinate mass:
\begin{equation}\label{eq:B-distribution}
  \Pp\set{B=\floor{\log_2\wt w_q}}=p_q,
  \qquad
  D(B):=\E|B-B'|.
\end{equation}
Also set
\begin{equation}\label{eq:Wy}
  W_y:=\prod_{q=1}^r\wt w_q^{-p_q}.
\end{equation}
If $J_a$ falls in part $q_a$, then
\[
  \operatorname{SB}(\wt Q)
  \le\frac1{\max_{a\le d}\wt w_{q_a}}
  \le2^{-\max_{a\le d}B_a}.
\]
Moreover,
\[
  2^{-\E B}\le2\prod_q\wt w_q^{-p_q}=2W_y.
\]
Thus \cref{lem:max-weight} and~\eqref{eq:actual-truncated-comparison} prove the main proposition.

\begin{proposition}[Averaged row-profile estimate]\label{prop:averaged-profile}
There are absolute constants $\eps_0,c,C>0$ such that, whenever
\[
  \diam(\supp B)\le\eps_0d,
\]
one has
\begin{equation}\label{eq:averaged-profile}
  \E_Q\operatorname{SB}(Q)
  \le Ce^{-cD(B)}W_y+e^{-cd^{2/3}}.
\end{equation}
Consequently, for every $\tau\ge1$,
\begin{equation}\label{eq:averaged-profile-small-ball}
  \mathcal Q\!\left(\widehat X\cdot y,\frac{\tau\sqrt d}{k}\right)
  \le C\tau^2\left(Ce^{-cD(B)}W_y+e^{-cd^{2/3}}\right).
\end{equation}
The same estimate holds after adding a deterministic scalar.
\end{proposition}

\begin{proof}
Equation~\eqref{eq:averaged-profile} follows from~\eqref{eq:actual-truncated-comparison} and \cref{lem:max-weight}. Averaging \cref{lem:conditional-row-small-ball} over $Q$ gives~\eqref{eq:averaged-profile-small-ball}.
\end{proof}

\subsection{Tensorization and conditioning to the slice}

Let $\widehat X_1,\ldots,\widehat X_K$ be independent with-replacement rows and let $\widehat M_K$ be the resulting $K\times n$ matrix. By \cref{lem:tensorization,prop:averaged-profile}, for deterministic $V\in\C^K$,
\begin{equation}\label{eq:profile-tensorization-general}
  \Pp\set{
    \norm{\widehat M_Ky+V}_2
    \le\frac{\tau\sqrt{dK}}{k}
  }
  \le\left[
    C\tau^2\left(Ce^{-cD(B)}W_y+e^{-cd^{2/3}}\right)
  \right]^K.
\end{equation}
Taking
\[
  \tau=C\sqrt{\frac nK}\log d
\]
yields
\begin{equation}\label{eq:profile-tensorization}
  \begin{aligned}
  \Pp\set{
    \norm{\widehat M_Ky+V}_2
    \le C\log d\frac{\sqrt{dn}}{k}
  }
  &\le\Bigg[
    C\frac nK\log^2d\\[-2pt]
  &\hspace{4em}\times
    \left(Ce^{-cD(B)}W_y+e^{-cd^{2/3}}\right)
  \Bigg]^K.
  \end{aligned}
\end{equation}

An ordered with-replacement row has distinct samples with probability $(n)_d/n^d$, where $(n)_d=n(n-1)\cdots(n-d+1)$. Conditional on distinctness, its indicator is a uniform $d$-subset row. Therefore, for every event $\mathcal E$ depending on $K$ independent rows,
\begin{equation}\label{eq:condition-to-slice}
  \Pp_{\mathrm{slice}}(\mathcal E)
  \le\left(\frac{n^d}{(n)_d}\right)^K
  \Pp_{\mathrm{wr}}(\mathcal E).
\end{equation}
Under~\eqref{eq:lower-band}, $K\le n$ and $d=o(\sqrt n)$, so
\begin{equation}\label{eq:conditioning-cost}
  K\log\frac{n^d}{(n)_d}
  =K\sum_{j=0}^{d-1}-\log\left(1-\frac jn\right)
  =O\left(\frac{Kd^2}{n}\right)
  =O(d^2).
\end{equation}
Hence the total conditioning cost is $e^{O(d^2)}=e^{o(n)}$, uniformly in the lower band.

\begin{remark}[Why the logarithmic lower constant must be large]\label{rem:large-log-constant}
By~\eqref{eq:weight-range},
\[
  \diam(\supp B)\le\log_2n+\frac32\log_2d+2.
\]
Under~\eqref{eq:lower-band},
\[
  \frac{\diam(\supp B)}d
  \le\frac1{C_*\log2}+o(1).
\]
Thus the hypothesis of \cref{lem:max-weight} holds uniformly throughout~\eqref{eq:lower-band} once $C_*$ is a sufficiently large absolute constant. This is the new profile-theoretic reason for requiring a large constant multiplying $\log n$. The final $C_*$ must also dominate the constants in the imported smallest-singular-value and restricted-norm inputs and in the rectangular right-kernel argument; see \cref{rem:final-Cstar}.
\end{remark}

\section{Rectangular right-kernel exclusions}\label{sec:kernels}

This section connects rectangular right kernels in the independent-row model to the exact gradual class of~\cite{LLTTY19}. The result is uniform for the two codimensions in the leave-one-row-out argument; no square theorem is invoked.

The proof has three stages. We first exclude the shifted very-steep classes using independent-row witness events and center-invariant equal-block anti-concentration. We then treat the almost-constant remainder on a uniform row-load event. Finally, the deterministic LLTTY alternatives identify every remaining nonzero kernel vector with either the permitted very-steep class or a scalar multiple of the normalized gradual class.

\subsection{Complex convention and row action}

We use the Hermitian inner product
\[
  \inner{u}{v}:=\sum_{j=1}^n u_j\ol{v_j},
\]
linear in the first variable. For matrix multiplication we use the bilinear row action
\[
  r\cdot x:=\sum_{j=1}^n r_jx_j.
\]
Thus $r\cdot x=\inner{x}{\ol r}$. If
\[
  r_i=X_i-\zeta e_i,
  \qquad i\in I,
\]
and
\[
  E_I:=\spanv\set{\ol{r_i}:i\in I},
\]
then
\begin{equation}\label{eq:kernel-conjugate-span}
  E_I^\perp=\ker A_{\zeta,I},
  \qquad
  (A_{\zeta,I}x)_i=r_i\cdot x
  =\sum_{j=1}^nX_{ij}x_j-\zeta x_i.
\end{equation}
No rank assumption is involved in~\eqref{eq:kernel-conjugate-span}.

\subsection{The codimension-dependent LLTTY classes}\label{subsec:classes}

Assume the lower-band condition~\eqref{eq:lower-band} and set
\begin{equation}\label{eq:kernel-codimension}
  m=\floor{\frac n{d^3}}-1,
  \qquad
  L\in\set{m,m+1},
  \qquad
  K=n-L.
\end{equation}
Let $I\subset[n]$ be deterministic with $|I|=K$, and assume throughout that $|\zeta|\le Z_0\sqrt d$. Write
\[
  P_I:\C^n\longrightarrow\C^I,
  \qquad
  P_Ix=x_I,
  \qquad
  A_{\zeta,I}:=M_I-\zeta P_I,
\]
where $M_I$ is the restriction of $M_n$ to the rows indexed by $I$. We record the classes from~\cite[Section~4]{LLTTY19}, displaying their dependence on $L$. To avoid collision with the Bernoulli density $p=d/n$, set
\begin{equation}\label{eq:steep-epsilon-p}
  \eps_{\mathrm{st}}=\sqrt{\frac{\log d}{d}},
  \qquad
  p_{\mathrm{st}}=\floor{\frac1{5\eps_{\mathrm{st}}}}
  =\floor{\frac15\sqrt{\frac d{\log d}}}.
\end{equation}
With the fixed absolute constant $a_3>0$ from the constant convention, put
\begin{equation}\label{eq:n1n2n3}
  n_1=\ceil{\frac n{d^{3/2}}},
  \qquad
  n_2=\floor{\frac n{d^{2/3}}},
  \qquad
  n_3=\floor{a_3n}.
\end{equation}
Let $r$ be determined by
\begin{equation}\label{eq:r-steep}
  p_{\mathrm{st}}^r<n_1\le p_{\mathrm{st}}^{r+1},
\end{equation}
and let $r_0=r_0(L)$ be the least nonnegative integer for which
\begin{equation}\label{eq:r0-steep}
  p_{\mathrm{st}}^{r_0}\ge\frac{20L}{d}.
\end{equation}
For all sufficiently large $n$, uniformly for $L=m,m+1$, one has $r_0<r$. Indeed,
\[
  p_{\mathrm{st}}^{r-1}\ge\frac{n_1}{p_{\mathrm{st}}^2}
  \ge\frac{25n\log d}{d^{5/2}},
  \qquad
  \frac{20L}{d}\le\frac{20n}{d^4}.
\]
Let
\[
  \mathcal K:=\set{(c,\ldots,c):c\in\C}
\]
be the one-dimensional space of constant vectors.

The permitted very-steep source class is
\begin{equation}\label{eq:T3}
  T_3(L):=
  \set{x:\exists i\le p_{\mathrm{st}}^{r_0},\ 
  x_i^*>\left(\frac ni\right)^3x_{p_{\mathrm{st}}^{r_0}}^*}.
\end{equation}
For $r_0\le i\le r-1$, recursively define
\begin{equation}\label{eq:T0i}
  T_{0,i}(L):=
  \set{x:x\notin T_3(L)\cup\bigcup_{j=r_0}^{i-1}T_{0,j}(L),\ 
  x_{p_{\mathrm{st}}^i}^*>4d\,x_{p_{\mathrm{st}}^{i+1}}^*}.
\end{equation}
At the terminal scale set
\begin{equation}\label{eq:T0r}
  T_{0,r}(L):=
  \set{x:x\notin T_3(L)\cup\bigcup_{j=r_0}^{r-1}T_{0,j}(L),\ 
  x_{\ceil{n_1/p_{\mathrm{st}}}}^*>4d\,x_{n_1}^*},
\end{equation}
and put
\begin{equation}\label{eq:T0}
  T_0(L)=\bigcup_{i=r_0}^rT_{0,i}(L).
\end{equation}
The remaining source steep classes are
\begin{equation}\label{eq:T1}
  T_1(L):=\set{x\notin T_0(L)\cup T_3(L):x_{n_1}^*>d^{3/2}x_{n_2}^*}
\end{equation}
and
\begin{equation}\label{eq:T2}
  T_2(L):=\set{x\notin T_0(L)\cup T_1(L)\cup T_3(L):x_{n_2}^*>d^{3/2}x_{n_3}^*}.
\end{equation}
We abbreviate
\begin{equation}\label{eq:T-union}
  T(L)=\bigcup_{i=0}^3T_i(L).
\end{equation}
For $0\le i\le3$, define the shifted class
\begin{equation}\label{eq:shifted-Ti}
  T_i^{\mathcal K}(L):=
  \set{v=x+y:x\in T_i(L),\ y\in\mathcal K,\ |y_1|\le x_{n_1}^*/10}.
\end{equation}
As in the source, the classes $T_0^{\mathcal K}$, $T_1^{\mathcal K}$, and $T_2^{\mathcal K}$ are excluded, whereas $T_3^{\mathcal K}$ is part of the permitted very-steep alternative.

For $\vartheta>0$, define
\begin{equation}\label{eq:B-vartheta}
  \mathcal B(\vartheta):=
  \set{x:\exists\lambda\in\C,\ 
  \#\set{j:|x_j-\lambda|\le\vartheta x_{n_3}^*}>n-n_3}.
\end{equation}
Set
\begin{equation}\label{eq:vartheta0}
  \vartheta_0=\frac{10}{d^3}.
\end{equation}
The normalized gradual class used in the LLTTY decomposition is literally the following.

\begin{definition}[LLTTY gradual class]\label{def:gradual-class}
\begin{equation}\label{eq:gradual-class}
  S_L:=\set{x\notin T(L)\cup\mathcal B(\vartheta_0):x_{n_3}^*=1}.
\end{equation}
Its dependence on $L$ is through $r_0(L)$ and the recursive steep classes.
\end{definition}

Let $c_{4.10}>0$ be the absolute constant from \cite[Lemma~4.10]{LLTTY21}. Fix a sufficiently small absolute constant $a\in(0,c_{4.10}]$ for the random-normal theorem, and put
\begin{equation}\label{eq:qL}
  q_L=\floor{aL},
\end{equation}
and define
\begin{equation}\label{eq:VL}
  V_L:=\set{0\ne x:\exists i\le q_L,\ 
  x_i^*>0.9\left(\frac ni\right)^3x_{q_L}^*}.
\end{equation}

\subsection{\texorpdfstring{The shifted $T_3$ class is permitted}{The shifted T3 class is permitted}}

The deterministic inclusion from~\cite[Section~4.1, equation~(8)]{LLTTY19} gives
\begin{equation}\label{eq:T3-shift-inclusion}
  T_3^{\mathcal K}(L)
  \subset
  \set{v:\exists i\le p_{\mathrm{st}}^{r_0},\ 
  v_i^*>0.9\left(\frac ni\right)^3v_{p_{\mathrm{st}}^{r_0}}^*}.
\end{equation}
For completeness, if $v=x+c\1$ and $x\in T_3(L)$, then
\[
  |c|\le x_{n_1}^*/10\le x_{p_{\mathrm{st}}^{r_0}}^*/10.
\]
If $i$ witnesses~\eqref{eq:T3}, order-statistic stability gives
\[
  v_i^*\ge x_i^*-|c|,
  \qquad
  v_{p_{\mathrm{st}}^{r_0}}^*\le x_{p_{\mathrm{st}}^{r_0}}^*+|c|.
\]
Since $(n/i)^3>10$ for all relevant indices and all large $n$, the constant $0.9$ in~\eqref{eq:T3-shift-inclusion} follows.

The minimality of $r_0$ and~\eqref{eq:steep-epsilon-p} yield
\begin{equation}\label{eq:pr0-upper}
  p_{\mathrm{st}}^{r_0}<\frac{20p_{\mathrm{st}}L}{d}
  \le\frac{4L}{\sqrt{d\log d}}.
\end{equation}
After increasing $C_*$, the last quantity is at most $aL/2$. Since $L\to\infty$,
\begin{equation}\label{eq:pr0-qL}
  p_{\mathrm{st}}^{r_0}\le q_L.
\end{equation}
Monotonicity of order statistics in~\eqref{eq:T3-shift-inclusion} therefore implies
\begin{equation}\label{eq:T3-in-VL}
  T_3^{\mathcal K}(L)\subset V_L.
\end{equation}
All estimates are uniform for $L=m,m+1$.

\subsection{\texorpdfstring{Unique-coordinate witnesses for $T_0^{\mathcal K}$}{Unique-coordinate witnesses for T0K}}

For $r_0\le i<r$, put
\begin{equation}\label{eq:si-ti}
  s_i=p_{\mathrm{st}}^i,
  \qquad
  t_i=p_{\mathrm{st}}^{i+1}.
\end{equation}
At the terminal scale put
\begin{equation}\label{eq:sr-tr}
  s_r=\ceil{\frac{n_1}{p_{\mathrm{st}}}},
  \qquad
  t_r=n_1.
\end{equation}
The choice $t_r=n_1$, rather than $p_{\mathrm{st}}\ceil{n_1/p_{\mathrm{st}}}$, is important: it matches the defining jump in~\eqref{eq:T0r} and gives direct control of the constant shift.

For disjoint sets $J^\ell,J^r\subset[n]$ satisfying $|J^\ell|=s_i$ and $|J^\ell\cup J^r|=t_i$, let
\[
  R=I\setminus(J^\ell\cup J^r)
\]
and define
\begin{equation}\label{eq:unique-witness-count}
  Z_{J^\ell,J^r}:=
  \#\set{u\in R:|\supp X_u\cap J^\ell|=1,\ \supp X_u\cap J^r=\varnothing}.
\end{equation}
Because the rows are independent,
\[
  Z_{J^\ell,J^r}\sim\operatorname{Bin}(|R|,\pi_i),
  \qquad
  \pi_i=s_i\frac{\binom{n-t_i}{d-1}}{\binom nd}.
\]
Equivalently,
\begin{equation}\label{eq:unique-witness-mean}
  \E Z_{J^\ell,J^r}
  =|R|\frac{ds_i}{n}\prod_{h=0}^{d-2}\frac{n-t_i-h}{n-1-h}.
\end{equation}
Since $t_i\le n_1\le2n/d^{3/2}$,
\[
  \prod_{h=0}^{d-2}\frac{n-t_i-h}{n-1-h}
  \ge1-\frac{(d-1)(t_i-1)}{n-d+1}
  \ge0.9
\]
for all large $n$. Also $|R|\ge n-L-n_1\ge0.9n$. Thus
\begin{equation}\label{eq:unique-witness-mean-lower}
  \E Z_{J^\ell,J^r}\ge0.8ds_i,
\end{equation}
and Chernoff's inequality gives
\begin{equation}\label{eq:unique-witness-chernoff}
  \Pp\set{Z_{J^\ell,J^r}<\tfrac12ds_i}\le e^{-cds_i}.
\end{equation}
At one scale the number of pairs is at most
\[
  \binom n{t_i}\binom{t_i}{s_i}
  \le\left(\frac{en}{t_i}\right)^{t_i}2^{t_i}.
\]
Here $t_i\le p_{\mathrm{st}}s_i$, including $i=r$, and
\[
  s_i\ge p_{\mathrm{st}}^{r_0}\ge\frac{20L}{d}\ge\frac{10n}{d^4}
\]
for all large $n$. Therefore
\[
  \log\left[\binom n{t_i}\binom{t_i}{s_i}\right]
  \le Cp_{\mathrm{st}}s_i\log d=o(ds_i).
\]
A union bound over the $O(\log n)$ scales gives an event $\mathcal E_0(L)$ such that
\begin{equation}\label{eq:E0-failure}
  \Pp(\mathcal E_0(L)^c)\le e^{-cn/d^3}.
\end{equation}

On $\mathcal E_0(L)$, let $v=x+c\1\in T_{0,i}^{\mathcal K}(L)$. Order the coordinates of $x$ by modulus and take the first $s_i$ coordinates as $J^\ell$ and the next $t_i-s_i$ coordinates as $J^r$. Then
\begin{equation}\label{eq:T0-jump-shift}
  x_{s_i}^*>4d\,x_{t_i}^*,
  \qquad
  |c|\le x_{n_1}^*/10\le x_{t_i}^*/10.
\end{equation}
For every witness row $u$, the row index and all selected coordinates other than the unique $j(u)\in J^\ell$ lie beyond the first $t_i$ coordinates. Using the exact row sum,
\[
  |(A_{\zeta,I}v)_u|
  \ge x_{s_i}^*-(d-1+|\zeta|)x_{t_i}^*-(d+|\zeta|)|c|
  \ge x_{s_i}^*-1.2d\,x_{t_i}^*
  \ge0.7x_{s_i}^*>0
\]
for all large $d$. Hence
\begin{equation}\label{eq:T0-exclusion}
  \ker A_{\zeta,I}\cap T_0^{\mathcal K}(L)=\varnothing
  \qquad\text{on }\mathcal E_0(L).
\end{equation}

\subsection{Uniform equal-block witness events}

Fix a sufficiently small absolute constant $c_w>0$. For an integer $s\le n/(100d)$, let $\mathcal W(s)$ be the event that for every pair of disjoint sets $J^\ell,J^r\subset[n]$ with $|J^\ell|=|J^r|=s$, at least $c_wds$ rows $i\in I$ satisfy
\begin{equation}\label{eq:equal-block-witness}
  |\supp X_i\cap(J^\ell\cup J^r)|=1.
\end{equation}
For a fixed pair, the witness probability equals
\[
  2s\frac{\binom{n-2s}{d-1}}{\binom nd}\ge c\frac{ds}{n},
\]
so the mean number of witnesses is at least $cds$. Chernoff's inequality and
\[
  \#\set{(J^\ell,J^r)}
  \le\binom n{2s}\binom{2s}s
  \le\exp\left(Cs\log\frac{en}{s}\right)
\]
give
\begin{equation}\label{eq:W-s-failure}
  \Pp(\mathcal W(s)^c)
  \le\exp\left[-cds+Cs\log\frac{en}{s}\right].
\end{equation}
We use
\begin{equation}\label{eq:witness-scales}
  s=n_1
  \qquad\text{and}\qquad
  s=m_0:=\floor{\frac n{200d}}.
\end{equation}
At these scales, respectively,
\begin{equation}\label{eq:witness-scale-failures}
  \Pp(\mathcal W(n_1)^c)\le e^{-cn/\sqrt d},
  \qquad
  \Pp(\mathcal W(m_0)^c)\le e^{-cn}.
\end{equation}
These uniform events are kept separate from the fixed-vector small-ball probabilities below; in particular, their failures are not multiplied by net cardinalities.

\subsection{A center-invariant equal-block lemma}

The next statement is the independent-row version of the fixed-vector mechanism in~\cite[Section~4.3]{LLTTY19}.

\begin{lemma}[Equal-block small ball]\label{lem:equal-block}
There is an absolute constant $C_{\mathrm{eb}}\ge1$ such that the following holds. Let $1\le m_1\le m_2\le n-m_1$, and suppose that a deterministic vector $u\in\C^n$ satisfies
\begin{equation}\label{eq:equal-block-separation}
  u_{m_1}^*>\frac23+u_j^*
  \qquad\text{for every }j>m_2.
\end{equation}
Let $w\in\mathcal K$. Put $m_*=\min\set{m_0,m_1}$. The following estimates hold uniformly in $w$:
\begin{enumerate}[label=\textup{(\roman*)}]
\item If $m_1\le m_0$, then
\begin{equation}\label{eq:equal-block-small}
  \Pp\set{\norm{A_{\zeta,I}(u+w)}_2\le c\sqrt{dm_1},\ \mathcal W(m_1)}
  \le e^{-cdm_1}.
\end{equation}
\item If $m_1\ge C_{\mathrm{eb}}m_0$, then
\begin{equation}\label{eq:equal-block-large}
  \Pp\set{\norm{A_{\zeta,I}(u+w)}_2\le c\sqrt{dm_0},\ \mathcal W(m_0)}
  \le\exp\left[-cdm_0\log\left(\frac{dm_1}{n}\right)\right].
\end{equation}
\end{enumerate}
\end{lemma}

\begin{proof}
For part~(i), set $q_0=1$ and $m_*=m_1$. For part~(ii), put
\[
  q_0=\floor{\frac{m_1}{m_0}},
  \qquad
  m_*=m_0.
\]
In the second case $q_0m_*\ge m_1/2$. Partition $q_0m_*$ of the largest coordinates into disjoint sets $J_1^\ell,\ldots,J_{q_0}^\ell$, and the same number of coordinates beyond $m_2$ into disjoint sets $J_1^r,\ldots,J_{q_0}^r$, all of size $m_*$. Set $J_q=J_q^\ell\cup J_q^r$, and write
\[
  \mathscr J=\set{(J_q^\ell,J_q^r):1\le q\le q_0}.
\]
Define the selected-block event
\[
  \mathcal W_{\mathscr J}
  :=\bigcap_{q=1}^{q_0}
  \set{\#\set{i\in I:|\supp X_i\cap J_q|=1}\ge c_wdm_*}.
\]
By definition, $\mathcal W(m_*)\subset\mathcal W_{\mathscr J}$. It is therefore enough to prove the asserted bounds with $\mathcal W_{\mathscr J}$ in place of $\mathcal W(m_*)$.

On $\mathcal W_{\mathscr J}$, if
\[
  b_i=\#\set{q:|\supp X_i\cap J_q|=1},
\]
then
\begin{equation}\label{eq:sum-bi}
  \sum_{i\in I}b_i\ge c_wdm_*q_0.
\end{equation}
In the first case, let $I_0$ be the witness rows for the sole pair. Then $|I_0|\ge c_wdm_1$ and $b_i=1$ for every $i\in I_0$. In the second case, use $b_i\le q_0$ and set
\[
  a_0=\frac{c_w}{4}\frac{dm_*}{n}q_0,
  \qquad
  I_0=\set{i\in I:b_i\ge a_0}.
\]
Then
\[
  c_wdm_*q_0\le\sum_{i\in I}b_i\le|I_0|q_0+na_0,
\]
so $|I_0|\ge(3c_w/4)dm_*$. Moreover, $q_0m_*\ge m_1/2$ gives $a_0\ge(c_w/8)dm_1/n$. Thus, after adjusting the absolute constant,
\begin{equation}\label{eq:I0-bi}
  |I_0|\ge cdm_*,
  \qquad
  b_i\ge c\frac{dm_1}{n}
  \quad(i\in I_0).
\end{equation}

Condition on the counts $Q_{i,q}=|\supp X_i\cap J_q|$ for all $i\in I$ and $q\le q_0$, on all selected coordinates outside $\bigcup_qJ_q$, and on all positions in blocks for which $Q_{i,q}\ne1$. The local event $\mathcal W_{\mathscr J}$, the values $b_i$, and the set $I_0$ are measurable with respect to this information. Conditional on it, whenever $Q_{i,q}=1$, the unique coordinate $U_{i,q}$ is uniform on $J_q^\ell\cup J_q^r$; these singleton positions are mutually independent across the disjoint blocks and across rows. By~\eqref{eq:equal-block-separation}, a disk of radius $1/3$ cannot contain coefficient values from both halves, so one witnessed block has concentration at most $1/2$. Identifying $\C$ with $\R^2$, the two-dimensional Esseen concentration inequality~\cite[Proposition~2.2]{LLTTY19} applied to the $b_i$ independent witnessed blocks gives
\begin{equation}\label{eq:block-Esseen}
  \sup_{\lambda\in\C}
  \Pp\set{\left|\sum_{q:Q_{i,q}=1}u_{U_{i,q}}-\lambda\right|<\frac13}
  \le\frac C{\sqrt{b_i}},
\end{equation}
with the sharper bound $1/2$ when $b_i=1$. Every remaining row contribution, including $-\zeta(u_i+w_i)$ and the exact row-sum contribution of $w$, is conditionally deterministic and merely changes the center.

If the Euclidean norm is below a sufficiently small multiple of $\sqrt{dm_*}$, at least $0.9|I_0|$ retained coordinates have modulus less than $1/3$. Conditional independence across rows and a binomial large-deviation bound give $e^{-c|I_0|}$ in the first case. In the second case,~\eqref{eq:I0-bi} implies
\[
  p_*\le C\sqrt{\frac n{dm_1}},
\]
and $C_{\mathrm{eb}}$ is chosen sufficiently large that $p_*\le1/2$ and $\log(dm_1/n)$ is bounded below by a positive absolute constant. Hence the binomial estimate is
\[
  \exp[-c|I_0|\log(1/p_*)]
  \le\exp\left[-cdm_0\log\left(\frac{dm_1}{n}\right)\right].
\]
Integrating the conditional estimate proves the same bounds with $\mathcal W_{\mathscr J}$; the inclusion $\mathcal W(m_*)\subset\mathcal W_{\mathscr J}$ then gives~\eqref{eq:equal-block-small} and~\eqref{eq:equal-block-large}.
\end{proof}

\subsection{\texorpdfstring{The shifted $T_1$ and $T_2$ classes}{The shifted T1 and T2 classes}}\label{subsec:shifted-T1T2}

Normalize the nonconstant component of $T_i^{\mathcal K}(L)$ by $x_{n_i}^*=1$, $i=1,2$. The nets in~\cite[Lemma~4.5]{LLTTY19} have the form
\[
  \mathcal N_i=\mathcal N_i'+\mathcal N_i'',
  \qquad
  \mathcal N_i''\subset\mathcal K,
\]
with
\begin{equation}\label{eq:shifted-net-size}
  |\mathcal N_i|\le\exp(7n_{i+1}\log d).
\end{equation}
Every $u\in\mathcal N_i'$ is supported on the first $n_{i+1}$ ordered coordinates, and every normalized $x+c\1\in T_i^{\mathcal K}(L)$ has an approximant $u+w\in\mathcal N_i$ satisfying
\begin{equation}\label{eq:shifted-net-error}
  \norm{x-u}_\infty\le d^{-3/2},
  \qquad
  \norm{c\1-w}_\infty\le d^{-3/2}.
\end{equation}
In particular,
\[
  u_{n_i}^*\ge1-d^{-3/2}>\frac23,
  \qquad
  u_j^*=0\quad(j\ge n_{i+1}).
\]
For the error $h=(x+c\1)-(u+w)$, each row obeys
\[
  |(A_{\zeta,I}h)_j|
  \le(d+|\zeta|)\norm h_\infty
  \le\frac C{\sqrt d},
\]
so
\begin{equation}\label{eq:shifted-net-matrix-error}
  \norm{A_{\zeta,I}h}_2\le C\sqrt{\frac nd}.
\end{equation}

For $T_1^{\mathcal K}(L)$, apply 
\cref{lem:equal-block} with $m_1=n_1$, $m_2=n_2$. Since $dn_1\asymp n/\sqrt d$ and $\sqrt{n/d}=o(\sqrt{dn_1})$, the net union bound, kept separate from the failure of $\mathcal W(n_1)$, gives
\begin{align}
  \Pp\set{\ker A_{\zeta,I}\cap T_1^{\mathcal K}(L)\ne\varnothing}
  &\le\Pp(\mathcal W(n_1)^c)+\exp(7n_2\log d-cdn_1)\notag\\
  &\le e^{-cn/\sqrt d}.
  \label{eq:T1-exclusion}
\end{align}
Here $n_2\log d/(dn_1)=O(\log d/d^{1/6})=o(1)$.

For $T_2^{\mathcal K}(L)$, apply the second part of 
\cref{lem:equal-block} with $m_1=n_2$, $m_2=n_3$. Then $n_2/m_0\asymp d^{1/3}\ge C_{\mathrm{eb}}$ for all sufficiently large $n$, $dm_0\asymp n$, $dn_2/n\asymp d^{1/3}$, and $\sqrt{n/d}=o(\sqrt n)$. By the fixed choice of $a_3$,
\begin{align}
  \Pp\set{\ker A_{\zeta,I}\cap T_2^{\mathcal K}(L)\ne\varnothing}
  &\le\Pp(\mathcal W(m_0)^c)+\exp(7n_3\log d-cn\log d)\notag\\
  &\le e^{-cn}+e^{-c'n\log d}
  \le e^{-c''n}.
  \label{eq:T2-exclusion}
\end{align}

\subsection{The exact almost-constant class}\label{subsec:almost-constant}

Set
\begin{equation}\label{eq:n0-S-t}
  n_0=\floor{\frac n{16d}},
  \qquad
  \mathcal S(t)=\set{x:0<x_{n_0}^*\le t x_{n_3}^*}.
\end{equation}
We first record the deterministic reduction used in~\cite[Lemma~4.11]{LLTTY19}.

\begin{lemma}[Deterministic almost-constant reduction]\label{lem:almost-constant-reduction}
For every $L\in\set{m,m+1}$,
\begin{equation}\label{eq:almost-constant-reduction}
  (\mathcal B(\vartheta_0)\setminus\mathcal S(12))\setminus\set0
  \subset\bigcup_{i=0}^3T_i^{\mathcal K}(L).
\end{equation}
\end{lemma}

\begin{proof}
Let $0\ne x\in\mathcal B(\vartheta_0)\setminus\mathcal S(12)$, and choose $\lambda_0$ such that more than $n-n_3$ coordinates satisfy $|x_j-\lambda_0|\le\vartheta_0x_{n_3}^*$. Put $y=\lambda_0\1$ and $w=x-y$.

First suppose $x_{n_3}^*>0$. Then
\begin{equation}\label{eq:lambda0-bounds}
  (1-\vartheta_0)x_{n_3}^*
  \le|\lambda_0|
  \le(1+\vartheta_0)x_{n_3}^*.
\end{equation}
Since $x\notin\mathcal S(12)$, one has $x_{n_0}^*>12x_{n_3}^*$. Therefore
\[
  w_{n_1}^*\ge w_{n_0}^*
  \ge x_{n_0}^*-|\lambda_0|
  >(11-\vartheta_0)x_{n_3}^*,
\]
while
\[
  w_{n_3}^*\le\vartheta_0x_{n_3}^*.
\]
For all sufficiently large $d$,
\[
  (11-\vartheta_0)x_{n_3}^*
  >10(1+\vartheta_0)x_{n_3}^*
  \ge10|\lambda_0|,
\]
and, because $\vartheta_0=10/d^3$,
\[
  d^3w_{n_3}^*\le10x_{n_3}^*
  <(11-\vartheta_0)x_{n_3}^*
  <w_{n_1}^*.
\]
If $w\notin T_0(L)\cup T_3(L)$ and also $w\notin T_1(L)\cup T_2(L)$, the definitions would give $w_{n_1}^*\le d^{3/2}w_{n_2}^*\le d^3w_{n_3}^*$, a contradiction. Thus $w\in T(L)$. Moreover, $|\lambda_0|\le w_{n_1}^*/10$, so $x=w+y$ belongs to the corresponding shifted class.

Now suppose $x_{n_3}^*=0$. More than $n-n_3$ coordinates of $x$ are zero, and more than $n-n_3$ coordinates equal $\lambda_0$; since $a_3<1/2$, the two sets intersect and $\lambda_0=0$. Thus $w=x$ and $y=0$. If $w_{n_1}^*>0$, exclusion from all four source steep classes would imply $w_{n_1}^*\le d^{3/2}w_{n_2}^*\le d^3w_{n_3}^*=0$, a contradiction. If $w_{n_1}^*=0$, then $w$ is supported on fewer than $n_1$ coordinates. If its support has size below $p_{\mathrm{st}}^{r_0}$, then $w\in T_3(L)$. Otherwise, if the support has size at least $\ceil{n_1/p_{\mathrm{st}}}$, the terminal jump in~\eqref{eq:T0r} is strict because $w_{n_1}^*=0$. If the support is smaller than $\ceil{n_1/p_{\mathrm{st}}}$, there is a first index $i\in\set{r_0,\ldots,r-1}$ for which the support lies below $p_{\mathrm{st}}^{i+1}$; the corresponding jump in~\eqref{eq:T0i} is strict. After accounting for the recursive exclusions, $w$ belongs to the first applicable $T_{0,i}(L)$. Hence $w\in T_0(L)$. In every case $x=w+0\in T_i^{\mathcal K}(L)$ for some $i$.
\end{proof}

We next construct the rectangular row-load event. For $J\subset[n]$ and $i\in I$, write
\[
  H_i(J)=|\supp X_i\cap J|.
\]
Let $\mathcal E_{\mathrm{load}}$ be the event that simultaneously
\begin{equation}\label{eq:load-small-J}
  \#\set{i\in I:H_i(J)\ge1}\le\frac K8
  \qquad\text{for every }|J|\le n_0,
\end{equation}
and
\begin{equation}\label{eq:load-n3-J}
  \#\set{i\in I:H_i(J)\ge16n_3d/n}\le\frac K8
  \qquad\text{for every }|J|\le n_3.
\end{equation}
For fixed $J$ with $|J|\le n_0$, the union bound within one row gives
\[
  \Pp\set{H_i(J)\ge1}\le\frac{d|J|}{n}\le\frac1{16}.
\]
For fixed $J$ with $|J|\le n_3$, the hypergeometric mean satisfies $\E H_i(J)=d|J|/n\le dn_3/n$, so Markov's inequality gives
\[
  \Pp\set{H_i(J)\ge16n_3d/n}\le\frac1{16}.
\]
The rows are independent. A binomial Chernoff bound therefore makes the probability that more than $K/8$ rows are bad at most $e^{-cK}$ for each fixed $J$. Moreover,
\[
  \log\sum_{j\le n_0}\binom nj
  =O\left(\frac{n\log d}{d}\right)=o(n),
  \qquad
  \log\sum_{j\le n_3}\binom nj
  \le nH(a_3)+o(n),
\]
where $H(t)=-t\log t-(1-t)\log(1-t)$. By the fixed choice of $a_3$, $H(a_3)<c/2$, and union bounds over all such $J$ give
\begin{equation}\label{eq:load-failure}
  \Pp(\mathcal E_{\mathrm{load}}^c)\le e^{-cn}.
\end{equation}

Take $x\in\mathcal B(\vartheta_0)\cap\mathcal S(12)$, choose $\lambda_0$ as above, and let $\sigma$ order the coordinates by modulus. Set
\[
  J_1=\set{j:|x_j-\lambda_0|\le\vartheta_0x_{n_3}^*},
  \qquad
  J_2=\sigma([n_0])\setminus J_1,
\]
\[
  J_3=\sigma([n_3])\setminus(J_1\cup J_2),
  \qquad
  J_4=[n]\setminus(J_1\cup\sigma([n_3])).
\]
Then $|J_2|\le n_0$, $|J_3|,|J_4|\le n_3$, and
\[
  |x_j|\le12x_{n_3}^*\quad(j\in J_3),
  \qquad
  |x_j|\le x_{n_3}^*\quad(j\in J_4).
\]
On $\mathcal E_{\mathrm{load}}$, delete rows intersecting $J_2$, rows with load at least $16n_3d/n$ in either $J_3$ or $J_4$, and row indices in $\sigma([n_3])$. The surviving set $G\subset I$ has
\begin{equation}\label{eq:surviving-G}
  |G|\ge K-\frac{3K}{8}-n_3\ge\frac n2.
\end{equation}
For $i\in G$, at least $d-32n_3d/n$ selected coordinates lie in $J_1$. Using~\eqref{eq:lambda0-bounds}, the contribution from $J_1$ has modulus at least
\[
  \left(d-\frac{32n_3d}{n}\right)(1-2\vartheta_0)x_{n_3}^*,
\]
whereas the $J_3\cup J_4$ contribution is at most
\[
  \frac{208n_3d}{n}x_{n_3}^*.
\]
Since $i\notin\sigma([n_3])$, the diagonal term is at most $Z_0\sqrt d\,x_{n_3}^*$. By the fixed choice of $a_3$, and then for all sufficiently large $n$,
\begin{equation}\label{eq:almost-constant-row-lower}
  |(A_{\zeta,I}x)_i|\ge\frac d2x_{n_3}^*
  \qquad(i\in G).
\end{equation}
Thus no vector in $\mathcal B(\vartheta_0)\cap\mathcal S(12)$ belongs to the kernel on $\mathcal E_{\mathrm{load}}$.

Combining 
\cref{lem:almost-constant-reduction} with~\eqref{eq:T0-exclusion}, \eqref{eq:T1-exclusion}, \eqref{eq:T2-exclusion}, and~\eqref{eq:load-failure} yields
\begin{align}
  &\Pp\set{(\ker A_{\zeta,I}\setminus\set0)
  \cap\bigl(\mathcal B(\vartheta_0)\setminus T_3^{\mathcal K}(L)\bigr)\ne\varnothing}
  \notag\\
  &\hspace{3cm}\le e^{-cn/d^3}+e^{-cn/\sqrt d}+e^{-c'n\log d}+e^{-cn}.
  \label{eq:almost-constant-exclusion}
\end{align}
The zero vector is excluded from~\eqref{eq:almost-constant-exclusion}: it belongs to $\mathcal B(\vartheta_0)$, while the compatibility theorem concerns only nonzero kernel vectors.

\subsection{Compatibility with the exact gradual class}

We can now state the precise theorem needed by the profile argument.

\begin{theorem}[Rectangular right-kernel compatibility]\label{thm:kernel-compatibility}
There exists an absolute $C_*>0$, depending only on the fixed constants in the LLTTY decomposition and on $a$, such that the following holds. Fix $\sigma\in(0,1)$ and $Z_0<\infty$. Let $d$ satisfy~\eqref{eq:lower-band}, let $L\in\set{m,m+1}$, let $I\subset[n]$ be deterministic with $|I|=n-L$, and let $|\zeta|\le Z_0\sqrt d$. Then, for all sufficiently large $n$, where the threshold may depend on $\sigma$ and $Z_0$,
\begin{equation}\label{eq:kernel-compatibility}
  \Pp\set{\ker A_{\zeta,I}\setminus\set0\subset V_L\cup\cone{S_L}}
  \ge1-\frac Cn.
\end{equation}
The constants and the probability bound are uniform for $L=m,m+1$.
\end{theorem}

\begin{proof}
Let $\mathcal E_L$ be the event on which
\[
  \ker A_{\zeta,I}\cap T_i^{\mathcal K}(L)=\varnothing,
  \qquad i=0,1,2,
\]
and
\[
  (\ker A_{\zeta,I}\setminus\set0)
  \cap\bigl(\mathcal B(\vartheta_0)\setminus T_3^{\mathcal K}(L)\bigr)
  =\varnothing.
\]
The estimates above give
\begin{equation}\label{eq:EL-failure}
  \Pp(\mathcal E_L^c)
  \le e^{-cn/d^3}+e^{-cn/\sqrt d}+e^{-c'n\log d}+e^{-cn}
  \le\frac Cn.
\end{equation}
The last inequality is uniform in the lower band because $d\le(\log n)^{1+\sigma}$ implies $n/d^3\gg\log n$.

Take $0\ne v\in\ker A_{\zeta,I}$ on $\mathcal E_L$. If $v\in T_3^{\mathcal K}(L)$, then~\eqref{eq:T3-in-VL} gives $v\in V_L$. Otherwise $v\notin\mathcal B(\vartheta_0)$, and $v\notin T(L)$: membership in any unshifted $T_i(L)$ would imply membership in $T_i^{\mathcal K}(L)$ by taking the constant vector equal to zero. Moreover, $v_{n_3}^*>0$, since otherwise $v\in\mathcal B(\vartheta_0)$ with center zero. By homogeneity,
\[
  \frac{v}{v_{n_3}^*}\in S_L.
\]
Thus $v\in\cone{S_L}$, proving~\eqref{eq:kernel-compatibility}.
\end{proof}

\begin{remark}\label{rem:final-Cstar}
The final lower bound on $C_*$ must dominate the constants in the support-diameter estimate, the smallest-singular-value theorem, the rectangular restricted-norm estimate, and the witness and net arguments above. It is not introduced solely by the max-weight lemma. The constant $C_*$ is independent of $\sigma$ and of the fixed shift bound $Z_0$; only the threshold from which the estimates hold may depend on them, because $Z_0\sqrt d=o(d)$.
\end{remark}

\section{The row-slice structural theorem}\label{sec:structural}

We now insert the averaged profile estimate into the deterministic LLTTY decomposition of the exact gradual class $S_L$. Assume~\eqref{eq:lower-band} and set
\begin{equation}\label{eq:structural-parameters-basic}
  m=\floor{\frac n{d^3}}-1,
  \qquad
  L\in\set{m,m+1},
  \qquad
  K=n-L,
  \qquad
  \theta_L=\frac Kn.
\end{equation}
Thus
\begin{equation}\label{eq:theta-L}
  1-\theta_L=\frac Ln=d^{-3}+O(n^{-1}),
  \qquad
  -\log\theta_L=O(d^{-3}).
\end{equation}
The lower-band upper bound gives, uniformly in $d$,
\begin{equation}\label{eq:lower-band-budgets}
  d^2=o(n),
  \qquad
  \frac n{d^3}\gg\log n,
  \qquad
  d^{1/96}=o(\log n).
\end{equation}

Let $\gamma=1/288$ and define the exact codimension-dependent parameters
\begin{equation}\label{eq:u-rho-delta-v}
  u_L=\left(\frac nL\right)^\gamma,
  \qquad
  \rho_L=e^{-u_L},
  \qquad
  \delta_L=\left(\frac Ln\right)^{\gamma/2},
  \qquad
  v_L=\floor{\frac{u_L}{\log d}}.
\end{equation}
Then $\rho_L\le d^{-v_L}$ and, uniformly for $L=m,m+1$,
\begin{equation}\label{eq:delta-v-growth}
  \delta_Lv_L\asymp\frac{d^{1/192}}{\log d}\gg\log d.
\end{equation}
More explicitly,
\[
  \delta_L=d^{-1/192+o(1)},
  \qquad
  v_L\asymp\frac{d^{1/96}}{\log d}.
\]
Consequently,
\[
  \cP a_3(v_L-4)\asymp\frac{d^{1/96}}{\log d},
\]
whereas
\[
  2\log_2d+2\log_2\left(\frac{72\sqrt d}{\delta_L}\right)
  +2-\log_2(\cP a_3)=O(\log d).
\]
Thus, for all sufficiently large $n$,
\begin{equation}\label{eq:Pv-hypotheses}
  \begin{gathered}
  0<\delta_L\le36\cP a_3,\\
  \cP a_3(v_L-4)
  \ge2\log_2d+2\log_2\left(\frac{72\sqrt d}{\delta_L}\right)
  +2-\log_2(\cP a_3).
  \end{gathered}
\end{equation}
Hence the hypotheses in part~(v) of 
\cref{prop:LLTTY-package} hold at $(v,\rho,\delta)=(v_L,\rho_L,\delta_L)$. By~\eqref{eq:lower-band-budgets},
\[
  d^{v_L}=\exp(O(d^{1/96}))=n^{o(1)},
  \qquad
  \frac{\sqrt n}{d^{3/2}}=n^{1/2-o(1)}.
\]
In particular, $d\le n^{1/3}$ and $d^{v_L}\le\sqrt n/d^{3/2}$ for all sufficiently large $n$, so all LLTTY net and equivalence-class bounds apply for $4\le u\le v_L$.

For later reference, define the clustered terminal class
\begin{equation}\label{eq:clustered-Pv}
  P_{v_L,\rho_L,\delta_L}(L)
  :=\set{x\in P_{v_L}(L):\exists\lambda\in\C,\ 
  \#\set{j:|x_j-\lambda|\le\rho_Lx_{n_3}^*}>\delta_Ln}.
\end{equation}

\subsection{The net reduction}

The following is the rectangular form of the deterministic net reduction in~\cite[Sections~5--7]{LLTTY19}. The only matrix estimate needed in passing from a continuous vector to its net point is the rectangular restricted norm from 
\cref{prop:restricted-norm}.

\begin{lemma}[Row-slice net reduction]\label{lem:net-reduction}
Let $\mathcal X$ be one of the families $K_u(L)$ or $P_{v_L,\rho_L,\delta_L}(L)$, and let $k=d^u$ be its discretization scale. There is a deterministic set $\mathcal A_k(\mathcal X)$ of $k$-vectors such that
\begin{equation}\label{eq:net-reduction}
  \begin{aligned}
  \Pp\set{\exists x\in\mathcal X:A_{\zeta,I}x=0}
  &\le2d^2\sum_{y\in\mathcal A_k(\mathcal X)}
  \sup_{V\in\C^K}\\[-2pt]
  &\hspace{3em}\times
  \Pp\set{\norm{M_Iy+V}_2\le C\log d\frac{\sqrt{dn}}{k}}
  +n^{-100}.
  \end{aligned}
\end{equation}
The $k$-vectors represented in $\mathcal A_k(\mathcal X)$ fall into at most $e^n$ LLTTY equivalence classes, and each represented class obeys~\eqref{eq:class-cardinality}. Moreover,
\begin{equation}\label{eq:crude-net-count}
  \left|\bigcup_{u=4}^{v_L}\mathcal A_{d^u}(L)\right|
  \le v_L\exp(Cnd^{1/3}).
\end{equation}
\end{lemma}

\begin{lemmafiveonea}[Coordinate ranges for gradual net points]\label{lem:coordinate-ranges}
Let $x\in S_L$, so $x_{n_3}^*=1$, and let $y$ be its canonical $k$-approximation. If $\sigma$ orders the coordinates of $x$ by nonincreasing modulus, then, for all sufficiently large $n$,
\[
  \max_{j\le n}|y_j|\le2n^6,
  \qquad
  \max_{j>n_2}|y_{\sigma(j)}|\le2d^{3/2}.
\]
\end{lemmafiveonea}

\begin{proof}[Proof of Lemma 5.1A]
Because $x\in S_L\subset T(L)^c$ and $x_{n_3}^*=1$, the deterministic order-statistic estimate in~\cite[Lemma~4.2]{LLTTY19} gives $x_1^*\le n^6x_{n_3}^*=n^6$. Since $x\notin T_2(L)$, the defining reverse inequality gives $x_{n_2}^*\le d^{3/2}x_{n_3}^*=d^{3/2}$. Canonical rounding changes each real and imaginary component by less than $1/k$, and $k\ge1$, so the displayed bounds follow after replacing the constants by $2$.
\end{proof}

\begin{proof}[Proof of 
\cref{lem:net-reduction}]
Work on the restricted-norm event from 
\cref{prop:restricted-norm}. Let $x\in\mathcal X$ satisfy $A_{\zeta,I}x=0$, and let $y$ be its canonical $k$-approximation. Put
\[
  \alpha=\frac1n\sum_{j=1}^n(x_j-y_j),
  \qquad
  w=x-y-\alpha\1.
\]
Then $w\perp\1$, while coordinatewise rounding gives
\begin{equation}\label{eq:net-w-alpha}
  \norm w_2\le C\frac{\sqrt n}{k},
  \qquad
  |\alpha|\le\frac Ck.
\end{equation}
Round the real and imaginary parts of $\alpha$ to the grid $(kd)^{-1}\Z$, obtaining $\wt\alpha$. There are at most $Cd^2$ possible values and
\begin{equation}\label{eq:alpha-rounding}
  |\alpha-\wt\alpha|\le\frac C{kd}.
\end{equation}
Writing $y_I=(y_i)_{i\in I}$, define
\[
  V(y,\wt\alpha)=-\zeta y_I+\wt\alpha(d-\zeta)\1_I.
\]
The exact row-sum identity $M_I\1=d\1_I$ and $A_{\zeta,I}x=0$ give
\[
  M_Iy+V(y,\wt\alpha)
  =-(M_I-\zeta P_I)w+(\wt\alpha-\alpha)(d-\zeta)\1_I.
\]
Since $w\perp\1$, one has $M_Iw=(M_I-\E M_I)w$. Therefore~\eqref{eq:net-w-alpha}, \eqref{eq:alpha-rounding}, and the restricted norm imply
\[
  \norm{M_Iy+V(y,\wt\alpha)}_2
  \le C_{Z_0}\frac{\sqrt{dn}}{k}
  \le C\log d\frac{\sqrt{dn}}{k}.
\]
Taking the union over $O(d^2)$ choices of $\wt\alpha$ proves~\eqref{eq:net-reduction}. We use 
\cref{prop:restricted-norm} with failure $n^{-100}$.

The equivalence-class assertions for the canonical approximations in $\mathcal A_k(\mathcal X)$ are~\cite[Lemmas~5.1 and 5.5]{LLTTY19}; no claim is made here about unrestricted $k$-vectors. For the crude count, first choose the positions of the top $n_2$ coordinates. Lemma 5.1A shows that each top coordinate has $O(k^2n^{12})$ possible complex grid values, while each remaining coordinate has $O(k^2d^3)$ possible values. Hence
\[
  |\mathcal A_k(L)|
  \le\binom n{n_2}(Ck^2n^{12})^{n_2}(Ck^2d^3)^{n-n_2}.
\]
Since $n_2\asymp n/d^{2/3}$ and $\log k=O(d^{1/96})$, the logarithm of the right-hand side is at most $Cnd^{1/3}$. Summing over the $v_L$ discretization scales proves~\eqref{eq:crude-net-count}.
\end{proof}

For a fixed net vector $y$, define $W_y$ and $B$ by~\eqref{eq:Wy} and~\eqref{eq:B-distribution}. Applying~\eqref{eq:profile-tensorization} and~\eqref{eq:condition-to-slice} gives
\begin{equation}\label{eq:fixed-net-slice}
  \Pp_{\mathrm{slice}}
  \set{\norm{M_Iy+V}_2\le C\log d\frac{\sqrt{dn}}{k}}
  \le e^{Cd^2}
  \left[\Lambda_L\bigl(Ce^{-cD(B)}W_y+e^{-cd^{2/3}}\bigr)\right]^K,
\end{equation}
where
\begin{equation}\label{eq:Lambda-L}
  \Lambda_L=C\frac nK\log^2d=C\theta_L^{-1}\log^2d.
\end{equation}
We split the two terms using $(A+B)^K\le2^{K-1}(A^K+B^K)$.

\subsection{The profile-exception term}

By~\eqref{eq:crude-net-count}, the additive term $e^{-cd^{2/3}}$, summed over all $4\le u\le v_L$, contributes at most
\[
  \exp\left[-cKd^{2/3}+Cnd^{1/3}+CK\log\log d+Cd^2+O(n)\right]
  \le e^{-c'nd^{2/3}}.
\]
This single estimate is used for every branch below.

\subsection{Entropy of dyadic weight orders}

\begin{lemma}\label{lem:weight-entropy}
For every integer-valued random variable $B$,
\[
  H(B)\le C+\log(1+D(B)),
  \qquad
  D(B)=\E|B-B'|.
\]
\end{lemma}

\begin{proof}
Let $r\in\Z$ be a median of $B$. Since a median minimizes expected absolute deviation,
\[
  M:=\E|B-r|\le\E|B-B'|=D(B).
\]
For $\alpha\in(0,1]$, compare the law of $B$ by relative entropy with
\[
  q_j=Z_\alpha^{-1}e^{-\alpha|j-r|},
  \qquad
  Z_\alpha\le C/\alpha.
\]
Then $H(B)\le\log Z_\alpha+\alpha M$. Taking $\alpha=(1+D(B))^{-1}$ proves the lemma.
\end{proof}

Let $T_b$ be the dyadic weight set of order $b$, and let $t_b=|T_b|$. By~\eqref{eq:factorial-compression},
\begin{equation}\label{eq:factorial-entropy}
  \frac{n!}{\prod_qs_q!}
  \le e^{Cn}\frac{n!}{\prod_bt_b!}
  \le e^{Cn+nH(B)}.
\end{equation}
Combining 
\cref{lem:weight-entropy} with $K/n\to1$ yields
\begin{equation}\label{eq:entropy-dispersion-cancel}
  \frac{n!}{\prod_qs_q!}e^{-cKD(B)}\le e^{C'n},
\end{equation}
because $\log(1+t)-c(K/n)t$ is uniformly bounded above.

\subsection{\texorpdfstring{The $K_u(L)$ branches}{The Ku(L) branches}}

Let $y$ be the $d^u$-approximation of a vector in $K_u(L)$, and put
\begin{equation}\label{eq:kappa-I-large}
  \kappa=\frac{\cK a_3}{16}<\frac13,
  \qquad
  \mathscr I=\set{q:s_q\ge d^{-\kappa}n}.
\end{equation}
The inequality $\kappa<1/3$ ensures that every part in $\mathscr I$ uses the large-part branch of~\eqref{eq:truncated-weights}.

\subsubsection{Case 1: large parts occupy almost all coordinates}

Assume
\[
  \sum_{q\in\mathscr I}s_q\ge n-\frac{\cK n_3}{4}.
\]
Let $\mathscr I_1$ be the large spread parts, $\mathscr I_2$ the large regular parts, and $\mathscr I_3$ the remaining parts. Then
\[
  \sum_{q\in\mathscr I_1}s_q\ge\frac{3\cK n_3}{4},
  \qquad
  \sum_{q\in\mathscr I_3}s_q\le\frac{\cK n_3}{4}.
\]
For $q\in\mathscr I_1$,
\[
  \wt w_q^{-s_q}
  \le d^{-s_q/2}\left(\frac n{s_q}\right)^{s_q}h_q^{-s_q};
\]
for $q\in\mathscr I_2$,
\[
  \wt w_q^{-s_q}
  =\left(\frac n{s_q}\right)^{s_q}h_q^{-s_q};
\]
and for $q\in\mathscr I_3$,
\[
  \wt w_q^{-s_q}\le d^{s_q}h_q^{-s_q}.
\]
Since $n/s_q\le d^\kappa$ on $\mathscr I$, multiplication gives
\[
  \prod_q\wt w_q^{-s_q}
  \le C^nd^{-\cK n_3/16}\prod_qh_q^{-s_q}.
\]
Thus, for some absolute $c_1>0$,
\begin{equation}\label{eq:Ku-case1-energy}
  \prod_qh_q^{s_q}W_y^n\le C^nd^{-c_1n}.
\end{equation}
Passing from $n$ to $K=n-L$ rows costs only
\begin{equation}\label{eq:Wy-K-n}
  W_y^K=W_y^nW_y^{-L},
  \qquad
  W_y^{-L}\le(n\sqrt d)^L=e^{O(n/d^2)}.
\end{equation}
For one equivalence class, multiply the main term in~\eqref{eq:fixed-net-slice} by~\eqref{eq:class-cardinality}. Using~\eqref{eq:entropy-dispersion-cancel}, \eqref{eq:Ku-case1-energy}, and~\eqref{eq:Wy-K-n}, its logarithm is
\[
  -c_1n\log d+O(n\log\log d+n+d^2+n/d^2).
\]
The $e^n$ equivalence classes and $v_L$ choices of $u$ are absorbed, so
\begin{equation}\label{eq:Ku-case1-prob}
  \Pp\set{\text{some Case-1 $K_u(L)$ net vector is small}}
  \le e^{-cn\log d}.
\end{equation}

\subsubsection{Case 2: a constant fraction lies in small parts}

Assume
\[
  \sum_{q\in\mathscr I^c}s_q\ge\frac{\cK n_3}{4}.
\]
By 
\cref{prop:LLTTY-package}, there are $b_1\le b_2$ such that
\[
  \Pp\set{B\le b_1}\ge c_0,
  \qquad
  \Pp\set{B\ge b_2}\ge c_0,
  \qquad
  b_2-b_1\ge c_1d^\kappa.
\]
Therefore
\begin{equation}\label{eq:Ku-case2-dispersion}
  D(B)\ge2c_0^2(b_2-b_1)\ge cd^\kappa.
\end{equation}
Using $\wt w_q\ge h_q/d$,
\[
  W_y^K\prod_qh_q^{s_q}
  \le d^K\prod_qh_q^{(1-K/n)s_q}
  \le d^Ke^{O(n/d^2)}.
\]
The factorial factor is at most $(3\log n)^n=e^{O(n\log d)}$. Thus the contribution of one class has logarithm
\[
  -cKd^\kappa+O(n\log d+n/d^2+d^2).
\]
Since $d^\kappa\gg\log d$,
\begin{equation}\label{eq:Ku-case2-prob}
  \Pp\set{\text{some Case-2 $K_u(L)$ net vector is small}}
  \le e^{-cnd^\kappa}.
\end{equation}
Combining the two cases and the profile-exception term gives
\begin{equation}\label{eq:Ku-total}
  \sum_{u=4}^{v_L}
  \sum_{y\in\mathcal A_{d^u}(K_u(L))}
  \sup_V\Pp\set{\norm{M_Iy+V}_2\le C\log d\frac{\sqrt{dn}}{d^u}}
  \le e^{-cn\log d}.
\end{equation}

\subsection{\texorpdfstring{The clustered $P_{v_L}(L)$ branch}{The clustered P-vL(L) branch}}

Let $x\in P_{v_L,\rho_L,\delta_L}(L)$, and let $y$ be its $d^{v_L}$-approximation. By part~(v) of 
\cref{prop:LLTTY-package}, with the hypotheses verified in~\eqref{eq:Pv-hypotheses}, there exist $b_1<b_2$ such that
\[
  \Pp\set{B\le b_1}\ge\frac{\delta_L}{36},
  \qquad
  \Pp\set{B\ge b_2}\ge c_0,
  \qquad
  b_2-b_1\ge\frac{\cP a_3}{2}(v_L-4)\ge cv_L.
\]
Consequently,
\begin{equation}\label{eq:Pv-dispersion}
  D(B)\ge c\delta_Lv_L.
\end{equation}
For one equivalence class, use $\wt w_q\ge h_q/d$ in the main term. Multiplication by~\eqref{eq:class-cardinality} leaves
\[
  e^{O(d^2)}(Cd\log^2d)^Ke^{-cK\delta_Lv_L}
  \frac{n!}{\prod_qs_q!}
  \prod_qh_q^{(1-K/n)s_q}.
\]
The factorial factor is $e^{O(n\log d)}$, the height mismatch is $e^{O(n/d^2)}$, and $\delta_Lv_L\gg\log d$. Therefore
\begin{equation}\label{eq:Pv-clustered-prob}
  \Pp\set{\text{some clustered $P_{v_L}(L)$ net vector is small}}
  \le e^{-cn\delta_Lv_L}+e^{-cnd^{2/3}}.
\end{equation}

\subsection{Assembly}

By 
\cref{prop:LLTTY-package}, the exact gradual class decomposes as
\begin{equation}\label{eq:SL-decomposition-vL}
  S_L=\bigcup_{u=4}^{v_L}K_u(L)\cup P_{v_L}(L).
\end{equation}
The estimates above show that, with failure $o(n^{-1})$, no vector in $\cone{S_L}$ lying in the right kernel has a cluster of more than $\delta_Ln$ coordinates at radius $\rho_Lx_{n_3}^*$. The compatibility theorem from 
\cref{thm:kernel-compatibility} identifies every remaining nonzero kernel vector with either $V_L$ or $\cone{S_L}$.

\begin{theorem}[Row-slice structural theorem]\label{thm:row-slice-structural}
There are absolute constants $a,C_*>0$ such that the following holds. Fix $\sigma\in(0,1)$, let $d$ satisfy~\eqref{eq:lower-band}, and set
\[
  m=\floor{n/d^3}-1,
  \qquad
  L\in\set{m,m+1}.
\]
Let $I\subset[n]$ be deterministic with $|I|=n-L$. Put $q_L=\floor{aL}$, and let
\[
  E_I=\spanv\set{\ol{X_i-\zeta e_i}:i\in I},
  \qquad
  |\zeta|\le Z_0\sqrt d.
\]
Then, with probability at least $1-C/n$, every nonzero $x\in E_I^\perp=\ker A_{\zeta,I}$ satisfies one of the following:
\begin{enumerate}[label=\textup{(\roman*)}]
\item there exists $i\le q_L$ such that
\[
  x_i^*>0.9\left(\frac ni\right)^3x_{q_L}^*;
\]
\item for all $i\le q_L$ the reverse inequality holds, and
\[
  \sup_{\lambda\in\C}
  \#\set{j:|x_j-\lambda|
  \le\exp\left[-2\left(\frac nL\right)^{1/288}\right]x_{q_L}^*}
  \le n\left(\frac Ln\right)^{1/576}.
\]
\end{enumerate}
The statement is for each fixed $I$, not simultaneously over all row sets. Its constants are uniform in $d$ throughout~\eqref{eq:lower-band} and for $L=m,m+1$; the threshold from which it holds may depend on $\sigma$ and $Z_0$.
\end{theorem}

\begin{proof}
By 
\cref{thm:kernel-compatibility}, outside an event of probability $C/n$, every nonzero kernel vector is in $V_L$ or in $\cone{S_L}$. The first class is alternative~(i). Consider a vector in the second class that is not in $V_L$. The deterministic order-statistic estimates built into the LLTTY gradual class~\cite[Lemma~4.2]{LLTTY19} give
\begin{equation}\label{eq:qL-n3-order-stat}
  x_{q_L}^*
  \le d^3\left(\frac n{q_L}\right)^6x_{n_3}^*
  \le C_ad^{21}x_{n_3}^*.
\end{equation}
Since $u_L/\log d\to\infty$,
\[
  C_ad^{21}\le e^{u_L}
\]
for all large $n$. Hence
\begin{equation}\label{eq:radius-comparison}
  e^{-2u_L}x_{q_L}^*
  \le e^{-u_L}x_{n_3}^*
  =\rho_Lx_{n_3}^*.
\end{equation}
A cluster forbidden by alternative~(ii) would therefore be a cluster in $P_{v_L,\rho_L,\delta_L}(L)$, which is excluded by~\eqref{eq:Ku-total}, \eqref{eq:Pv-clustered-prob}, and 
\cref{lem:net-reduction} with failure $o(n^{-1})$. Since $\delta_Ln=n(L/n)^{1/576}$, the stated bound follows.
\end{proof}

\section{Gaussian normals and exact-slice anti-concentration}\label{sec:normals}

The structural theorem is converted into anti-concentration for a random normal. This stage is model-independent once the row span has the structural property.

\subsection{The Gaussian cluster theorem}

Let $E\subset\C^n$ be a subspace and let $G$ be a standard complex Gaussian vector. We write
\[
  Y=P_{E^\perp}G.
\]
We use the following fixed measurable enlargement rule. If a subspace $F\subset\C^n$ has dimension below a prescribed target $r$, set $F_0=F$ and recursively choose the least index $j$ for which $P_{F_{k-1}^\perp}e_j\ne0$, normalize that projection, and adjoin it to $F_{k-1}$. Stop at dimension $r$, and denote the resulting subspace by $\Ext_r(F)$. Orthogonal projection, the least-index choice, and Gram--Schmidt normalization are Borel measurable on each rank stratum, so $F\mapsto\Ext_r(F)$ is measurable. It contains $F$. Consequently, if $F\subset E_0$, then $E_0^\perp\subset F^\perp$, every structural statement valid on $F^\perp$ remains valid on $E_0^\perp$, and
\[
  \dist(R,E_0)\le\dist(R,F)
\]
for every row $R$. Thus rank deficiency creates no obstruction.

The following is the fixed-subspace content used from~\cite[Theorem~4.2 and Lemmas~4.8--4.10]{LLTTY21}. More precisely, 
\cite[Lemma~4.10]{LLTTY21} supplies the deterministic implication from an oversized correlation-cluster exceptional set to a forbidden vector in $E^\perp$, while~\cite[Lemmas~4.8 and 4.9]{LLTTY21} supply the Gaussian anti-clustering estimate after the large clusters have been removed. The proposition below isolates only this conditional fixed-subspace interface: $E$ is deterministic throughout, and all probability is over $G$. No regular-digraph probability estimate is imported at this stage.

\begin{proposition}[Fixed-subspace Gaussian anti-clustering]\label{prop:Gaussian-anti-clustering}
Let $\gamma=1/288$ and
\[
  \frac{n}{\log^{1/\gamma}n}\le L\le\frac n{d^3},
  \qquad
  q_L=\floor{aL},
  \qquad
  a\le c_{4.10}.
\]
Let $E\subset\C^n$ be deterministic with $\dim E=n-L$. Suppose that every nonzero $x\in E^\perp$ satisfies at least one of the following exact alternatives:
\begin{enumerate}[label=\textup{(\roman*)}]
\item there is $i\le q_L$ such that
\[
  x_i^*>0.9\left(\frac ni\right)^3x_{q_L}^*;
\]
\item for every $i\le q_L$ the reverse inequality holds, and
\[
  \sup_{\lambda\in\C}
  \#\set{j:|x_j-\lambda|
  \le e^{-2(n/L)^\gamma}x_{q_L}^*}
  \le n\left(\frac Ln\right)^{\gamma/2}.
\]
\end{enumerate}
Let $G$ be a standard complex Gaussian vector and $Y=P_{E^\perp}G$. Then, with probability at least $1-CL/n$ over $G$, there is $J\subset[n]$ such that
\begin{equation}\label{eq:Gaussian-exceptional-size}
  |J|\le2n\left(\frac Ln\right)^{\gamma/2}
\end{equation}
and
\begin{equation}\label{eq:Gaussian-anti-cluster}
  \sup_{\lambda\in\C}
  \#\set{j\notin J:|Y_j-\lambda|\le e^{-C(n/L)^\gamma}}
  \le L.
\end{equation}
\end{proposition}

\begin{proof}
Only the fixed-subspace Gaussian results of~\cite{LLTTY21} are used. The notation and hypotheses are mapped as follows.
\begin{center}
\small
\setlength{\tabcolsep}{3pt}
\begin{tabular}{p{0.22\linewidth}p{0.26\linewidth}p{0.43\linewidth}}
\toprule
\textbf{Source quantity} & \textbf{Choice here} & \textbf{Verification}\\
\midrule
$m=\dim E^\perp$ & $m=L$ & $L\ge n/\log^{288}n\ge n^{3/4}$ for large $n$\\
number of Gaussian copies & $p_{\src}=\ell_0-1$ & $p_{\src}<n/b\le3n/L\le n^{1/4}$ for large $n$\\
cluster cutoff $b$ & $b=\floor{L/2}$ & $1\le b<L$ and $2b\le L$\\
small-ball parameter $s$ & $s=\tfrac12(L/n)^{3/2}$ & positive and deterministic\\
strong-correlation radius $\alpha$ & $e^{-C_G(n/L)^\gamma}$ & the radius used in the source construction\\
correlation probability $\beta$ & $\tfrac14(L/(4n))^3$ & the source Lemma~4.10 scale\\
exceptional-size cutoff $V$ & $2n(L/n)^{\gamma/2}$ & the source Lemma~4.10 scale\\
Lemma~4.10 index $\ell_0$ & $p_{\src}+1$ & $\ell_0<n/b+1\le4n/L$ for large $L$\\
source steep cutoff & $\floor{c_{4.10}L}$ & $q_L\le\floor{c_{4.10}L}$\\
\bottomrule
\end{tabular}
\end{center}
The hypotheses of~\cite[Propositions~4.4 and 4.7]{LLTTY21} are therefore both satisfied by the displayed stronger bound $p_{\src}\le n^{1/4}$; the dimension requirement $m\ge n^{1/2}$ follows from $L\ge n^{3/4}$.

Call a pair $(i,j)$ $(\alpha,\beta)$-strongly correlated if
\[
  \Pp_G\set{|Y_i-Y_j|\ge\alpha}\le\beta.
\]
Fix once and for all a total order on the subsets of $[n]$ that first compares cardinality and then compares indicator vectors lexicographically. At each stage of the greedy construction, choose the greatest admissible cluster in this order. This is the deterministic tie-break used immediately before~\cite[Lemma~4.8]{LLTTY21}; it makes every cluster $U_\ell$ and the exceptional set below measurable functions of the covariance of $Y$.

Construct the disjoint strong-correlation clusters $U_1,U_2,\ldots$ with this tie-break. Let $\ell_0$ be the first index with $|U_{\ell_0}|\le b$, and put
\[
  J=\bigcup_{\ell<\ell_0}U_\ell.
\]
If $|J|>V$, 
\cite[Lemma~4.10]{LLTTY21} produces a nonzero $x\in E^\perp$ satisfying the reverse steepness inequalities through $\floor{c_{4.10}L}$ and having more than $n(L/n)^{\gamma/2}$ coordinates in a disk of radius
\[
  e^{-2(n/L)^\gamma}x_{\floor{c_{4.10}L}}^*.
\]
Because $q_L\le\floor{c_{4.10}L}$, order-statistic monotonicity gives
\[
  x_{\floor{c_{4.10}L}}^*\le x_{q_L}^*.
\]
The vector is therefore neither in alternative~(i) nor in alternative~(ii), a contradiction. Hence $|J|\le V$ under the structural hypothesis.

Outside $J$, all greedy clusters have size at most $b$. Apply~\cite[Lemmas~4.8 and 4.9]{LLTTY21} with the displayed $\alpha,\beta,b,s$. The probability of a disk containing at least $2b$ coordinates outside $J$ at radius
\[
  r_G=\frac{\alpha s}{\sqrt{16\log(1/\beta)}}
\]
is at most
\[
  \frac{(sn)^2}{2b^2}\le C\frac Ln.
\]
Since $2b\le L$, absence of such a disk implies the cardinality bound in~\eqref{eq:Gaussian-anti-cluster}. Finally, writing $q=n/L\ge2$,
\[
  \log r_G
  =-C_Gq^\gamma-\frac32\log q-O(\log(e+\log q))
  \ge-Cq^\gamma,
\]
so $r_G\ge e^{-C(n/L)^\gamma}$. Anti-clustering at radius $r_G$ implies the claim at the smaller displayed radius.
\end{proof}

Under~\eqref{eq:lower-band},
\[
  L\ge m\ge\frac{n}{2d^3}
  \ge\frac{n}{\log^{3(1+\sigma)+1}n}
  \ge\frac{n}{\log^{288}n}
\]
for all sufficiently large $n$, uniformly for $L=m,m+1$. Thus the lower codimension hypothesis in 
\cref{prop:Gaussian-anti-clustering} is satisfied throughout the lower band. Combining 
\cref{thm:row-slice-structural,prop:Gaussian-anti-clustering} with $L\in\set{m,m+1}$ and $m=n/d^3+O(1)$ gives
\begin{equation}\label{eq:normal-scales}
  \rho:=e^{-Cd^{1/96}},
  \qquad
  |J|\le2nd^{-1/192+o(1)},
\end{equation}
and failure probability $O(L/n)+O(n^{-1})=O(d^{-3})+O(n^{-1})$.

\subsection{Quotienting the constant direction}

For $y\in\C^n$, let
\[
  \mathfrak m(y)=\frac1n\sum_{j=1}^ny_j,
  \qquad
  y^\circ=y-\mathfrak m(y)\1.
\]
If $X$ is a uniform $d$-subset row, then
\[
  X\cdot y=d\mathfrak m(y)+X\cdot y^\circ.
\]
The deterministic term $d\mathfrak m(y)$ changes only the center in the L\'evy concentration function. For a shifted row $R_i=X-\zeta e_i$,
\[
  R_i\cdot y=X\cdot y^\circ+d\mathfrak m(y)-\zeta y_i,
\]
and the last two terms are deterministic after conditioning on $y$ and $i$. Coordinate clusters are likewise invariant under subtracting $\mathfrak m(y)$. Thus all small-ball estimates are naturally statements modulo $\spanv\set{\1}$.

\subsection{A double-counting lemma for an exact row slice}

\begin{lemma}[Slice anti-concentration from coordinate anti-clustering]\label{lem:slice-double-count}
Let $J\subset[n]$, $K=J^c$, and suppose $|K|\ge n/2$. Assume that for some $m\ge1$ and $\rho>0$,
\begin{equation}\label{eq:coordinate-anti-cluster}
  \sup_{\lambda\in\C}\#\set{j\in K:|y_j-\lambda|\le\rho}\le m.
\end{equation}
If $X$ is the indicator of a uniform $d$-subset of $[n]$ and $d\le n/2$, then
\begin{equation}\label{eq:slice-anti-concentration}
  \sup_{\lambda\in\C}\Pp\set{|X\cdot y-\lambda|\le\rho}
  \le\frac{4m}{n-d+1}+e^{-cd}.
\end{equation}
\end{lemma}

\begin{proof}
Let $S$ be the random $d$-subset and put $H(S)=|S\cap K|$. Since $|K|\ge n/2$, hypergeometric concentration gives
\begin{equation}\label{eq:hypergeometric-H}
  \Pp\set{H<d/4}\le e^{-cd}.
\end{equation}
Fix $\lambda\in\C$ and let
\[
  \mathcal A=\set{S:|S|=d,\ \left|\sum_{j\in S}y_j-\lambda\right|\le\rho}.
\]
Double-count pairs $(S,t)$ with $S\in\mathcal A$ and $t\in S\cap K$:
\[
  \sum_{|S|=d}\1_{\mathcal A}(S)H(S)
  =\sum_{|T|=d-1}\sum_{\substack{t\in K\\t\notin T}}
  \1\set{\left|y_t+\sum_{j\in T}y_j-\lambda\right|\le\rho}.
\]
For each fixed $T$, the inner sum is at most $m$ by~\eqref{eq:coordinate-anti-cluster}. Therefore
\[
  \sum_{|S|=d}\1_{\mathcal A}(S)H(S)
  \le m\binom n{d-1}.
\]
On $H(S)\ge d/4$,
\[
  \frac d4\#\set{S\in\mathcal A:H(S)\ge d/4}
  \le m\binom n{d-1}.
\]
Dividing by $\binom nd$ and using~\eqref{eq:hypergeometric-H} gives
\[
  \Pp(\mathcal A)
  \le\frac{4m}{d}\frac{\binom n{d-1}}{\binom nd}+e^{-cd}
  =\frac{4m}{n-d+1}+e^{-cd}.
\]
Taking the supremum over $\lambda$ proves the lemma.
\end{proof}

At $m=n/d^3+O(1)$, the right-hand side of~\eqref{eq:slice-anti-concentration} is $O(d^{-3})$.

\subsection{The averaged random-normal small-ball bound}

Fix a deterministic exposed row set $I$ of one of the two required cardinalities, and let
\[
  F_I=\spanv\set{\ol{X_i-\zeta e_i}:i\in I}
\]
be the actual conjugated row span. Define
\[
  E_I=\Ext_{|I|}(F_I),
  \qquad
  Y=P_{E_I^\perp}G,
\]
where $G$ is an independent standard complex Gaussian. Let
\[
  \mathcal F_I=\sigma\bigl((X_i)_{i\in I},G\bigr).
\]
Then $E_I$, $Y$, the greedy clusters, and the exceptional set $J$ are $\mathcal F_I$-measurable. By~\eqref{eq:kernel-conjugate-span}, $E_I^\perp\subset F_I^\perp=\ker A_{\zeta,I}$, so $Y$ is a Gaussian vector in the right kernel of the exposed shifted matrix. Conjugation is an isometry and preserves coordinate-cluster cardinalities.

A fresh row slice $X$ is independent of $\mathcal F_I$. Conditional on $\mathcal F_I$, 
\cref{thm:row-slice-structural,prop:Gaussian-anti-clustering} give coordinate anti-clustering outside $J$, and 
\cref{lem:slice-double-count} therefore yields
\begin{equation}\label{eq:averaged-normal-small-ball}
  \E\left[
  \sup_{\lambda\in\C}
  \Pp_X\set{|X\cdot Y-\lambda|\le e^{-Cd^{1/96}}\given\mathcal F_I}
  \right]
  \le Cd^{-3}+\frac Cn.
\end{equation}
The structural and Gaussian failure events contribute $O(n^{-1})+O(d^{-3})$, and the concentration function is at most one. The exceptional coordinate set from~\eqref{eq:normal-scales} has size $o(n)$, hence its complement has size at least $n/2$ for large $n$. The estimate is uniform for each deterministic fixed $I$ of the two required cardinalities. No simultaneous event over all choices of $I$ is asserted.

\section{Intermediate singular values}\label{sec:intermediate}

We convert~\eqref{eq:averaged-normal-small-ball} into a lower bound for an intermediate singular value. Fix $z\in\C$ and write
\[
  A_z=M_n-z\sqrt d\,I_n.
\]
Let $N=n-m$ with
\[
  m=\floor{n/d^3}-1,
\]
and expose the first $N$ rows $R_1,\ldots,R_N$ of $A_z$.

For each $i\le N$, let
\[
  F_i=\spanv\set{\ol{R_j}:j\le N,\ j\ne i}
\]
and define the actual leave-one-row-out distance
\[
  D_i=\dist(\ol{R_i},F_i)
  =\dist\bigl(R_i,\spanv\set{R_j:j\le N,\ j\ne i}\bigr).
\]
If $\dim F_i<N-1$, set $E_i=\Ext_{N-1}(F_i)$; otherwise set $E_i=F_i$. Put
\[
  \wt D_i=\dist(\ol{R_i},E_i)\le D_i,
  \qquad
  Y_i=P_{E_i^\perp}G_i,
\]
where $G_i$ is an independent standard complex Gaussian. Conditional on $R_i$ and $E_i$,
\[
  R_i\cdot Y_i
  =\inner{Y_i}{\ol{R_i}}
  =\inner{G_i}{P_{E_i^\perp}\ol{R_i}}
\]
is a centered complex Gaussian with variance $\wt D_i^2$. Hence there are absolute $c,c_0>0$ such that
\begin{equation}\label{eq:Gaussian-distance-implication}
  \wt D_i\le c\rho
  \quad\Longrightarrow\quad
  \Pp_{G_i}\set{|R_i\cdot Y_i|\le\rho\given R_i,E_i}\ge c_0,
\end{equation}
where
\[
  \rho=e^{-C_{\isv}(z)d^{1/96}}
\]
for a sufficiently large shift-dependent constant $C_{\isv}(z)$.

Conditional on $E_i,Y_i$, the omitted row $R_i=X_i-z\sqrt d\,e_i$ remains an independent row slice. The diagonal shift changes only the center of the concentration function, so~\eqref{eq:averaged-normal-small-ball} gives
\begin{equation}\label{eq:row-normal-small-ball}
  \Pp\set{|R_i\cdot Y_i|\le\rho}\le Cd^{-3}+C/n.
\end{equation}
Combining~\eqref{eq:Gaussian-distance-implication} and~\eqref{eq:row-normal-small-ball},
\begin{equation}\label{eq:leave-one-out-distance-prob}
  \Pp\set{D_i\le c\rho}
  \le\Pp\set{\wt D_i\le c\rho}
  \le Cd^{-3}+C/n.
\end{equation}

For each fixed index $i$, the exposed set $\set{1,\ldots,N}\setminus\set i$ is deterministic, so the fixed-$I$ estimate in 
\cref{sec:normals} applies. We do not intersect these $N$ structural events and do not take a union bound over the leave-one-out indices. Instead, the per-index bound is summed only after taking expectations.

Let
\[
  B=\#\set{i\le N:D_i<c\rho}.
\]
Linearity of expectation, not independence, gives
\[
  \E B\le C\frac n{d^3}+C.
\]
Set
\[
  t=\ceil{\frac n{d^2}}.
\]
Markov's inequality yields
\begin{equation}\label{eq:bad-row-Markov}
  \Pp\set{B>t}
  =O\left(\frac1d+\frac{d^2}{n}\right)
  =O(1/d)=o(1).
\end{equation}

On $B\le t$, delete all bad rows and, if necessary, additional arbitrary rows so that the resulting matrix $H$ has $k=N-t$ rows. For every surviving row $h_i$,
\[
  \dist\bigl(h_i,\spanv\set{h_j:j\ne i}\bigr)\ge c\rho,
\]
because deleting rows only shrinks the comparison span. By the negative second-moment identity~\eqref{eq:negative-second-moment},
\begin{equation}\label{eq:NSM-H}
  \sum_{j=1}^ks_j(H)^{-2}\le\frac{k}{c^2\rho^2}.
\end{equation}
Let $j_0=N-2t=k-t$. At least $t+1$ singular values of $H$ are no larger than $s_{j_0}(H)$, and therefore
\[
  \frac{t+1}{s_{j_0}(H)^2}
  \le\sum_{j=1}^ks_j(H)^{-2}.
\]
Using~\eqref{eq:NSM-H},
\begin{equation}\label{eq:intermediate-H}
  s_{j_0}(H)
  \ge c\rho\sqrt{\frac{t+1}{k}}
  \ge\frac{c\rho}{d}.
\end{equation}
Adding the deleted rows replaces the common $n\times n$ column Gram matrix $H^*H$ by
\[
  H^*H+\sum_{u\in\mathcal D}R_u^*R_u,
\]
a positive-semidefinite increment. Hence the relevant singular values can only increase, and
\[
  s_{j_0}(A_z)\ge s_{j_0}(H).
\]
Define
\[
  q_0=m+2t.
\]
Then $j_0=n-q_0$ and
\begin{equation}\label{eq:q0-scale}
  q_0=O\left(\frac n{d^2}\right)
  =o\left(\frac n{d\log d}\right).
\end{equation}
Since the polynomial factor $d^{-1}$ is absorbed into the exponential,~\eqref{eq:intermediate-H} proves the desired estimate.

\begin{theorem}[Lower-band intermediate singular value]\label{thm:intermediate}
Fix $\sigma\in(0,1)$ and let $d$ satisfy~\eqref{eq:lower-band}. Let $M_n$ be the independent row-slice matrix. For every fixed $z\in\C$, there are constants $C_{\isv}(z),C'_{\isv}(z)>0$ and integers $q_0=q_0(n,z)$ with
\[
  q_0\le C'_{\isv}(z)\frac n{d^2}
\]
such that
\[
  \Pp\set{s_{n-q_0}(M_n-z\sqrt d\,I_n)
  \ge e^{-C_{\isv}(z)d^{1/96}}}=1-o(1).
\]
More precisely, the failure probability is $O(1/d)$ uniformly throughout~\eqref{eq:lower-band}; the term $d^2/n$ and the exponentially small net and witness errors are absorbed into this bound.
\end{theorem}

\begin{remark}
The choice
\[
  m=\floor{n/d^3}-1=\frac n{d^3}+O(1)
\]
is tuned to the LLTTY exponent $\gamma=1/288$:
\[
  (n/m)^\gamma=d^{1/96+o(1)},
  \qquad
  (m/n)^{\gamma/2}=d^{-1/192+o(1)}.
\]
The deletion scale $t=n/d^2$ is larger than the expected number $n/d^3$ of bad rows, but remains small enough for the lower-band logarithmic-tail budget.
\end{remark}

\section{Generalized regularized row replacement}\label{sec:replacement}

We prove the comparison estimate used in both the lower-band and superlogarithmic regimes. The argument is a regularized version of the prefix-volume replacement in~\cite[Section~3]{LLY26a}; the codimension is allowed to be any $r=n-k$ with $p(1-p)(r+1)\ge1$.

\subsection{Regularized distances associated with positive contractions}

Let $P$ be an $n\times n$ Hermitian matrix satisfying $0\le P\le I$. For $\eta>0$ and $u\in\C^n$, define
\[
  f_P(x):=\sqrt{\eta^2+\inner{P(x+u)}{x+u}}.
\]
Let $\boldsymbol\xi=(\xi_1,\ldots,\xi_n)$ have i.i.d. $\Ber(p)$ coordinates, let $X$ be a uniform $d$-subset row with $p=d/n$, and put
\[
  \mu_P:=\E f_P(\boldsymbol\xi),
\]
\begin{equation}\label{eq:DP}
  D_P^2:=\eta^2+p(1-p)\tr P
  +\inner{P(u+p\1)}{u+p\1}.
\end{equation}

\begin{lemma}[Convex concentration on a product cube]\label{lem:convex-concentration}
Let $K>0$, let $\zeta_1,\ldots,\zeta_n$ be independent complex random variables with $|\zeta_j|\le K$ almost surely, and let $F:\C^n\to\R$ be convex and $1$-Lipschitz. If $m_F$ is a median of $F(\zeta)$, then, for every $t>0$,
\begin{align}
  \Pp\set{|F(\zeta)-m_F|\ge tK}&\le Ce^{-ct^2},
  \label{eq:convex-median}\\
  \Pp\set{|F(\zeta)-\E F(\zeta)|\ge tK}&\le Ce^{-ct^2}.
  \label{eq:convex-mean}
\end{align}
The constants are absolute.
\end{lemma}

\begin{proof}
The median-centered estimate is~\cite[Theorem~2.1.13]{Tao12}. To record the recentering used below, integrating~\eqref{eq:convex-median} gives
\[
  |\E F(\zeta)-m_F|
  \le\E|F(\zeta)-m_F|
  \le K\int_0^\infty Ce^{-ct^2}\dd t
  \le CK.
\]
If $t\ge2C$, then
\[
  \set{|F(\zeta)-\E F(\zeta)|\ge tK}
  \subset
  \set{|F(\zeta)-m_F|\ge tK/2};
\]
for $0<t<2C$, the estimate follows after increasing the leading constant. This proves~\eqref{eq:convex-mean}.
\end{proof}

\begin{lemma}[Regularized distance concentration]\label{lem:regularized-distance}
Assume $D_P\ge1$. Then
\[
  c_0D_P\le\mu_P\le D_P,
\]
and, for every $t>0$,
\begin{align}
  \Pp\set{|f_P(\boldsymbol\xi)-\mu_P|\ge t}&\le Ce^{-ct^2},
  \label{eq:Bernoulli-distance-tail}\\
  \Pp\set{|f_P(X)-\mu_P|\ge t}&\le C\sqrt d\,e^{-ct^2}.
  \label{eq:slice-distance-tail}
\end{align}
\end{lemma}

\begin{proof}
Write $x=\boldsymbol\xi-p\1$ and $u_0=u+p\1$. Since $\E x=0$ and the coordinates are independent,
\[
  \E f_P(\boldsymbol\xi)^2
  =\eta^2+p(1-p)\tr P+\inner{Pu_0}{u_0}
  =D_P^2.
\]
We write out the positive-contraction variance calculation. Set
\[
  Q_P:=\inner{Px}{x}-p(1-p)\tr P,
  \qquad
  L_P:=2\Re\inner{Px}{u_0}.
\]
Put $\sigma^2=p(1-p)$ and $m_4=\E x_1^4$. Since $x$ is real and $P=P^*$,
\[
  \inner{Px}{x}
  =\sum_jP_{jj}x_j^2+2\sum_{j<\ell}\Re(P_{j\ell})x_jx_\ell.
\]
The centered diagonal terms are orthogonal in $L^2$ to the off-diagonal monomials, and distinct off-diagonal monomials are likewise orthogonal. Hence
\begin{align}
  \Var(Q_P)
  &=(m_4-\sigma^4)\sum_jP_{jj}^2
  +4\sigma^4\sum_{j<\ell}(\Re P_{j\ell})^2\notag\\
  &\le C\sigma^2\sum_{j,\ell}|P_{j\ell}|^2
  =Cp(1-p)\tr(P^2).
  \label{eq:QP-variance}
\end{align}
Here
\[
  m_4=p(1-p)^4+(1-p)p^4=\sigma^2(1-3\sigma^2),
\]
so $0\le m_4-\sigma^4\le\sigma^2$, while $4\sigma^4\le\sigma^2$. Furthermore, self-adjointness under our linear-first convention gives
\[
  \inner{Px}{u_0}=\inner{x}{Pu_0}
  =\sum_jx_j\ol{(Pu_0)_j}.
\]
Consequently,
\begin{align}
  \Var(L_P)
  &=4\sigma^2\sum_j\bigl(\Re(Pu_0)_j\bigr)^2\notag\\
  &\le4\sigma^2\norm{Pu_0}_2^2
  =4p(1-p)\inner{P^2u_0}{u_0}.
  \label{eq:LP-variance}
\end{align}
Since $f_P(\boldsymbol\xi)^2-D_P^2=Q_P+L_P$, it follows that
\begin{align*}
  \Var\bigl(f_P(\boldsymbol\xi)^2\bigr)
  &\le2\Var(Q_P)+2\Var(L_P)\\
  &\le C\left[p(1-p)\tr(P^2)+p(1-p)\inner{P^2u_0}{u_0}\right]\\
  &\le C\left[p(1-p)\tr P+\inner{Pu_0}{u_0}\right]
  \le CD_P^2.
\end{align*}
Here we used $p(1-p)\le1$, $\tr(P^2)\le\tr P$, and $P^2\le P$. Thus, when $D_P\ge1$,
\[
  \E f_P(\boldsymbol\xi)^4
  =D_P^4+\Var\bigl(f_P(\boldsymbol\xi)^2\bigr)
  \le CD_P^4.
\]
H\"older's inequality yields
\[
  D_P^2=\E f_P^2
  \le(\E f_P)^{2/3}(\E f_P^4)^{1/3},
\]
which implies $\mu_P\ge c_0D_P$ for an absolute constant $c_0>0$. The upper bound is Jensen's inequality.

The map $x\mapsto f_P(x)$ is convex and $1$-Lipschitz because $\norm{P^{1/2}}\le1$. Since the Bernoulli coordinates satisfy $|\xi_j|\le1$, 
\cref{lem:convex-concentration}, in its mean-centered form~\eqref{eq:convex-mean}, gives~\eqref{eq:Bernoulli-distance-tail}.

Finally, the row-slice law is the law of $\boldsymbol\xi$ conditioned on $\sum_j\xi_j=d$. The local central-limit bound
\[
  \Pp\set{\sum_j\xi_j=d}\ge\frac c{\sqrt d}
\]
follows directly from Stirling's formula; the explicit computation appears in the proof of~\cite[Corollary~3.5]{LLY26a}. Dividing by this conditioning probability transfers~\eqref{eq:Bernoulli-distance-tail} to~\eqref{eq:slice-distance-tail}.
\end{proof}

\begin{lemma}[One-row logarithmic comparison]\label{lem:one-row-log}
Let $0\le r<n$ be an integer. Assume $\tr P\ge r+1$ and put
\[
  Q=p(1-p)(r+1)\ge1.
\]
Suppose $f_P(\boldsymbol\xi),f_P(X)\le U$ almost surely. Then
\begin{equation}\label{eq:one-row-log}
  \E|\log f_P(X)-\log f_P(\boldsymbol\xi)|
  \le C\sqrt{\frac{\log(ed)}{Q}}
  +C\sqrt d\log\left(\frac U\eta\right)e^{-cQ}.
\end{equation}
\end{lemma}

\begin{proof}
Since $D_P^2\ge Q$ and $D_P\ge1$, 
\cref{lem:regularized-distance} gives $\mu_P\ge c_0\sqrt Q$. Let
\[
  \mathcal G=\set{f_P(X)\ge\mu_P/2,\ f_P(\boldsymbol\xi)\ge\mu_P/2}.
\]
On $\mathcal G$,
\[
  |\log f_P(X)-\log f_P(\boldsymbol\xi)|
  \le\frac2{\mu_P}|f_P(X)-f_P(\boldsymbol\xi)|.
\]
Integrating the tails in~\eqref{eq:Bernoulli-distance-tail}--\eqref{eq:slice-distance-tail},
\[
  \E|f_P(\boldsymbol\xi)-\mu_P|\le C,
\]
while
\[
  \E|f_P(X)-\mu_P|
  \le\int_0^\infty\min(1,C\sqrt d\,e^{-ct^2})\dd t
  \le C\sqrt{\log(ed)}.
\]
Thus the contribution from $\mathcal G$ is at most $C\sqrt{\log(ed)/Q}$.

By 
\cref{lem:regularized-distance},
\[
  \Pp(\mathcal G^c)
  \le C\sqrt d\,e^{-c\mu_P^2}
  \le C\sqrt d\,e^{-cQ}.
\]
Since $\eta\le f_P\le U$, the logarithmic difference on $\mathcal G^c$ is at most $\log(U/\eta)$. This proves~\eqref{eq:one-row-log}.
\end{proof}

\subsection{The hybrid argument}

For a $k\times n$ matrix $H$, define
\[
  \Psi_\eta(H):=\frac12\log\det(HH^*+\eta^2I_k).
\]
Let $R_1,\ldots,R_k$ be independent uniform row slices. Let $\widetilde R_1,\ldots,\widetilde R_k$ be i.i.d. Bernoulli$(p)$ rows, and assume the two families are mutually independent. Let $u_1,\ldots,u_k\in\C^n$ be deterministic.

Define a hybrid $H^{(i)}$ whose first $i$ rows are $R_j+u_j$ and whose remaining rows are $\widetilde R_j+u_j$. At the $i$th replacement, let $K_i$ be the matrix of the other $k-1$ rows and put
\begin{equation}\label{eq:Pi-hybrid}
  P_i=I-K_i^*(K_iK_i^*+\eta^2I_{k-1})^{-1}K_i,
  \qquad
  \widehat P_i:=\ol{P_i}.
\end{equation}
For a candidate $i$th row with coordinate column $v\in\C^n$, let $H_v$ denote the corresponding hybrid matrix. The Schur-complement identity, together with our convention that matrix rows act bilinearly, gives
\begin{align}
  \det(H_vH_v^*+\eta^2I_k)
  &=\det(K_iK_i^*+\eta^2I_{k-1})
  \bigl(\eta^2+v^\top P_i\ol v\bigr)\notag\\
  &=\det(K_iK_i^*+\eta^2I_{k-1})
  \bigl(\eta^2+\inner{\widehat P_iv}{v}\bigr).
  \label{eq:Schur-hybrid}
\end{align}
Indeed, since $P_i$ is Hermitian,
\[
  \inner{\widehat P_iv}{v}
  =v^\top\widehat P_i^\top\ol v
  =v^\top P_i\ol v.
\]
The determinant involving $K_i$ is common to the two hybrids and cancels. Since $\Psi_\eta=\frac12\log\det$, the remaining factor gives
\begin{align}
  \Psi_\eta(H^{(i)})-\Psi_\eta(H^{(i-1)})
  &=\log\sqrt{\eta^2+\inner{\widehat P_i(R_i+u_i)}{R_i+u_i}}\notag\\
  &\quad-\log\sqrt{\eta^2+\inner{\widehat P_i(\widetilde R_i+u_i)}{\widetilde R_i+u_i}}.
  \label{eq:hybrid-increment}
\end{align}
The matrix $P_i$ is a positive contraction. On $\ker K_i$, it is the identity; hence
\begin{equation}\label{eq:trace-Pi}
  \tr P_i\ge\dim\ker K_i\ge n-k+1=r+1.
\end{equation}
Entrywise conjugation preserves Hermitian positivity and the contraction property. Thus $\widehat P_i$ is also a positive contraction; moreover, since $\tr P_i$ is real,
\[
  \tr\widehat P_i=\ol{\tr P_i}=\tr P_i\ge r+1.
\]
Also, for either row type,
\[
  \sqrt{\eta^2+\inner{\widehat P_i(x+u_i)}{x+u_i}}
  \le\sqrt{\eta^2+(\sqrt n+\norm{u_i}_2)^2}.
\]
Conditional on all other rows, $\widehat P_i$ is deterministic and the two candidate rows retain their respective slice and Bernoulli laws. Therefore 
\cref{lem:one-row-log} applies with $P=\widehat P_i$ and
\[
  U=\sqrt{\eta^2+\left(\sqrt n+\max_{i\le k}\norm{u_i}_2\right)^2}.
\]
Telescoping~\eqref{eq:hybrid-increment}, taking absolute values, conditioning, and summing gives the desired result. Notice that no independence among the increments is required.

\begin{proposition}[Generalized regularized row replacement]\label{prop:row-replacement}
Let $p=d/n$, $1\le k<n$, and $r=n-k$. Define $H_{\mathrm{slice}}$ from independent row slices and $H_{\mathrm{iid}}$ from i.i.d. Bernoulli$(p)$ rows, with the same deterministic shifts $u_i$, and assume that the two row families are mutually independent. Put
\[
  Q=p(1-p)(r+1),
  \qquad
  U=\sqrt{\eta^2+\left(\sqrt n+\max_{i\le k}\norm{u_i}_2\right)^2}.
\]
If $Q\ge1$, then
\begin{equation}\label{eq:row-replacement}
  \frac1n\E|\Psi_\eta(H_{\mathrm{slice}})-\Psi_\eta(H_{\mathrm{iid}})|
  \le C\sqrt{\frac{\log(ed)}{Q}}
  +C\sqrt d\log\left(\frac U\eta\right)e^{-cQ}.
\end{equation}
\end{proposition}

\begin{proof}
Apply 
\cref{lem:one-row-log} conditionally to each increment in~\eqref{eq:hybrid-increment}. There are $k\le n$ increments, so division by $n$ leaves the same bound.
\end{proof}

\section{Hermitization and the superlogarithmic bridge}\label{sec:hermitization}

We combine \cref{thm:intermediate,prop:row-replacement} with the smallest-singular-value inputs and the optimal sparse i.i.d. circular law. Throughout this section, fix an arbitrary deterministic integer sequence $d=d(n)$ satisfying the degree assumptions of the result under consideration, and set $p_n=d(n)/n$. Every assertion is made for this fixed sequence. The Bernoulli logarithmic-potential interface below is valid for every fixed $z\ne0$, so no sequence-dependent exceptional set of shifts is needed.

\subsection{A logarithmic-tail lemma}

\begin{lemma}[Negative logarithmic uniform integrability]\label{lem:negative-log-tail}
Let
\[
  \nu_n=\frac1n\sum_{j=1}^n\delta_{s_{j,n}}
\]
be a random probability measure on $[0,\infty)$, where $s_{1,n}\ge\cdots\ge s_{n,n}\ge0$. Suppose that, for a deterministic probability measure $\nu$,
\begin{enumerate}[label=\textup{(\roman*)}]
\item $\nu_n$ converges weakly to $\nu$ in probability;
\item
\[
  \int\log t\,\dd\nu_n(t)
  \xrightarrow{\Pp}
  \int\log t\,\dd\nu(t)\in\R;
\]
\item
\[
  \int t^2\,\dd\nu_n(t)=\bigOP.
\]
\end{enumerate}
Then, for every $\eps>0$,
\begin{equation}\label{eq:negative-log-UI}
  \lim_{T\to\infty}\limsup_{n\to\infty}
  \Pp\set{
    \int_{\set{\logm t>T}}\logm t\,\dd\nu_n(t)>\eps
  }=0.
\end{equation}
Consequently, if $r_n=o(n)$, then
\begin{equation}\label{eq:bottom-r-log-tail}
  \frac1n\sum_{j=n-r_n+1}^n\logm s_{j,n}
  \xrightarrow{\Pp}0.
\end{equation}
\end{lemma}

\begin{proof}
We first control positive logarithms. For $R\ge e$, put
\[
  c_R:=\sup_{t\ge R}\frac{\log t}{t^2},
\]
so $c_R\to0$. Then
\[
  \int_{t>R}\log t\,\dd\nu_n(t)
  \le c_R\int t^2\,\dd\nu_n(t),
\]
which tends to zero in the iterated limit $n\to\infty$, $R\to\infty$, in probability. Weak convergence applied to $\min(\logp t,\log R)$ therefore gives
\begin{equation}\label{eq:positive-log-convergence}
  \int\logp t\,\dd\nu_n(t)
  \xrightarrow{\Pp}
  \int\logp t\,\dd\nu(t).
\end{equation}

The bounded-in-probability second moments also imply
\[
  \int t^2\,\dd\nu(t)<\infty.
\]
Indeed, for each $L<\infty$, weak convergence in probability gives
\[
  \int(t^2\wedge L)\,\dd\nu_n(t)
  \xrightarrow{\Pp}
  \int(t^2\wedge L)\,\dd\nu(t).
\]
Choose $M<\infty$ such that
\[
  \limsup_{n\to\infty}
  \Pp\set{\int t^2\,\dd\nu_n(t)>M}<\frac12.
\]
Because the truncated integral is bounded above by the full second moment, its deterministic probability limit cannot exceed $M$. Hence
\[
  \int(t^2\wedge L)\,\dd\nu(t)\le M
\]
for every $L$, and monotone convergence proves the claim.

Since $\log t=\logp t-\logm t$, assumption~(ii) and \eqref{eq:positive-log-convergence} imply
\begin{equation}\label{eq:negative-log-total}
  \int\logm t\,\dd\nu_n(t)
  \xrightarrow{\Pp}
  \int\logm t\,\dd\nu(t)<\infty.
\end{equation}
For fixed $T$, the function
\[
  f_T(t)=\min(\logm t,T),
  \qquad f_T(0)=T,
\]
is bounded and continuous. Thus
\[
  \int f_T\,\dd\nu_n
  \xrightarrow{\Pp}
  \int f_T\,\dd\nu.
\]
Subtracting from \eqref{eq:negative-log-total} yields
\[
  \int(\logm-f_T)\,\dd\nu_n
  \xrightarrow{\Pp}
  \int(\logm-f_T)\,\dd\nu.
\]
The deterministic right-hand side tends to zero by monotone convergence. On $\set{\logm t>2T}$,
\[
  \logm t\le2(\logm t-T)=2(\logm t-f_T(t)).
\]
This proves \eqref{eq:negative-log-UI}, after renaming $2T$ as $T$.

Finally, for fixed $T$,
\[
  \frac1n\sum_{j=n-r_n+1}^n\logm s_{j,n}
  \le T\frac{r_n}{n}
  +\int_{\set{\logm t>T}}\logm t\,\dd\nu_n(t).
\]
Let $n\to\infty$ and then $T\to\infty$.
\end{proof}

\subsection{The Bernoulli hermitization interface}

Let $B_n$ denote an $n\times n$ matrix with i.i.d. $\Ber(d/n)$ entries. For fixed $z\in\C$, write
\[
  H_{n,z}:=\frac{B_n}{\sqrt d}-zI_n,
  \qquad
  \nu^B_{n,z}:=\frac1n\sum_{j=1}^n\delta_{s_j(H_{n,z})}.
\]
Let $\nu_z^G$ denote the deterministic limiting singular-value law of
$n^{-1/2}G_n-zI_n$, where $G_n$ has i.i.d. standard complex Gaussian entries, and put
\begin{equation}\label{eq:Gaussian-log-potential}
  U^\circ(z):=-\int_\C\log\abs{z-w}\,\dd\mu_{\mathrm{circ}}(w)
  =-\int_0^\infty\log t\,\dd\nu_z^G(t).
\end{equation}

\begin{proposition}[Bernoulli hermitization interface]\label{prop:Bernoulli-hermitization}
Assume $d\to\infty$ and $d=o(n)$. For every fixed $z\in\C\setminus\set{0}$:
\begin{enumerate}[label=\textup{(\roman*)}]
\item $\nu^B_{n,z}\Rightarrow\nu_z^G$ in probability;
\item
\begin{equation}\label{eq:Bernoulli-full-log}
  -\int_0^\infty\log t\,\dd\nu^B_{n,z}(t)
  \xrightarrow{\Pp}U^\circ(z),
\end{equation}
equivalently, the full untruncated logarithmic integral converges to the finite value
$\int\log t\,\dd\nu_z^G(t)$;
\item
\begin{equation}\label{eq:Bernoulli-second-moment}
  \int_0^\infty t^2\,\dd\nu^B_{n,z}(t)=\bigOP.
\end{equation}
\end{enumerate}
\end{proposition}

\begin{proof}
Part~(i) is \cite[Fact~10.2]{SSS25}, applied to the deterministic sequence $p_n=d/n$ and to the Bernoulli model. The Gaussian limit in that fact is $\nu_z^G$. By \cite[Fact~10.1]{SSS25}, the measure $\nu_z^G$ is compactly supported and satisfies
\[
  \nu_z^G([0,t))\le Ct,
  \qquad t>0,
\]
and \eqref{eq:Gaussian-log-potential} holds. The mass bound at zero and compact support imply that the limiting logarithmic integral is finite.

We reproduce the squeeze that gives the full, untruncated convergence in part~(ii). Define
\begin{equation}\label{eq:Bernoulli-Un}
  U_n(z):=-\frac1n\sum_{j=1}^n\log s_j(H_{n,z}),
\end{equation}
with the extended value $+\infty$ when $H_{n,z}$ is singular. Choose a deterministic sequence $\eps_n\downarrow0$ sufficiently slowly for the estimates in Sections~8--10 of \cite{SSS25}, and write $\eps=\eps_n$. Set
\[
  m=\floor{(1-\eps)n},
  \qquad
  \ell=\floor{(1-\eps/4)m},
  \qquad
  a=\ceil{2\eps n},
  \qquad
  b=\floor{(1-\eps)n}.
\]
Let $H_{m,z}$ be the top-left $m\times m$ principal minor of $H_{n,z}$. Thus its entries retain the normalization $d^{-1/2}$ used for $H_{n,z}$. Define
\begin{align}
  T_n(z)&:=-\frac1n\sum_{j=1}^{\ell}\log s_j(H_{m,z}),
  \label{eq:SSS-T}\\
  T_n^{(1)}(z)&:=-\frac1n\sum_{j=1}^{\ell}\log s_j(H_{n,z}),
  \label{eq:SSS-T1}\\
  T_n^{(2)}(z)&:=-\frac1n\sum_{j=a}^{b}\log s_j(H_{n,z})
  -\frac{\ell-\floor{(1-3\eps)n}}{n}\log s_{\ell}(H_{m,z}).
  \label{eq:SSS-T2}
\end{align}
These are the integer-rounded versions of the quantities in Sections~8--10 of \cite{SSS25}.

For every fixed $z\ne0$, \cite[Lemma~8.2]{SSS25} gives
\begin{equation}\label{eq:SSS-upper-squeeze}
  U_n(z)\le T_n(z)+\littleoP.
\end{equation}
The first deterministic comparison in \cite[Fact~9.1]{SSS25} gives
\begin{equation}\label{eq:SSS-T-upper}
  T_n(z)\le T_n^{(2)}(z).
\end{equation}
For completeness, we derive the lower comparison directly from monotonicity of the ordered singular values. Put
\[
  x_j=-\log s_j(H_{n,z}),\qquad 1\le j\le n.
\]
Since the singular values are nonincreasing, $x_1\le\cdots\le x_n$. Hence
\begin{equation}\label{eq:SSS-U-lower-exact}
  U_n(z)=\frac1n\sum_{j=1}^n x_j
  \ge\frac1\ell\sum_{j=1}^{\ell}x_j
  =\frac n\ell T_n^{(1)}(z).
\end{equation}
Moreover,
\[
  \frac\ell n=(1-\eps/4)(1-\eps)+O(n^{-1}).
\]
Finally, \cite[Lemma~9.3]{SSS25}, proved there from Facts~10.1--10.2 and Lemmas~10.3--10.5, yields
\begin{equation}\label{eq:SSS-truncated-convergence}
  T_n^{(1)}(z)\xrightarrow{\Pp}U^\circ(z),
  \qquad
  T_n^{(2)}(z)\xrightarrow{\Pp}U^\circ(z)
\end{equation}
when $\eps_n\downarrow0$ sufficiently slowly. Since $n/\ell\to1$, \eqref{eq:SSS-U-lower-exact} and the first convergence in \eqref{eq:SSS-truncated-convergence} give
\[
  U_n(z)\ge U^\circ(z)-\littleoP.
\]
On the other hand, \eqref{eq:SSS-upper-squeeze}, \eqref{eq:SSS-T-upper}, and the second convergence in \eqref{eq:SSS-truncated-convergence} give
\[
  U_n(z)\le U^\circ(z)+\littleoP.
\]
Therefore $U_n(z)\to U^\circ(z)$ in probability. In particular,
\[
  \Pp\set{H_{n,z}\text{ is singular}}\to0.
\]
Redefining $U_n(z)$ arbitrarily on that event does not change the conclusion. This proves \eqref{eq:Bernoulli-full-log} for every fixed $z\ne0$.

For part~(iii),
\begin{align*}
  \int t^2\,\dd\nu^B_{n,z}(t)
  &=\frac1n\norm{H_{n,z}}_{\HS}^2\\
  &\le \frac{2}{nd}\sum_{i,j=1}^n(B_n)_{ij}+2\abs z^2.
\end{align*}
The sum has mean $nd$ and is $O_{\Pp}(nd)$, proving \eqref{eq:Bernoulli-second-moment}.
\end{proof}

\begin{corollary}[Bernoulli bottom tail]\label{cor:Bernoulli-bottom-tail}
Fix $z\in\C\setminus\set{0}$. For every deterministic integer sequence $r_n=o(n)$,
\begin{equation}\label{eq:Bernoulli-bottom-tail}
  \frac1n\sum_{j=n-r_n+1}^n
  \logm s_j\!\left(\frac{B_n}{\sqrt d}-zI_n\right)
  =\littleoP.
\end{equation}
\end{corollary}

\begin{proof}
Apply \cref{lem:negative-log-tail} to the three conclusions of
\cref{prop:Bernoulli-hermitization}.
\end{proof}

\subsection{Lower-band regularized comparison of a large prefix}

Fix $\sigma\in(0,1)$, assume \eqref{eq:lower-band}, fix $z\in\C\setminus\set{0}$, and define the unnormalized shifted matrices
\begin{equation}\label{eq:shifted-comparators}
  A=M_n-z\sqrt d\,I_n,
  \qquad
  \wt A=B_n-z\sqrt d\,I_n.
\end{equation}
Here $B_n$ is independent of $M_n$ and has i.i.d. $\Ber(p)$ entries, $p=d/n$.

By \cref{prop:LLY-smallest}, for a sufficiently large absolute constant $C_L$,
\begin{equation}\label{eq:lower-tau}
  \tau_n=n^{-C_L\log\log n}
\end{equation}
satisfies
\begin{equation}\label{eq:lower-good-event}
  \Pp(\mathcal L_n(z)^c)=O(d^{-1/2})=o(1),
  \qquad
  \mathcal L_n(z)=\set{s_n(A)\ge\tau_n}\cap\set{s_n(\wt A)\ge\tau_n}.
\end{equation}
The shift restriction in \cref{prop:LLY-smallest} holds because
$\abs z\sqrt d\le\sqrt d\log\log d$ for fixed $z\ne0$ and all sufficiently large $n$.

Choose
\begin{equation}\label{eq:lower-prefix-parameters}
  r=\ceil{\frac n{\sqrt d}},
  \qquad
  k=n-r,
  \qquad
  \eta=\frac{\tau_n}{n}.
\end{equation}
Let $C_M$ and $C_B$ be the first $k$ rows of $A$ and $\wt A$, respectively. In \cref{prop:row-replacement},
\[
  Q_n=p(1-p)(r+1)\asymp\sqrt d,
\]
and the deterministic shifts are $u_i=-z\sqrt d\,e_i$. Thus
\[
  U\le\sqrt{\eta^2+(\sqrt n+\abs z\sqrt d)^2}\le2(1+\abs z)\sqrt n
\]
and
\begin{equation}\label{eq:lower-U-over-eta}
  \log\frac U\eta=O(\log n\,\log\log n)=O(d\log d).
\end{equation}
Consequently,
\begin{align}
  \frac1n\E\abs{\Psi_\eta(C_M)-\Psi_\eta(C_B)}
  &\le O\!\left(d^{-1/4}\sqrt{\log d}\right)
  +O\!\left(d^{3/2}\log d\,e^{-c\sqrt d}\right)
  \notag\\
  &=o(1).
  \label{eq:lower-regularized-comparison}
\end{align}
Markov's inequality gives
\begin{equation}\label{eq:lower-regularized-op}
  \frac1n\abs{\Psi_\eta(C_M)-\Psi_\eta(C_B)}=\littleoP.
\end{equation}

\subsection{Removing the lower-band regularization}

Let $t_1(C)\ge\cdots\ge t_k(C)$ be the singular values of a $k\times n$ prefix. Cauchy interlacing gives
\[
  t_i(C_M)\ge s_{i+r}(A)\ge s_n(A),
\]
and similarly for $C_B$. On $\mathcal L_n(z)$, every prefix singular value is at least $\tau_n$. Hence
\begin{align*}
  0\le\Psi_\eta(C)-\log\vol(C)
  &=\frac12\sum_{i=1}^k\log\left(1+\frac{\eta^2}{t_i(C)^2}\right)\\
  &\le\frac{k\eta^2}{2\tau_n^2}
  \le\frac1{2n}.
\end{align*}
After division by $n$, the error is at most $1/(2n^2)$. Combining this with
\eqref{eq:lower-regularized-op} and \eqref{eq:lower-good-event}, and extending both logarithmic-volume random variables by zero on $\mathcal L_n(z)^c$, gives
\begin{equation}\label{eq:lower-prefix-volume-comparison}
  \frac1n\abs{\log\vol(C_M)-\log\vol(C_B)}=\littleoP.
\end{equation}

\subsection{From lower-band prefix volumes to square determinants}

Let $H$ be an invertible $n\times n$ matrix, and let $C$ be its first $n-r$ rows. Interlacing implies
\[
  \sum_{j=r+1}^n\log s_j(H)
  \le\log\vol(C)
  \le\sum_{j=1}^{n-r}\log s_j(H).
\]
Therefore
\begin{equation}\label{eq:prefix-square-error}
  \abs{\log\abs{\det H}-\log\vol(C)}
  \le\sum_{j=1}^r\logp s_j(H)
  +\sum_{j=n-r+1}^n\logm s_j(H).
\end{equation}
For the positive logarithms, concavity gives
\begin{equation}\label{eq:positive-prefix-tail}
  \frac1n\sum_{j=1}^r\logp s_j(H)
  \le\frac r{2n}\log\left(1+\frac{\norm H_{\HS}^2}{r}\right).
\end{equation}
Since $\norm A_{\HS}^2\le2(1+\abs z^2)nd$ deterministically and
$\norm{\wt A}_{\HS}^2=O_{\Pp}(nd)$,
\begin{equation}\label{eq:lower-positive-tail}
  \frac1n\sum_{j=1}^r\logp s_j(A)
  =O\!\left(\frac{\log d}{\sqrt d}\right)=o(1),
\end{equation}
with the analogous $o_{\Pp}(1)$ estimate for $\wt A$.

For the negative row-slice tail, use bottom-up indexing
$a_j=s_{n-j+1}(A)$. On $\mathcal L_n(z)$ and the event in \cref{thm:intermediate},
\begin{equation}\label{eq:lower-bottom-q0}
  \frac1n\sum_{j=1}^{q_0}\logm a_j
  \le\frac{q_0}{n}\log\frac1{\tau_n}
  =O\!\left(\frac{\log n\,\log\log n}{d^2}\right)=o(1),
\end{equation}
and
\begin{equation}\label{eq:lower-bottom-intermediate}
  \frac1n\sum_{j=q_0+1}^{r}\logm a_j
  \le\frac rn C_{\isv}(z)d^{1/96}
  =O(d^{-47/96})=o(1).
\end{equation}
Thus
\begin{equation}\label{eq:lower-row-bottom-tail}
  \frac1n\sum_{j=n-r+1}^n\logm s_j(A)=\littleoP.
\end{equation}
By \cref{cor:Bernoulli-bottom-tail}, and because multiplication by $\sqrt d\ge1$ can only decrease $\logm$,
\begin{equation}\label{eq:lower-Bernoulli-bottom-tail}
  \frac1n\sum_{j=n-r+1}^n\logm s_j(\wt A)=\littleoP.
\end{equation}
Applying \eqref{eq:prefix-square-error} to $(A,C_M)$ and $(\wt A,C_B)$, and using \eqref{eq:lower-prefix-volume-comparison}, yields
\begin{equation}\label{eq:lower-unnormalized-det-comparison}
  \frac1n\log\abs{\det A}
  -\frac1n\log\abs{\det\wt A}
  \xrightarrow{\Pp}0.
\end{equation}
On the event where a determinant vanishes, define its normalized logarithm arbitrarily; by \eqref{eq:lower-good-event}, this changes a random variable only on an event of probability $o(1)$. Since the common factor $d^{-n/2}$ cancels,
\begin{equation}\label{eq:lower-normalized-det-comparison}
  \frac1n\log\abs{\det\!\left(\frac{M_n}{\sqrt d}-zI_n\right)}
  -\frac1n\log\abs{\det\!\left(\frac{B_n}{\sqrt d}-zI_n\right)}
  \xrightarrow{\Pp}0
\end{equation}
for every fixed $z\ne0$, hence for Lebesgue-almost every fixed shift.

\subsection{The lower-polylogarithmic circular law}

\begin{theorem}[Lower-polylogarithmic regime]\label{thm:lower-band-circular}
Fix $\sigma\in(0,1)$. There exists an absolute constant $C_*>0$ such that, whenever
\[
  C_*\log n\le d\le(\log n)^{1+\sigma},
\]
one has
\[
  \mu_{M_n/\sqrt{d(1-d/n)}}\xrightarrow{\Pp}\mu_{\mathrm{circ}}.
\]
The constant $C_*$ is independent of $\sigma$; the threshold from which the conclusion holds may depend on $\sigma$.
\end{theorem}

\begin{proof}
Set
\[
  X_n=\frac{M_n}{\sqrt d},
  \qquad
  Y_n=\frac{B_n}{\sqrt d}.
\]
Then
\[
  \frac1n\norm{X_n}_{\HS}^2=1
\]
deterministically and
\[
  \E\frac1n\norm{Y_n}_{\HS}^2=1.
\]
Thus the Hilbert--Schmidt condition in the replacement principle
\cite[Theorem~2.1]{TV10} is satisfied. Equation~\eqref{eq:lower-normalized-det-comparison} gives its logarithmic-determinant condition for every fixed $z\ne0$, and therefore for Lebesgue-almost every fixed $z$. No estimate uniform in $z$ and no union bound over shifts is required. Consequently,
\begin{equation}\label{eq:BL-row-Bernoulli}
  d_{\BL}(\mu_{X_n},\mu_{Y_n})\xrightarrow{\Pp}0.
\end{equation}
The optimal sparse i.i.d. circular law \cite{RT19,SSS25} applies because
\[
  p=d/n\to0,
  \qquad
  np=d\to\infty,
\]
and gives
\[
  d_{\BL}(\mu_{Y_n},\mu_{\mathrm{circ}})\xrightarrow{\Pp}0.
\]
The triangle inequality therefore yields
\begin{equation}\label{eq:sqrt-d-circular}
  \mu_{M_n/\sqrt d}\xrightarrow{\Pp}\mu_{\mathrm{circ}}.
\end{equation}

For every square matrix $H$, Schur triangularization gives
\[
  \sum_{j=1}^n\abs{\lambda_j(H)}^2\le\norm H_{\HS}^2.
\]
Consequently,
\[
  \int\abs z\,\dd\mu_H(z)
  \le\left(\frac1n\norm H_{\HS}^2\right)^{1/2}.
\]
Let $c_n=(1-d/n)^{-1/2}\to1$ and write $\mu_n=\mu_{M_n/\sqrt d}$. For any bounded $1$-Lipschitz function $f$,
\begin{align*}
  \abs{\int f(c_nz)\,\dd\mu_n(z)-\int f(z)\,\dd\mu_n(z)}
  &\le\abs{c_n-1}\int\abs z\,\dd\mu_n(z)\\
  &\le\abs{c_n-1}
  \left(\frac1n\norm{M_n/\sqrt d}_{\HS}^2\right)^{1/2}\\
  &=\abs{c_n-1}.
\end{align*}
Thus the bounded-Lipschitz distance between the ESDs at the two normalizations tends to zero deterministically. Combining this with \eqref{eq:sqrt-d-circular} proves the theorem.
\end{proof}

\subsection{The superlogarithmic bridge}\label{subsec:superlog-bridge}

\begin{proposition}[Superlogarithmic bridge]\label{prop:superlog-bridge}
Fix $\eps>0$. Let $d=d(n)$ be an integer sequence satisfying
\[
  (\log n)^{1+\eps}\le d,
  \qquad
  d=o(n).
\]
Then
\[
  \mu_{M_n/\sqrt{d(1-d/n)}}\xrightarrow{\Pp}\mu_{\mathrm{circ}}.
\]
\end{proposition}

\begin{proof}
For all sufficiently large $n$, one has $d\le n/2$. Let $B_\eps$ be supplied by \cref{cor:uniform-polynomial}, fix $z\in\C\setminus\set{0}$, and define $A,\wt A$ by \eqref{eq:shifted-comparators}. Put
\[
  \tau_{n,\eps}=n^{-B_\eps},
  \qquad
  \mathcal L_{n,\eps}(z)
  =\set{s_n(A)\ge\tau_{n,\eps}}
  \cap\set{s_n(\wt A)\ge\tau_{n,\eps}}.
\]
By \cref{cor:uniform-polynomial},
\begin{equation}\label{eq:superlog-good-event}
  \Pp(\mathcal L_{n,\eps}(z)^c)=O(d^{-1/2})=o(1).
\end{equation}
Choose the balanced exponent
\begin{equation}\label{eq:balanced-a}
  a=a_\eps:=\frac{2+\eps}{2(1+\eps)}.
\end{equation}
Then
\begin{equation}\label{eq:a-identities}
  1-a=\frac{\eps}{2(1+\eps)},
  \qquad
  a(1+\eps)-1=\frac\eps2.
\end{equation}
Set
\begin{equation}\label{eq:superlog-prefix-parameters}
  r=\ceil{\frac n{d^a}},
  \qquad
  k=n-r,
  \qquad
  \eta=\frac{\tau_{n,\eps}}n.
\end{equation}
Let $C_M^{(\eps)}$ and $C_B^{(\eps)}$ be the first $k$ rows of $A$ and $\wt A$. Since $p=d/n=o(1)$,
\begin{equation}\label{eq:superlog-Q}
  Q=p(1-p)(r+1)\asymp d^{1-a}
  \ge c(\log n)^{\eps/2}\to\infty.
\end{equation}
Moreover, $U\le2(1+\abs z)\sqrt n$ and
\[
  \log\frac U\eta=O_\eps(\log n).
\]
Therefore \cref{prop:row-replacement} gives
\begin{align}
  \frac1n\E\abs{
    \Psi_\eta(C_M^{(\eps)})-\Psi_\eta(C_B^{(\eps)})
  }
  &\le C\sqrt{\frac{\log(ed)}{d^{1-a}}}
  +C_\eps\sqrt d\,\log n\,e^{-cd^{1-a}}
  \notag\\
  &=o(1).
  \label{eq:superlog-replacement}
\end{align}
Indeed, $\log d=o(d^{1-a})$. Also
\[
  \log(\sqrt d\log n)=\tfrac12\log d+\log\log n
  =O_\eps(\log d)=o(d^{1-a}),
\]
where $d\ge(\log n)^{1+\eps}$ is used in the middle estimate. Thus the exponential term dominates its prefactor uniformly in the stated range.

Markov's inequality and the same interlacing calculation used in the lower band yield, after extending both log-volume random variables by zero on $\mathcal L_{n,\eps}(z)^c$,
\begin{equation}\label{eq:superlog-prefix-volume}
  \frac1n\abs{
    \log\vol(C_M^{(\eps)})-\log\vol(C_B^{(\eps)})
  }=\littleoP.
\end{equation}
Here the regularization is removed on $\mathcal L_{n,\eps}(z)$ exactly as before because $\eta=\tau_{n,\eps}/n$.

It remains to pass from prefixes to square determinants. By \eqref{eq:positive-prefix-tail},
\begin{equation}\label{eq:superlog-positive-tail}
  \frac1n\sum_{j=1}^r\logp s_j(A)
  =O(d^{-a}\log d)=o(1),
\end{equation}
with the analogous $o_{\Pp}(1)$ estimate for $\wt A$. On $\mathcal L_{n,\eps}(z)$, every one of the bottom $r$ row-slice singular values may pay the polynomial smallest-singular-value cost:
\begin{align}
  \frac1n\sum_{j=n-r+1}^n\logm s_j(A)
  &\le\frac rn B_\eps\log n
  \notag\\
  &\le C_\eps d^{-a}\log n
  \le C_\eps(\log n)^{-\eps/2}=o(1),
  \label{eq:superlog-row-bottom-tail}
\end{align}
where \eqref{eq:a-identities} is used in the last step. Since $r=o(n)$, \cref{cor:Bernoulli-bottom-tail} gives the corresponding estimate for $B_n/\sqrt d-zI_n$; multiplication by $\sqrt d\ge1$ can only decrease $\logm$, so the same estimate holds for $\wt A$.

Applying \eqref{eq:prefix-square-error} to both matrices and using \eqref{eq:superlog-prefix-volume}, we obtain
\begin{equation}\label{eq:superlog-det-comparison}
  \frac1n\log\abs{\det\!\left(\frac{M_n}{\sqrt d}-zI_n\right)}
  -\frac1n\log\abs{\det\!\left(\frac{B_n}{\sqrt d}-zI_n\right)}
  \xrightarrow{\Pp}0
\end{equation}
for every fixed $z\ne0$.

The Hilbert--Schmidt estimates, the replacement principle, and the optimal sparse i.i.d. circular law now give
\[
  \mu_{M_n/\sqrt d}\xrightarrow{\Pp}\mu_{\mathrm{circ}}
\]
exactly as in the proof of \cref{thm:lower-band-circular}. Finally, $d=o(n)$ implies $(1-d/n)^{-1/2}\to1$, and the same Schur/Hilbert--Schmidt argument transfers the conclusion to the variance normalization.
\end{proof}

\begin{remark}[Relation to the Li--Litvak--Yu range]
For every fixed $\eta\in(0,1)$, \cite[Theorem~1.3]{LLY26a} already gives the circular law when
\[
  d\ge(\log n)^{2+\eta},
  \qquad
  d=o(n).
\]
Thus, to connect the lower band to the previously available result, \cref{prop:superlog-bridge} is needed only on the intervening range between $(\log n)^{1+\eps}$ and a threshold of the form $(\log n)^{2+\eta}$. We stated and proved the bridge for the entire superlogarithmic range because the same argument is uniform there and requires no additional input.
\end{remark}

\subsection{Completion of the main theorem}

\begin{proof}[Proof of \cref{thm:main}]
Fix $\sigma=1/2$ and split the indices into two sets:
\[
  d\le(\log n)^{3/2}
  \qquad\text{and}\qquad
  d>(\log n)^{3/2}.
\]
On the first set, apply \cref{thm:lower-band-circular}. On the second set, apply \cref{prop:superlog-bridge} with $\eps=1/2$. Each subsequence converges in probability to the same deterministic circular law. Therefore the original sequence does as well.
\end{proof}

\section{Further directions}\label{sec:further-directions}

\subsection{Why the lower-band argument is logarithmic}

The superlogarithmic bridge removes any need to run the structural machinery throughout the full range $d=o(n)$ and supplies the missing connection to the Li--Litvak--Yu circular-law regime. Within the lower polylogarithmic band, where the LLTTY discretization, the codimension choice $m\asymp n/d^3$, the Gaussian-normal source range, and the conditioning-to-slice estimate are used, the numerical structural and intermediate-singular-value budgets improve with $d$:
\[
  q_0=O(n/d^2),
  \qquad
  \frac{q_0}{n}\,d\log d=O\!\left(\frac{\log d}{d}\right)\to0,
  \qquad
  \delta_Lv_L\asymp\frac{d^{1/192}}{\log d}\gg\log d.
\]
The explicit lower logarithmic restriction enters in \cref{lem:max-weight}. In the present $\ell$-decomposition,
\[
  \diam(\supp B)\le\log_2n+O(\log d),
\]
whereas that lemma requires the diameter to be $O(d)$. Below logarithmic degree, the exact occupancy functional
\[
  \Phi_d(B)=\E 2^{-\max_{a\le d}B_a}
\]
can no longer be compressed uniformly into
$2^{-\E B}e^{-cD(B)}$. A sublogarithmic proof would need to retain $\Phi_d(B)$ or establish a more refined entropy--occupancy inequality that separates bins of mass above and below $1/d$.

The current proof also does not reach every fixed constant in $d=C\log n$. The support-diameter inequality and the imported restricted-norm and smallest-singular-value estimates require a sufficiently large absolute lower constant. Reducing that constant is a quantitative problem distinct from passing from the logarithmic band to larger degrees, which is handled by \cref{prop:superlog-bridge}.

A separate question concerns the sharp scale of the unshifted smallest singular value. Li, Litvak, and Yu proved
$s_n(M_n)=O_{\Pp}(\sqrt d/n)$ in the dense regime $d=pn$ for each fixed $p\in(0,1/2]$, and conjectured the same upper scale whenever
$\min(d,n-d)\ge(1+\eps)\log n$~\cite[Theorem~1.1, Remark~1.4, and Conjecture~1.5]{LLY26b}. The present argument instead uses lower bounds for shifted matrices and controls a block of intermediate singular values; it neither resolves nor requires that sharp-scale conjecture.

\subsection{Defective columns below the invertibility threshold}

A fixed column of $M_n$ is zero with probability
\[
  \left(1-\frac dn\right)^n
  =\exp\!\left(-d+O\!\left(\frac{d^2}{n}\right)\right)
  =e^{-(1+o(1))d},
\]
where the last equality uses $d=o(n)$. Thus below the logarithmic scale the matrix has zero columns with substantial probability, and $s_n(M_n)=0$ at $z=0$. This does not by itself rule out the circular law when $d\to\infty$, because the proportion of zero columns is $e^{-(1+o(1))d}\to0$ and hermitization is an almost-everywhere statement in the shift. It does show that any sublogarithmic proof must tolerate an exceptional coordinate set rather than demand global invertibility or expansion.

The structural theorem already allows an exceptional set of size
\[
  n d^{-1/192+o(1)},
\]
much larger than the expected number
\[
  n\exp\!\left(-d+O\!\left(\frac{d^2}{n}\right)\right)
\]
of zero columns when $d\to\infty$. This suggests a defect-tolerant version of the present argument: first remove low-degree columns, then prove many-level geometry and random-normal anti-clustering on the complement. The operator-norm scale also changes from $O(\sqrt d)$ to $O(\sqrt d+\log n)$, so the net radii would have to be rebalanced.

\subsection{Conjectural optimal range}

\begin{conjecture}
The independent row-slice model obeys the circular law whenever
\[
  d\to\infty,
  \qquad
  d=o(n).
\]
\end{conjecture}

This is consistent with the i.i.d. Bernoulli and doubly regular models. For fixed $d$, a positive limiting proportion of zero columns prevents convergence to the circular law. \Cref{thm:main} reaches every sparse degree above a sufficiently large logarithmic constant. The remaining asymptotic gap is therefore the sublogarithmic regime, together with the quantitative question of how small the constant in $d=C\log n$ can be. The present paper identifies two ingredients likely needed below logarithmic degree: an occupancy-sensitive replacement for the support-diameter max-weight lemma and a defect-tolerant structural theorem.

\appendix

\section{Source extraction for the LLTTY deterministic consequences}\label{app:LLTTY-extraction}

This appendix proves \cref{prop:LLTTY-package} as an author-derived result. Page and equation locators refer to \cite{LLTTY19}, arXiv:1801.05575v4. The table below separates exact source statements, immediate corollaries, and conclusions reconstructed from source proofs. No item is presented as a verbatim theorem unless it is one.

\small
\setlength{\LTleft}{0pt}
\setlength{\LTright}{0pt}
\setlength{\tabcolsep}{3pt}
\begin{longtable}{>{\raggedright\arraybackslash}p{0.13\linewidth}>{\raggedright\arraybackslash}p{0.24\linewidth}>{\raggedright\arraybackslash}p{0.20\linewidth}>{\raggedright\arraybackslash}p{0.35\linewidth}}
\toprule
\textbf{Item} & \textbf{Source locator} & \textbf{Status} & \textbf{Deduction recorded here}\\
\midrule
\endfirsthead
\toprule
\textbf{Item} & \textbf{Source locator} & \textbf{Status} & \textbf{Deduction recorded here}\\
\midrule
\endhead
\eqref{eq:class-cardinality} and class count
& pp.~34--35, Lemma~5.1, (14)--(16); p.~37, Lemmas~5.4--5.5
& direct restatement after reindexing
& Restriction to $\mathcal A_k(\mathcal X)$ cannot increase the number of nonempty equivalence classes.\\
$S_L=\bigcup K_u(L)\cup P_v(L)$
& p.~38, Theorem~5.6
& verbatim after notation change
& Codimension enters only through the previously defined gradual class.\\
\eqref{eq:factorial-compression}
& p.~55, Lemma~7.2
& stated ``in particular'' consequence
& Reindex the four source families by the manuscript's part index $q$.\\
\eqref{eq:small-bin-mass}
& p.~61, start of Case~2 in the proof of Lemma~7.7
& proof-level consequence
& Each of four geometrically decreasing families contributes $O(d^{-\kappa}n)$ below the cutoff.\\
\eqref{eq:K1} and \eqref{eq:K1-product}
& pp.~60--61, proof of Lemma~7.7, source (32)--(33), (38), Case~1
& author-derived packaging
& Multiply the three source families and use $n/s_q\le d^\kappa$.\\
\eqref{eq:K2} and \eqref{eq:K2-separation}
& p.~61, Case~2 in the proof of Lemma~7.7
& intermediate proof conclusion
& Define the first lower and last upper occupied weight orders and divide the total small-part mass by the per-order bound.\\
Low-order conclusion in part~(v)
& p.~62, Lemma~7.8
& direct source conclusion
& A repeated rounded value produces a weight set of mass at least $\delta/36$.\\
High-order conclusion and \eqref{eq:Pv-separation}
& pp.~62--63, Lemma~7.9 and its proof
& intermediate proof conclusion
& The height lower bound and $\wt w_q\ge h_q/d$ yield the separated upper order.\\
\bottomrule
\end{longtable}
\normalsize

Throughout the appendix, the source's gradual class $S$ is our codimension-dependent class $S_L$, and its families $K_u,P_u$ are denoted by $K_u(L),P_u(L)$. The codimension enters only through the definition of the steep classes and hence through $S_L$; the arguments in source Sections~5--7 are otherwise unchanged.

For a fixed canonical $k$-approximation, reindex all nonempty regular and spread $\ell$-parts as $L^{(q)}$, and write
\[
  s_q=\abs{L^{(q)}},
  \qquad
  h_q=h(L^{(q)}).
\]
The source's truncated weight $\wt w_q$ is exactly \eqref{eq:truncated-weights}; its $w$-set $W_b$ is the union of the parts with $2^b\le\wt w_q<2^{b+1}$. Thus the random dyadic order $B$ used here satisfies
\[
  \Pp\set{B=b}=\frac{\abs{W_b}}n.
\]
No rescaling or rebinning of the source weights is made.

\subsection{Equivalence classes and the gradual decomposition}

Source Lemma~5.1 states that, when $d\le n^{1/3}$ and
$1\le k\le\sqrt n/d^{3/2}$, the normalized canonical $k$-approximations form at most $e^n$ equivalence classes under equality of the coordinate multiset. This is precisely the relation $\sim_k$ on $\mathcal A_k(L)$. Restricting an equivalence relation to a subset cannot increase the number of nonempty classes, which gives the corresponding assertion for $\mathcal A_k(\mathcal X)$.

Source Lemma~5.4 shows that permutation of a canonical approximation preserves, order by order, the type, height, and cardinality of every nonempty $\ell$-part. Hence these data are well defined on each equivalence class. Source Lemma~5.5 then separates spread and regular parts by their orders and gives, in its notation,
\[
  \abs{\mathcal C}
  \le n!\prod_{j\ge0}
  \frac{(h_j^S)^{c_j^S}(h_j^R)^{c_j^R}}{c_j^S!c_j^R!}.
\]
After reindexing every nonempty spread or regular part by $q$, this is exactly
\[
  \abs{\mathcal C}
  \le\frac{n!}{\prod_qs_q!}\prod_qh_q^{s_q},
\]
which is \eqref{eq:class-cardinality}. Finally, source Theorem~5.6 states, for every integer $v\ge5$,
\[
  S_L=\bigcup_{u=4}^{v}K_u(L)\cup P_v(L).
\]
This proves parts~(i) and~(ii) of \cref{prop:LLTTY-package}.

\subsection{Factorial compression inside one weight order}

Fix a dyadic order $b$, put
\[
  I(b)=\set{q:2^b\le\wt w_q<2^{b+1}},
  \qquad
  t_b=\sum_{q\in I(b)}s_q=\abs{W_b}.
\]
Source Lemma~7.2 shows that the part sizes in a fixed $w$-set, arranged in nonincreasing order, are controlled by a geometric sequence. Its stated ``in particular'' conclusion is
\[
  \prod_{q\in I(b)}s_q!\ge t_b!e^{-Ct_b},
\]
which is \eqref{eq:factorial-compression}.

We also need the small-part consequence used in the second case of the proof of source Lemma~7.7. Let
\[
  T=d^{-\kappa}n,
  \qquad
  J(b)=\set{q\in I(b):s_q<T}.
\]
The proof of source Lemma~7.2 splits the parts in one $w$-set into four families according to regular/spread type and large/small status. Within each family, the sizes decrease geometrically with the order at which the parts were created. Consequently, the sum of the members of that family below the cutoff $T$ is $O(T)$. Summing the four families gives
\[
  \abs{\bigcup_{q\in J(b)}L^{(q)}}\le CT=Cd^{-\kappa}n.
\]
This is also the estimate written explicitly at the start of Case~2 in the proof of source Lemma~7.7. It proves part~(iii) of \cref{prop:LLTTY-package}.

\subsection{\texorpdfstring{The two $K_u$ alternatives}{The two Ku alternatives}}

Let $y$ be the $d^u$-approximation of a vector in $K_u(L)$. By the source definition of $K_u$, the union of its spread $\ell$-parts has size at least $\cK n_3$. Set
\[
  \kappa=\frac{\cK a_3}{16}<\frac13,
  \qquad
  \mathscr I=\set{q:s_q\ge d^{-\kappa}n}.
\]
Either
\[
  \sum_{q\in\mathscr I}s_q\ge n-\frac{\cK n_3}{4}
\]
or its complement has mass at least $\cK n_3/4$. These are alternatives \eqref{eq:K1} and \eqref{eq:K2}.

Assume \eqref{eq:K1}. Write $\mathscr I_1$ for the large spread parts, $\mathscr I_2$ for the large regular parts, and $\mathscr I_3=(\mathscr I_1\cup\mathscr I_2)^c$. Since the total spread mass is at least $\cK n_3$,
\[
  \sum_{q\in\mathscr I_1}s_q\ge\frac{3\cK n_3}{4},
  \qquad
  \sum_{q\in\mathscr I_3}s_q\le\frac{\cK n_3}{4}.
\]
For these three families, the definition of the truncated weights and the source lower bounds corresponding to its equations~(32)--(33) give
\begin{align*}
  \wt w_q^{-s_q}
  &\le d^{-s_q/2}\left(\frac n{s_q}\right)^{s_q/2}h_q^{-s_q},
  &&q\in\mathscr I_1,\\
  \wt w_q^{-s_q}
  &=\left(\frac n{s_q}\right)^{s_q}h_q^{-s_q},
  &&q\in\mathscr I_2,\\
  \wt w_q^{-s_q}
  &\le d^{s_q}h_q^{-s_q},
  &&q\in\mathscr I_3.
\end{align*}
Because $n/s_q\le d^\kappa$ on $\mathscr I_1\cup\mathscr I_2$, the same three-family calculation as in Case~1 of the proof of source Lemma~7.7 yields
\begin{align*}
  \prod_qh_q^{s_q}\prod_q\wt w_q^{-s_q}
  &\le C^n d^{-3\cK n_3/8}d^{\kappa n}d^{\cK n_3/4}\\
  &\le C^n d^{-\cK n_3/16}
  \le C^n d^{-c_1n}
\end{align*}
for an absolute $c_1>0$, after a harmless adjustment for $n_3=\floor{a_3n}$. This is \eqref{eq:K1-product}.

Assume now \eqref{eq:K2}. For each $b$, let
\[
  J(b)=\mathscr I^c\cap I(b).
\]
The preceding subsection gives
\[
  \abs{\bigcup_{q\in J(b)}L^{(q)}}\le C'd^{-\kappa}n.
\]
Following Case~2 in the proof of source Lemma~7.7, define $b_1$ to be the first order for which the cumulative lower $w$-sets have mass at least $\cK n_3/16$, and $b_2$ to be the last order for which the cumulative upper $w$-sets have the same mass. Then
\[
  \Pp\set{B\le b_1}\ge\frac{\cK n_3}{16n},
  \qquad
  \Pp\set{B\ge b_2}\ge\frac{\cK n_3}{16n}.
\]
The small parts between these two cutoffs have total mass at least $\cK n_3/8$. Since each intervening order contributes at most $C'd^{-\kappa}n$,
\[
  b_2-b_1+1
  \ge\frac{\cK n_3}{8C'd^{-\kappa}n}
  \ge cd^\kappa.
\]
After reducing fixed constants and taking $n$ large, this gives
\[
  \Pp\set{B\le b_1}\ge c_0,
  \qquad
  \Pp\set{B\ge b_2}\ge c_0,
  \qquad
  b_2-b_1\ge c_1d^\kappa.
\]
This proves part~(iv) of \cref{prop:LLTTY-package}.

\subsection{\texorpdfstring{The clustered $P_v$ alternative}{The clustered Pv alternative}}

Let $x\in P_v(L)$ have at least $\delta n$ coordinates in a disk of radius $\rho\le d^{-v}$, and let $y$ be its $d^v$-approximation. Source Lemma~7.8 shows that one value of $y$ is repeated on at least $\delta n/9$ coordinates and therefore produces a $w$-set of mass at least $\delta n/36$ and order at most
\[
  b_1=\log_2\left(\frac{72\sqrt d}{\delta}\right),
\]
with the integer rounding used in the source.

By the source definition of $P_v$, at least $\cP n_3$ coordinates belong to parts whose heights are at least
\[
  \cP a_3\,2^{\cP a_3(v-4)}.
\]
Since every truncated weight satisfies $\wt w_q\ge h_q/d$, these coordinates lie in $w$-sets above the high order used in the proof of source Lemma~7.9. The separation hypothesis \eqref{eq:Pv-separation-hyp} is exactly the hypothesis under which that proof establishes the intermediate conclusion
\[
  b_2-b_1\ge\frac{\cP a_3}{2}(v-4).
\]
Thus
\[
  \Pp\set{B\le b_1}\ge\frac\delta{36},
  \qquad
  \Pp\set{B\ge b_2}\ge\frac{\cP n_3}{n}\ge c_0
\]
for all sufficiently large $n$, after reducing the common absolute constant $c_0$ used above. The size condition \eqref{eq:Pv-delta} and the displayed separation condition are precisely the remaining assumptions of source Lemma~7.9. This proves part~(v) of \cref{prop:LLTTY-package} and completes the source extraction.

\section{Dependency ledger and source interfaces}\label{app:dependencies}

For clarity, we list the imported results and the new arguments proved in this paper. This appendix is not logically necessary, but it makes the proof interfaces auditable.

\subsection{Imported statements}

\paragraph{Li--Litvak--Yu \cite{LLY26a}.}
All theorem and remark numbers cited from this preprint refer to arXiv:2604.10446v1 (April 2026). The source mappings used here are: Theorem~1.3 for the circular law above $\log^{2+\eps}n$; Theorem~1.7 and Remark~1.9 for the row-slice smallest-singular-value estimate and its refined threshold; Theorem~3.1 and Remark~3.2 for the Bernoulli counterpart; and Theorem~4.1 for the rectangular restricted operator norm. The deduction of a single polynomial threshold above $(\log n)^{1+\eps}$ is not quoted as a source theorem: it is carried out in \cref{cor:uniform-polynomial} from the formula in source Remark~1.9. Section~3 and Corollary~3.5 of the source are methodological antecedents for the replacement and conditioning steps, respectively; the arbitrary-codimension regularized comparison used here is proved in \cref{prop:row-replacement}. We do not import a codimension-$L$ right-kernel theorem from this source.

\paragraph{LLTTY19 \cite{LLTTY19}.}
Source checks use arXiv:1801.05575v4. Section~4 supplies the codimension-dependent steep classes $T_i(L)$, their constant-shifted versions, and the exact gradual class $S_L$; the particular locators used later are Section~4.1, equation~(8), for the shifted very-steep inclusion, Lemma~4.5 for the shifted nets, and Lemma~4.11 for the deterministic almost-constant reduction. Sections~5--7 supply the canonical $k$-approximation, the $\ell$-decomposition, source nets, and equivalence classes. The author-derived consequences in \cref{prop:LLTTY-package} are extracted from Lemmas~5.1, 5.4, 5.5, 7.2, 7.7--7.9 and Theorem~5.6. A lemma-by-lemma derivation of the factorial, geometric, and clustered-$P_v$ conclusions, including the intermediate conclusions taken from the proofs of source Lemmas~7.7--7.9, is given in \cref{app:LLTTY-extraction}. The row-local small-ball input is Lemma~6.4, and the tensorization form is Lemma~2.7. The clustered-$P_v$ hypotheses are reproduced in \eqref{eq:Pv-delta}--\eqref{eq:Pv-separation-hyp}.

\paragraph{LLTTY21 \cite{LLTTY21}.}
Source checks use arXiv:1801.05576v2. \Cref{prop:Gaussian-anti-clustering} isolates a deterministic fixed-subspace consequence of Theorem~4.2, Propositions~4.4 and~4.7, and Lemmas~4.8--4.10. Its proof gives a separate verification of the dimension requirements, the source linear-combination count $p_{\src}$, the cutoff $\ell_0\le4n/L$, and the radius, cluster-size, and steepness comparisons. The deterministic tie-break is the one explicitly suggested immediately before source Lemma~4.8. No random-regular-digraph probability estimate is imported.

\paragraph{Rudelson--Tikhomirov \cite{RT19}.}
The optimal sparse i.i.d. circular law under the minimal divergence condition is used only as a comparator result. No intermediate estimate from that paper is imported into the row-slice structural argument.

\paragraph{SSS25 \cite{SSS25}.}
Source checks use arXiv:2310.17600v2. \Cref{prop:Bernoulli-hermitization} reproduces the complete logarithmic-potential squeeze needed here. Fact~10.2 supplies bounded-continuous convergence to the deterministic shifted Gaussian singular-value law $\nu_z^G$, and Fact~10.1 identifies $U^\circ(z)=-\int\log t\,\dd\nu_z^G(t)$ and proves finiteness. Lemma~8.2, Fact~9.1, and Lemma~9.3 give the explicit upper and lower bounds for the full, untruncated logarithmic potential. The resulting statement holds for every fixed $z\ne0$, rather than only on an unspecified full-measure set. The Hilbert--Schmidt bound in \cref{prop:Bernoulli-hermitization} is elementary.

\paragraph{Tao's concentration theorem \cite{Tao12}.}
The convex concentration statement used for regularized Bernoulli distances is Theorem~2.1.13. It gives Gaussian tails for a convex $1$-Lipschitz function of bounded independent complex coordinates, centered at a median. The median-to-mean calculation is recorded in \cref{lem:convex-concentration}.

\paragraph{Tao--Vu \cite{TV10}.}
The replacement principle used in the lower-band and superlogarithmic arguments is Theorem~2.1. Its Hilbert--Schmidt and logarithmic-determinant hypotheses are verified explicitly in \cref{thm:lower-band-circular,prop:superlog-bridge}.

\subsection{New statements proved here}

The following components are specific to the independent row-slice lower-band argument and its superlogarithmic completion.
\begin{enumerate}
\item The max-weight inequality, \cref{lem:max-weight}.
\item The averaged multinomial row-profile estimate and its conditioning back to a row slice, \cref{prop:averaged-profile} and \eqref{eq:condition-to-slice}.
\item The entropy--energy cancellation $H(B)\le C+\log(1+D(B))$ and the completion of the $K_u$ and clustered-$P_v$ class sums uniformly for $C_*\log n\le d\le(\log n)^{1+\sigma}$.
\item The rectangular LLTTY compatibility theorem, \cref{thm:kernel-compatibility}: direct independent-row exclusion of $T_0^{\mathcal K}(L)$, $T_1^{\mathcal K}(L)$, $T_2^{\mathcal K}(L)$, and the nonzero vectors in $\mathcal B(10/d^3)\setminus T_3^{\mathcal K}(L)$, with literal identification of the remaining nonzero kernel vectors with $\cone{S_L}$.
\item The rectangular row-slice structural theorem, \cref{thm:row-slice-structural}, uniformly for codimensions $L=m$ and $L=m+1$ throughout the lower band.
\item The exact-slice double-counting lemma, \cref{lem:slice-double-count}.
\item The lower-band intermediate-singular-value bound, \cref{thm:intermediate}.
\item The generalized regularized row-replacement estimate, \cref{prop:row-replacement}, valid for arbitrary codimension with $p(1-p)(r+1)\ge1$.
\item The abstract negative-log tail lemma, \cref{lem:negative-log-tail}, and the lower-band prefix-to-square synthesis proving \cref{thm:lower-band-circular}.
\item The superlogarithmic bridge, \cref{prop:superlog-bridge}, which combines the uniform polynomial smallest-singular-value scale with the codimension $r=\ceil{n/d^{a_\eps}}$ and connects the lower band to the Li--Litvak--Yu circular-law range. The proof remains valid throughout the full superlogarithmic range.
\end{enumerate}

\subsection{Measurability and probability bookkeeping}

All row-set statements are uniform for each deterministic fixed set of the indicated size. Measurable enlargements are produced by the fixed standard-basis extension rule in \cref{sec:normals}, and greedy Gaussian clusters use the fixed total-order tie-break. For a fresh slice row, independence is stated relative to the sigma-algebra generated by the exposed rows and all auxiliary Gaussian randomness.

The structural theorem holds for each fixed exposed row set with failure $O(1/n)$, not simultaneously over all such sets. In the leave-one-row-out argument, no union bound over $n$ indices is taken. Instead, the per-index structural failure contributes only $O(1)$ to the expected number of bad indices. Markov's inequality is then applied at the much larger threshold $n/d^2$.

All convergence statements in both hermitization regimes are in probability. Random errors arise only through expectation bounds followed by Markov's inequality, deterministic estimates on events of probability $1-o(1)$, or imported Bernoulli convergence in probability. The proof is carried out separately for each fixed nonzero shift $z$; there is no union bound over uncountably many shifts. Excluding the single point $z=0$ is immaterial for the almost-everywhere hypothesis in the replacement principle.

For the lower band, the prefix codimension is $r=\ceil{n/\sqrt d}$, and only $q_0=O(n/d^2)$ bottom singular values pay the full smallest-singular-value cost. In the superlogarithmic bridge, the codimension is $r=\ceil{n/d^{a_\eps}}$, and all $r$ bottom singular values may pay the uniform polynomial cost.

\subsection{Normalization}

All determinant comparisons are performed at the $\sqrt d$ normalization. The variance normalization is recovered only after the circular law has been proved:
\[
  \frac{M_n}{\sqrt{d(1-d/n)}}
  =(1-d/n)^{-1/2}\frac{M_n}{\sqrt d},
\]
and the deterministic scalar tends to one because $d=o(n)$. This avoids invoking any unproved uniformity of the Bernoulli logarithmic-potential theorem under varying shifts.

\begin{thebibliography}{LLTTY21}

\bibitem[Bai97]{Bai97}
Zhidong D. Bai,
\emph{Circular law},
Ann. Probab. \textbf{25} (1997), no.~1, 494--529.

\bibitem[BC12]{BC12}
Charles Bordenave and Djalil Chafa\"i,
\emph{Around the circular law},
Probab. Surv. \textbf{9} (2012), 1--89.

\bibitem[Coo19]{Coo19}
Nicholas Cook,
\emph{The circular law for random regular digraphs},
Ann. Inst. Henri Poincar\'e Probab. Stat. \textbf{55} (2019), no.~4, 2111--2167.

\bibitem[FKS22]{FKS22}
Asaf Ferber, Matthew Kwan, and Lisa Sauermann,
\emph{Singularity of sparse random matrices: simple proofs},
Combin. Probab. Comput. \textbf{31} (2022), no.~1, 21--28.

\bibitem[Gir85]{Gir85}
V.~L. Girko,
\emph{Circular law},
Theory Probab. Appl. \textbf{29} (1985), no.~4, 694--706.

\bibitem[LLTTY19]{LLTTY19}
Alexander E. Litvak, Anna Lytova, Konstantin Tikhomirov,
Nicole Tomczak-Jaegermann, and Pierre Youssef,
\emph{Structure of eigenvectors of random regular digraphs},
Trans. Amer. Math. Soc. \textbf{371} (2019), no.~11, 8097--8172;
source checks use arXiv:1801.05575v4.

\bibitem[LLTTY21]{LLTTY21}
Alexander E. Litvak, Anna Lytova, Konstantin Tikhomirov,
Nicole Tomczak-Jaegermann, and Pierre Youssef,
\emph{Circular law for sparse random regular digraphs},
J. Eur. Math. Soc. \textbf{23} (2021), no.~2, 467--501;
source checks use arXiv:1801.05576v2.

\bibitem[LLY26a]{LLY26a}
Dongbin Li, Alexander E. Litvak, and Tingzhou Yu,
\emph{The circular law for sparse random combinatorial matrices},
arXiv:2604.10446v1 (April 2026).

\bibitem[LLY26b]{LLY26b}
Dongbin Li, Alexander E. Litvak, and Tingzhou Yu,
\emph{An upper bound on the smallest singular value of dense random combinatorial matrices},
J. Complexity \textbf{95} (2026), 102031.

\bibitem[NV13]{NV13}
Hoi H. Nguyen and Van H. Vu,
\emph{Circular law for random discrete matrices of given row sum},
J. Comb. \textbf{4} (2013), no.~1, 1--30.

\bibitem[RT19]{RT19}
Mark Rudelson and Konstantin Tikhomirov,
\emph{The sparse circular law under minimal assumptions},
Geom. Funct. Anal. \textbf{29} (2019), no.~2, 561--637.

\bibitem[SSS25]{SSS25}
Ashwin Sah, Julian Sahasrabudhe, and Mehtaab Sawhney,
\emph{The sparse circular law, revisited},
Bull. Lond. Math. Soc. \textbf{57} (2025), no.~2, 330--358;
source checks use arXiv:2310.17600v2.

\bibitem[Tao12]{Tao12}
Terence Tao,
\emph{Topics in random matrix theory},
Graduate Studies in Mathematics, vol.~132,
American Mathematical Society, Providence, RI, 2012.

\bibitem[Tra20]{Tra20}
Tuan Tran,
\emph{The smallest singular value of random combinatorial matrices},
arXiv:2007.06318 (2020).

\bibitem[TV10]{TV10}
Terence Tao and Van Vu,
\emph{Random matrices: universality of empirical spectral distributions and the circular law},
Ann. Probab. \textbf{38} (2010), no.~5, 2023--2065,
with an appendix by Manjunath Krishnapur.

\end{thebibliography}


\end{document}
